\documentclass{article}
\usepackage[letterpaper,margin=1.0in]{geometry}
%%%%% NEW MATH DEFINITIONS %%%%%

\usepackage{amsmath}\allowdisplaybreaks
\usepackage{amsfonts,bm}
\usepackage{amssymb}
\usepackage{amsthm}
\usepackage{amsmath}
\usepackage{algorithm}
\usepackage{algorithmic}
\usepackage{multirow}
\usepackage{booktabs}
\usepackage{xcolor}

% Mark sections of captions for referring to divisions of figures
\newcommand{\figleft}{{\em (Left)}}
\newcommand{\figcenter}{{\em (Center)}}
\newcommand{\figright}{{\em (Right)}}
\newcommand{\figtop}{{\em (Top)}}
\newcommand{\figbottom}{{\em (Bottom)}}
\newcommand{\captiona}{{\em (a)}}
\newcommand{\captionb}{{\em (b)}}
\newcommand{\captionc}{{\em (c)}}
\newcommand{\captiond}{{\em (d)}}

% Highlight a newly defined term
\newcommand{\newterm}[1]{{\bf #1}}


% Figure reference, lower-case.
\def\figref#1{figure~\ref{#1}}
% Figure reference, capital. For start of sentence
\def\Figref#1{Figure~\ref{#1}}
\def\twofigref#1#2{figures \ref{#1} and \ref{#2}}
\def\quadfigref#1#2#3#4{figures \ref{#1}, \ref{#2}, \ref{#3} and \ref{#4}}
% Section reference, lower-case.
\def\secref#1{section~\ref{#1}}
% Section reference, capital.
\def\Secref#1{Section~\ref{#1}}
% Reference to two sections.
\def\twosecrefs#1#2{sections \ref{#1} and \ref{#2}}
% Reference to three sections.
\def\secrefs#1#2#3{sections \ref{#1}, \ref{#2} and \ref{#3}}
% Reference to an equation, lower-case.
\def\eqref#1{equation~(\ref{#1})}
% Reference to an equation, upper case
\def\Eqref#1{Equation~\ref{#1}}
% A raw reference to an equation---avoid using if possible
\def\plaineqref#1{\ref{#1}}
% Reference to a chapter, lower-case.
\def\chapref#1{chapter~\ref{#1}}
% Reference to an equation, upper case.
\def\Chapref#1{Chapter~\ref{#1}}
% Reference to a range of chapters
\def\rangechapref#1#2{chapters\ref{#1}--\ref{#2}}
% Reference to an algorithm, lower-case.
\def\algref#1{algorithm~\ref{#1}}
% Reference to an algorithm, upper case.
\def\Algref#1{Algorithm~\ref{#1}}
\def\twoalgref#1#2{algorithms \ref{#1} and \ref{#2}}
\def\Twoalgref#1#2{Algorithms \ref{#1} and \ref{#2}}
% Reference to a part, lower case
\def\partref#1{part~\ref{#1}}
% Reference to a part, upper case
\def\Partref#1{Part~\ref{#1}}
\def\twopartref#1#2{parts \ref{#1} and \ref{#2}}

\def\ceil#1{\left\lceil #1 \right\rceil}
\def\floor#1{\left\lfloor #1 \right\rfloor}
\def\1{\bf{1}}
\newcommand{\train}{\mathcal{D}}
\newcommand{\valid}{\mathcal{D_{\mathrm{valid}}}}
\newcommand{\test}{\mathcal{D_{\mathrm{test}}}}
\newcommand{\Norm}[1]{\left\| #1 \right\|}
\newcommand{\norm}[1]{\left\| #1 \right\|_2}
\def\inner#1#2{\langle #1, #2 \rangle}

\def\eps{{\varepsilon}}

% Vectors
\def\vzero{{\bf{0}}}
\def\vone{{\bf{1}}}
\def\vmu{{\bm{\mu}}}
\def\vphi{{\bm{\phi}}}
\def\veta{{\bm{\eta}}}
\def\valpha{{\bm{\alpha}}}
\def\vbeta{{\bm{\beta}}}
\def\vtheta{{\bm{\theta}}}
\def\va{{\bf{a}}}
\def\vb{{\bf{b}}}
\def\vc{{\bf{c}}}
\def\vd{{\bf{d}}}
\def\ve{{\bf{e}}}
\def\vf{{\bf{f}}}
\def\vg{{\bf{g}}}
\def\vh{{\bf{h}}}
\def\vi{{\bf{i}}}
\def\vj{{\bf{j}}}
\def\vk{{\bf{k}}}
\def\vl{{\bf{l}}}
\def\vm{{\bf{m}}}
\def\vn{{\bf{n}}}
\def\vo{{\bf{o}}}
\def\vp{{\bf{p}}}
\def\vq{{\bf{q}}}
\def\vr{{\bf{r}}}
\def\vs{{\bf{s}}}
\def\vt{{\bf{t}}}
\def\vu{{\bf{u}}}
\def\vv{{\bf{v}}}
\def\vw{{\bf{w}}}
\def\vx{{\bf{x}}}
\def\vy{{\bf{y}}}
\def\vz{{\bf{z}}}


% Graph
\def\fA{{\mathcal{A}}}
\def\fB{{\mathcal{B}}}
\def\fC{{\mathcal{C}}}
\def\fD{{\mathcal{D}}}
\def\fE{{\mathcal{E}}}
\def\fF{{\mathcal{F}}}
\def\fG{{\mathcal{G}}}
\def\fH{{\mathcal{H}}}
\def\fI{{\mathcal{I}}}
\def\fJ{{\mathcal{J}}}
\def\fK{{\mathcal{K}}}
\def\fL{{\mathcal{L}}}
\def\fM{{\mathcal{M}}}
\def\fN{{\mathcal{N}}}
\def\fO{{\mathcal{O}}}
\def\fP{{\mathcal{P}}}
\def\fQ{{\mathcal{Q}}}
\def\fR{{\mathcal{R}}}
\def\fS{{\mathcal{S}}}
\def\fT{{\mathcal{T}}}
\def\fU{{\mathcal{U}}}
\def\fV{{\mathcal{V}}}
\def\fW{{\mathcal{W}}}
\def\fX{{\mathcal{X}}}
\def\fY{{\mathcal{Y}}}
\def\fZ{{\mathcal{Z}}}

% Sets
\def\sA{{\mathbb{A}}}
\def\sB{{\mathbb{B}}}
\def\sC{{\mathbb{C}}}
\def\sD{{\mathbb{D}}}
\def\BE{{\mathbb{E}}}
% Don't use a set called E, because this would be the same as our symbol
% for expectation.
\def\BF{{\mathbb{F}}}
\def\BG{{\mathbb{G}}}
\def\BH{{\mathbb{H}}}
\def\BI{{\mathbb{I}}}
\def\BJ{{\mathbb{J}}}
\def\BK{{\mathbb{K}}}
\def\BL{{\mathbb{L}}}
\def\BM{{\mathbb{M}}}
\def\BN{{\mathbb{N}}}
\def\BO{{\mathbb{O}}}
\def\BP{{\mathbb{P}}}
\def\BQ{{\mathbb{Q}}}
\def\BR{{\mathbb{R}}}
\def\BS{{\mathbb{S}}}
\def\BT{{\mathbb{T}}}
\def\BU{{\mathbb{U}}}
\def\BV{{\mathbb{V}}}
\def\BW{{\mathbb{W}}}
\def\BX{{\mathbb{X}}}
\def\BY{{\mathbb{Y}}}
\def\BZ{{\mathbb{Z}}}

\def\mA {{\bf A}}
\def\mB {{\bf B}}
\def\mC {{\bf C}}
\def\mD {{\bf D}}
\def\mE {{\bf E}}
\def\mF {{\bf F}}
\def\mG {{\bf G}}
\def\mH {{\bf H}}
\def\mI {{\bf I}}
\def\mJ {{\bf J}}
\def\mK {{\bf K}}
\def\mL {{\bf L}}
\def\mM {{\bf M}}
\def\mN {{\bf N}}
\def\mO {{\bf O}}
\def\mP {{\bf P}}
\def\mQ {{\bf Q}}
\def\mR {{\bf R}}
\def\mS {{\bf S}}
\def\mT {{\bf T}}
\def\mU {{\bf U}}
\def\mV {{\bf V}}
\def\mW {{\bf W}}
\def\mX {{\bf X}}
\def\mY {{\bf Y}}
\def\mZ {{\bf Z}}

\def\AUC{{\rm{AUC}}}
\def\SVI{{\rm{SVI}}}
\def\VI{{\rm{VI}}}
\def\ApproxMax{{\rm{ApproxMax}}}
\def\Gap{{\rm{Gap}}}

% The true underlying data generating distribution
\newcommand{\pdata}{p_{\rm{data}}}
% The empirical distribution defined by the training set
\newcommand{\ptrain}{\hat{p}_{\rm{data}}}
\newcommand{\Ptrain}{\hat{P}_{\rm{data}}}
% The model distribution
\newcommand{\pmodel}{p_{\rm{model}}}
\newcommand{\Pmodel}{P_{\rm{model}}}
\newcommand{\ptildemodel}{\tilde{p}_{\rm{model}}}
% Stochastic autoencoder distributions
\newcommand{\pencode}{p_{\rm{encoder}}}
\newcommand{\pdecode}{p_{\rm{decoder}}}
\newcommand{\precons}{p_{\rm{reconstruct}}}

\newcommand{\laplace}{\mathrm{Laplace}} % Laplace distribution

\newcommand{\E}{\mathbb{E}}
\newcommand{\Ls}{\mathcal{L}}
\newcommand{\R}{\mathbb{R}}
\newcommand{\emp}{\tilde{p}}
\newcommand{\lr}{\alpha}
\newcommand{\reg}{\lambda}
\newcommand{\rect}{\mathrm{rectifier}}
\newcommand{\softmax}{\mathrm{softmax}}
\newcommand{\sigmoid}{\sigma}
\newcommand{\softplus}{\zeta}
\newcommand{\KL}{D_{\mathrm{KL}}}
\newcommand{\Var}{\mathrm{Var}}
\newcommand{\standarderror}{\mathrm{SE}}
\newcommand{\Cov}{\mathrm{Cov}}
% Wolfram Mathworld says $L^2$ is for function spaces and $\ell^2$ is for vectors
% But then they seem to use $L^2$ for vectors throughout the site, and so does
% wikipedia.
\newcommand{\normlzero}{L^0}
\newcommand{\normlone}{L^1}
\newcommand{\normltwo}{L^2}
\newcommand{\normlp}{L^p}
\newcommand{\normmax}{L^\infty}

\newcommand{\parents}{Pa} % See usage in notation.tex. Chosen to match Daphne's book.

\DeclareMathOperator*{\argmax}{arg\,max}
\DeclareMathOperator*{\argmin}{arg\,min}

\DeclareMathOperator{\sign}{sign}
\DeclareMathOperator{\Tr}{Tr}
\let\ab\allowbreak

\DeclareMathOperator{\spn}{span}
\DeclareMathOperator{\prox}{prox}

%\usepackage{amsthm}
%\theoremstyle{plain}
\newtheorem{thm}{Theorem}[section]
\newtheorem{dfn}{Definition}[section]
\newtheorem{exmp}{Example}[section]
\newtheorem{lem}{Lemma}[section]
\newtheorem{asm}{Assumption}[section]
\newtheorem{rmk}{Remark}[section]
\newtheorem{cor}{Corollary}[section]
\newtheorem{prop}{Proposition}[section]
\newtheorem{cond}{Condition}[section]
\makeatletter
\def\Ddots{\mathinner{\mkern1mu\raise\p@
\vbox{\kern7\p@\hbox{.}}\mkern2mu
\raise4\p@\hbox{.}\mkern2mu\raise7\p@\hbox{.}\mkern1mu}}
\makeatother

\def\pr#1{\sP\left( #1 \right)}

\makeatletter
\newcommand*{\rom}[1]{\expandafter\@slowromancap\romannumeral #1@}
\makeatother



\usepackage{hyperref}
\usepackage{enumitem}
%\usepackage{natbib}
\usepackage[utf8]{inputenc} % allow utf-8 input
\usepackage[T1]{fontenc}    % use 8-bit T1 fonts
\usepackage{hyperref}       % hyperlinks
\usepackage{url}            % simple URL typesetting
\usepackage{booktabs}       % professional-quality tables
\usepackage{amsfonts}       % blackboard math symbols
\usepackage{nicefrac}       % compact symbols for 1/2, etc.
\usepackage{microtype}      % microtypography
%\usepackage{lipsum}
%\usepackage{fancyhdr}       % header
\usepackage{graphicx}       % graphics
\usepackage[numbers,sort]{natbib}
\usepackage{multirow}
\usepackage{makecell}
\usepackage{bbding}
\usepackage{wrapfig}
\graphicspath{{media/}}     % organize your images and

%\pagestyle{fancy}
%\thispagestyle{empty}
%\rhead{ \textit{ }} 

\usepackage{thmtools}
\usepackage{thm-restate}
\usepackage{threeparttable}
% Definitions of handy macros can go here
\usepackage{makecell}

\title{Near-Optimal Nonconvex-Strongly-Convex
Bilevel Optimization with Fully First-Order Oracles}

\author{
  Lesi Chen\textsuperscript{* \dag}\\
 { \small{\texttt{chenlc23@mails.tsinghua.edu.cn}}}
  \and
  Yaohua Ma\textsuperscript{* \dag}\\
 { \small{\texttt{ma-yh21@mails.tsinghua.edu.cn}}}
  \and 
  Jingzhao Zhang\textsuperscript{ \dag}\\
 { \small{\texttt{jingzhaoz@mail.tsinghua.edu.cn}}}
}

%  \author{Lesi Chen\thanks{School of Data Science, Fudan University; lschen19@fudan.edu.cn}
% }

\date{}
      
\begin{document}
\maketitle
\begingroup
\begin{NoHyper}
\renewcommand\thefootnote{*}
\footnotetext{Equal contributions.}
% \renewcommand\thefootnote{\S}
% \footnotetext{
% School of Data Science, Fudan University, China.}
\renewcommand\thefootnote{\dag}
\footnotetext{Institute for Interdisciplinary Information Sciences, Tsinghua University, China.}
\end{NoHyper}
\endgroup
\begin{abstract}
Bilevel optimization has wide applications such as hyperparameter tuning, neural architecture search, and meta-learning.
Designing efficient algorithms for bilevel optimization is challenging because the lower-level problem defines a feasibility set implicitly via another optimization problem. 
In this work, we consider one
tractable case when the lower-level problem is strongly convex. Recent works show that with a Hessian-vector product oracle, one can provably find
an $\epsilon$-first-order stationary point
within $\tilde{\mathcal{O}}(\epsilon^{-2})$ oracle calls. However, Hessian-vector product may be inaccessible or expensive in practice.
Kwon et al. (ICML 2023) addressed this issue by proposing a  first-order method that can achieve the same goal at a slower
rate of $\tilde{\mathcal{O}}(\epsilon^{-3})$. 
In this work, we provide a tighter analysis demonstrating that this method can converge at the near-optimal $\tilde {\mathcal{O}}(\epsilon^{-2})$ rate as second-order methods.
Our analysis further leads to
simple first-order algorithms that achieve similar convergence rates for finding second-order stationary points and for distributed bilevel problems.
\end{abstract}

\section{Introduction}

In this paper, we consider the following bilevel optimization problem:
\begin{align} \label{hyper-refo}
    \min_{x \in \BR^{d_x}} \varphi(x) := f(x,y^*(x)) , \quad {\rm where}~ y^*(x) = \arg \min_{y \in \BR^{d_y}} g(x,y).
\end{align}
We call $f$ the upper-level problem, $g$ the lower-level problem, and $\varphi$ the hyper-objective.
This formulation covers many applications, including but not limited to hyperparameter tuning~\citep{franceschi2018bilevel,pedregosa2016hyperparameter,gao2022value,bae2020delta,mackay2018self}, neural architecture search~\citep{liu2018darts,wang2022zarts,zoph2016neural,zhang2021idarts}, meta-learning~\citep{finn2017model,fallah2020convergence,andrychowicz2016learning,mishchenko2023convergence,rajeswaran2019meta,chandra2022gradient,ji2022theoretical,zhou2019efficient}, adversarial training~\citep{zhang2022revisiting,bruckner2011stackelberg,wang2021fast,bishop2020optimal,wang2022solving,robey2023adversarial} and sample reweighting~\citep{gao2022self,zhou2022model,ren2018learning,yong2022holistic,shu2019meta,li2022ood}.

This work focuses on the oracle complexity~\citep{nesterov2018lectures} of solving Problem \ref{hyper-refo}.
Since the hyper-objective $\varphi(x)$ is usually a nonconvex function, 
a common goal for non-asymptotic analysis is to find an approximate
stationary point~\citep{carmon2020lower,carmon2021lower,arjevani2023lower,fang2018spider,fang2019sharp}. In general, $\varphi(x)$ can be non-differentiable or even discontinuous, so one cannot define a stationary point~\citep{chen2023bilevel}.
A common condition to ensure the differentiability of $\varphi(x)$  adopted by many existing works~\citep{dagreou2022framework,ji2021bilevel,yang2021provably,ji2022will,ji2021lower,chen2021closing,khanduri2021near,hong2023two} 
is to assume the lower-level strong convexity.
When the lower-level problem is strongly convex in $y$,
the so-called hyper-gradient $\nabla \varphi(x)$~\citep{pedregosa2016hyperparameter,grazzi2020iteration} can be expressed by:
\begin{align} \label{hyper-grad} 
\begin{split}
    \nabla \varphi(x) 
    &= 
    \nabla_x f(x,y^*(x)) + \nabla y^*(x) \nabla_y f(x,y^*(x)) \\
    &=
    \nabla_x f(x,y^*(x)) - \nabla_{xy}^2 g(x,y^*(x)) [ \nabla_{yy}^2 g(x,y^*(x))]^{-1} \nabla_y f(x,y^*(x)).
\end{split}
\end{align}
This equation enables one to implement
second-order methods \citep{ji2021bilevel,ghadimi2018approximation} for nonconvex-strongly-convex bilevel optimization. 
These methods query 
second-order information of $g$ to 
estimate $\nabla \varphi(x)$ via \Eqref{hyper-grad}, and then
perform the so-called hyper-gradient descent on $\varphi(x)$. 
By using known convergence results of gradient descent for smooth nonconvex objectives, these methods can find an $\epsilon$-first-order stationary point of $\varphi(x)$ with $\tilde \fO(\epsilon^{-2})$ first-order oracle calls from $f$ and $\tilde \fO(\epsilon^{-2})$ second-order oracle calls from $g$.  
In practice, these methods are usually implemented with the Jacobian/Hessian-vector product oracle\footnote{For function $g(x,y)$ and vector $v$, the Jacobian-vector product oracle returns $ \nabla_{xy}^2 g(x,y) v $, and the Hessian-vector product oracle returns $\nabla_{yy}^2 g(x,y)v$.}.
However, in many applications~\citep{sow2022convergence,song2019maml,finn2017model,nichol2018first}, calling the Jacobian/Hessian-vector product oracle may be very costly and become the computational bottleneck in bilevel optimization.

Very recently,  \citet{kwon2023fully}  proposed a novel first-order algorithm that can find an $\epsilon$-first-order stationary point of $\varphi(x)$ without resorting to second-order information.  
The key idea is to exploit the following value-function reformulation~\citep{outrata1990numerical,dempe2020bilevel,lin2014solving,gao2022value,ye2023difference} for \Eqref{hyper-refo}:
\begin{align} \label{value-refo}
    \min_{x \in \BR^{d_x}, y \in \BR^{d_y}} f(x,y) , \quad {\rm subject ~to}~~ g(x,y) - g^*(x) \le 0.
\end{align}
Here $g^*(x) = g(x,y^*(x))$ is usually called the value-function~\citep{ye1995optimality,ye2010new,dempe2002foundations}. The 
Lagrangian with a multiplier $\lambda >0$ for this inequality-constrained problem takes the form of:
\begin{align} \label{eq:L}
    \fL_{\lambda}(x,y):= f(x,y)+ \lambda ( g(x,y) - g^*(x)).
\end{align}
\citet{kwon2023fully} further define the following auxiliary function:
\begin{align}\label{eq:L-lambda}
\begin{split}
     \fL_{\lambda}^*(x) := \fL_{\lambda}(x,y_{\lambda}^*(x)),~
     {\rm where }~ y_{\lambda}^*(x) = \arg \min_{y \in \BR^{d_y}} \fL_{\lambda}(x,y),
\end{split}
\end{align}
and showed that $\fL_{\lambda}^*(x)$ is a good proxy of $\varphi(x)$ with
\begin{align*}
    \Vert \nabla \fL_{\lambda}^*(x) - \nabla \varphi(x)  \Vert = \fO \left( \kappa^3 /\lambda \right).
\end{align*}
This indicates that when we set $\lambda \asymp \epsilon^{-1}$, an $\epsilon$-first-order stationary point of $\fL_{\lambda}^*(x)$ is also an $\fO(\epsilon)$-first-order stationary point of $\varphi(x)$. Therefore, we
can reduce the problem of finding an $\epsilon$-first-order stationary point of $\varphi(x)$ to finding that of $\fL_{\lambda}^*(x)$.
The advantage of this reduction is that 
 $\nabla \fL_{\lambda}^*(x)$ can be evaluated with only first-order information of $f$ and $g$:
 \begin{align} \label{eq:nabla-Lag}
 \begin{split}
     \nabla \fL_{\lambda}^*(x) &= 
     \nabla_x \fL_{\lambda}^*(x,y_{\lambda}^*(x)) + \nabla y_{\lambda}^*(x) \underbrace{\nabla_y \fL_{\lambda}^*(x,y_{\lambda}^*(x))}_{=0} 
     \\
     &=\nabla_x f(x, y_{\lambda}^*(x)) + \lambda ( \nabla_x g(x,y_{\lambda}^*(x)) - \nabla_x g(x,y^*(x))).
 \end{split}
 \end{align}
 Based on this property, \citet{kwon2023fully} proposed the Fully First-order Bilevel Approximation\footnote{
\citet{kwon2023fully} called  both their 
stochastic and deterministic method the same name Fully First-order Stochastic Approximation
(F${}^2$SA). Here we 
call the deterministic method Fully First-order Bilevel Approximation (F${}^2$BA) for ease of distinction.} (F${}^2$BA), whose procedure is presented in Algorithm \ref{alg:F2BA}. 
Since $\fL_{\lambda}(x,y)$ in \Eqref{eq:L} is $\fO(\lambda)$-gradient Lipschitz, direct analysis shows that it requires a complexity of $\tilde \fO(\lambda \epsilon^{-2}) = \tilde \fO(\epsilon^{-3} )$ to find an $\epsilon$-first-order stationary point, which is slower than the $\tilde \fO(\epsilon^{-2})$ complexity of second-order methods. Then \citet{kwon2023fully} conjectured that a fundamental gap may exist between first-order and second-order methods.  

However, this paper refutes this conjecture by showing that first-order methods can also achieve the near-optimal $\tilde \fO(\epsilon^{-2})$ rate as second-order methods, even though less information is used.
We prove that for a sufficiently large $\lambda$, the proxy $\fL_{\lambda}^*(x)$ not only has the same approximate first-order stationary point of $\varphi(x)$ , but also has almost the same gradient Lipschitz coefficient of $\varphi(x)$ given by:
\begin{align*}
    \Vert \nabla^2 \fL_{\lambda}^*(x) \Vert \asymp \Vert \nabla^2 \varphi(x) \Vert \asymp \kappa^3,
\end{align*}
where $\kappa$ is the condition number that will be formally defined later. This improved analysis indicates that although $\lambda$ is large, the gradient Lipschitz coefficient of $\fL_{\lambda}^*(x)$ would remain constant and not depend on $\epsilon$. As a consequence, we can improve the complexity of F${}^2$BA to $\tilde \fO(\epsilon^{-2})$. Meanwhile, we prove this result
under more relaxed assumptions than previous analyses for first-order methods, and we
summarize the comparison in Table \ref{tab:res-1st}.
As nonconvex-strongly-concave bilevel optimization problems subsume standard nonconvex optimization problems, the $\tilde \fO(\epsilon^{-2})$ oracle call is also optimal up to logarithmic factors according to the lower bound provided by ~\cite{carmon2020lower}.

Our analysis technique is general and it can simplify algorithms in multiple setups.
If we additionally assume the Hessian Lipschitz continuity of $f$ and the third-order derivative Lipschitz continuity of $g$, we can show that
\begin{align*}
    \Vert \nabla^2 \fL_{\lambda}^*(x) - \nabla^2 \varphi(x) \Vert = \fO(\kappa^6 /\lambda) ~~{\rm and}~~  \Vert \nabla^3 \fL_{\lambda}^*(x) \Vert \asymp \Vert \nabla^3 \varphi(x) \Vert \asymp \kappa^5.
\end{align*} 
Based on this observation, we propose the perturbed F${}^2$BA (Algorithm \ref{alg:PF2BA}), which can provably find an $\epsilon$-second-order stationary point of $\varphi(x)$ within $\tilde \fO(\epsilon^{-2})$ first-order oracle calls.
%matching the complexity of F${}^2$BA for finding an $\epsilon$-first-order stationary point, up to logarithmic factors. 
Our result shows that first-order methods can also escape saddle points in bilevel optimization like second-order methods~\citep{huang2022efficiently}.

Our algorithms can also be easily applied to the distributed scenario, when both the upper and lower-level functions adopt the following finite-sum structure:
\begin{align}  \label{prob:fn}
    f(x,y) := \frac{1}{m} \sum_{i=1}^m f_i(x,y), \quad g(x,y) := \frac{1}{m} \sum_{i=1}^m g_i(x,y),
\end{align}
where $f_i$ and $g_i$ denote the local function on the $i$-th agent. 
One of the primal challenges for second-order methods in distributed bilevel optimization arises from the Hessian inverse in the expression of hyper-gradient (\Eqref{hyper-grad}) since directly aggregating the Hessian matrix requires an expensive $\fO(d_{y}^2)$ communication complexity~\citep{yangdecentralized,lu2022stochastic}. To address this issue, existing works reformulate the Hessian inverse as a sub-problem and solve it with distributed algorithms that only require communicating Hessian-vector products~\citep{tarzanagh2022fednest,chen2022decentralized,xiao2023communication}.
In contrast, the extension of F${}^2$BA to the distributed setting (Algorithm~\ref{alg:CF2BA}) naturally avoids this additional distributed sub-solver for matrix inverse. The proposed method achieves near-optimal computation and communication complexity simultaneously.

\begin{table*}[t]
    \centering
    \caption{
    We present the oracle complexities of different methods for finding an $\epsilon$-first-order stationary point of the hyper-objective $\varphi(x):= f(x,y^*(x))$ under Assumption \ref{asm:sc}.}
    \label{tab:res-1st}
    \begin{tabular}{c c c c}
    \hline 
    Oracle     & Method  & 
    Additional Assumption  & Oracle Calls    \\
    \hline \hline 
    \addlinespace
         & AID \citep{ghadimi2018approximation} & No & $\fO( \epsilon^{-2.5} )$ 
    \\
     \addlinespace
    1st+2nd & AID \citep{ji2021bilevel} & No & $\fO ( \epsilon^{-2}   )$  \\
    \addlinespace
    & ITD \citep{ji2021bilevel} & No & $ \fO ( \epsilon^{-2} \log(\nicefrac{1}{\epsilon}))$  \\
    \addlinespace
    \hline
    \addlinespace
    & PZOBO~\citep{sow2022convergence}  & Asmp.(a) & $\fO( d_{x}^2\epsilon^{-4} \log (\nicefrac{1}{\epsilon}))$
    \\
    \addlinespace
    & BOME~\citep{liu2022bome} &  Asmp.(b)  &$ \fO(  \epsilon^{-6} \log(\nicefrac{1}{\epsilon}))$   \\[-0.12cm]
    1st \\[-0.12cm]
    & F${}^2$BA~\citep{kwon2023fully} &  Asmp.(c) &$  \fO(  \epsilon^{-3} \log(\nicefrac{1}{\epsilon}))$  \\
    \addlinespace
     & 
F${}^2$BA (Our Analysis)
      & No & $  \fO( \epsilon^{-2} \log(\nicefrac{1}{\epsilon}))$  \\ 
     \addlinespace
     \hline
    \end{tabular}
    \begin{tablenotes}
		{\scriptsize   
  \item \textbf{Asmp.(a):} Assume $ \Vert \nabla^2 g(x,y) - \nabla^2 g(x',y') \Vert_F^2 \le \rho_g^2 (\Vert x-x' \Vert^2 +  \Vert y - y' \Vert^2)$, stronger than Assumption \ref{asm:sc}c. 
  \item \textbf{Asmp.(b):} Additionally assume both $\vert f(x,y) \vert $ and $ \vert g(x,y) \vert$ are upper bounded for any $x \in \BR^{d_x}$ and $y \in \BR^{d_y}$.
  \item \textbf{Asmp.(c):} Additionally assume $\nabla^2 f$ is Lipschitz and both $\Vert \nabla_x f(x,y)\Vert, \Vert \nabla_x g(x,y) \Vert $ are upper bounded.
   }
    \end{tablenotes}
\end{table*}

% \begin{table*}[t]
%     \centering
%     \caption{
%     We present the oracle complexities of different methods for finding an $\epsilon$-first-order stationary point of the hyper-objective $\varphi(x):= f(x,y^*(x))$ under Assumption \ref{asm:sc}.}
%     \label{tab:res-1st}
%     \begin{tabular}{c c c c c }
%     \hline 
%     Oracle     & Method  & 
%     \makecell[c]{Additional \\ Assumption} & \makecell[c]{Runtime Per \\ Iteration in $x$}
%       & \makecell[c]{Oracle \\ Calls}  \\
%     \hline \hline 
%     \addlinespace
%          & \makecell[c]{AID \\ \citep{ghadimi2018approximation}} & No & $\fO((d_xd_y +d_y^2)\epsilon^{-0.5})$ & $\fO( \epsilon^{-2.5} )$ 
%     \\
%      \addlinespace
%     1st+2nd &  \makecell[c]{AID \\ \citep{ji2021bilevel}} & No & $\fO(d_xd_y +d_y^2)$ & $\fO ( \epsilon^{-2}   )$  \\
%     \addlinespace
%     & \makecell[c]{ ITD \\ \citep{ji2021bilevel} } & No & $\tilde \fO(d_x d_y + d_y^2)$ &$ \tilde \fO ( \epsilon^{-2} )$  \\
%     \addlinespace
%     \hline
%     \addlinespace
%     & \makecell[c]{PZOBO \\ \citep{sow2022convergence}}  & Asmp.(a) & $ \tilde \fO(d_x d_y )$ & $\tilde \fO( d_{x}^2\epsilon^{-4})$
%     \\
%     \addlinespace
%     & \makecell[c]{BOME \\ \citep{liu2022bome} }  &  Asmp.(b)  & $\tilde \fO(d_x +d_y) $ & $ \tilde \fO(  \epsilon^{-6} )$   \\[-0.12cm]
%     1st \\[-0.12cm]
%     & \makecell[c]{F${}^2$BA \\ \citep{kwon2023fully}}  &  Asmp.(c) & $\tilde \fO(d_x +d_y) $ &$  \tilde \fO(  \epsilon^{-3} )$  \\
%     \addlinespace
%      & 
% \makecell[c]{F${}^2$BA \\ (Our Analysis, Algorithm \ref{alg:F2BA})}
%       & No & $\tilde \fO(d_x +d_y) $ & $  \tilde \fO( \epsilon^{-2} )$  \\ 
%      \addlinespace
%      \hline
%     \end{tabular}
%     \begin{tablenotes}
% 		{\scriptsize   
%   \item \textbf{Asmp.(a)} Assume $ \Vert \nabla^2 g(x,y) - \nabla^2 g(x',y') \Vert_F^2 \le \rho_g^2 (\Vert x-x' \Vert^2 +  \Vert y - y' \Vert^2)$, stronger than Assumption \ref{asm:sc}c. 
%   \item \textbf{Asmp.(b)} Additionally assume both $\vert f(x,y) \vert $ and $ \vert g(x,y) \vert$ are upper bounded for any $x \in \BR^{d_x}$ and $y \in \BR^{d_y}$.
%   \item \textbf{Asmp.(c)} Additionally assume $f(x,y)$ is $C_f$-Lipschitz in $x$ and $\rho_f$-Hessian Lipschitz in $(x,y)$.
%   }
%     \end{tablenotes}
% \end{table*}

% \begin{table}[t]
%     \centering
%     \begin{tabular}{c c c }
%     \hline 
%     Oracle     & Method   & Oracle Calls    \\
%     \hline \hline 
%     \addlinespace
%     \multirow{1}{*}{1st+2nd}     & Perturbed AID \citep{huang2022efficiently} & $\tilde \fO(\kappa^4 \epsilon^{-2} )$
%     \\
%     \addlinespace
%     \multirow{1}{*}{1st} 
%     & Perturbed F${}^2$BA (Algorithm \ref{alg:PF2BA}) 
%       & $  \tilde \fO( \kappa^4 \epsilon^{-2}) $  \\ 
%      \addlinespace
%      \hline
%     \end{tabular}
%     \caption{
%     We present the oracle complexities of different methods for finding $(\epsilon,\kappa^{2.5} \sqrt{\epsilon})$-second-order stationary points under Assumption \ref{asm:sc} and \ref{asm:third}.
%     Note that notation $\tilde \fO(\,\cdot\,)$ hides  polylogarithmic factors of $d_x, \kappa,\epsilon$.
%     } 
%     \label{tab:res-2nd}
% \end{table}

\paragraph{Notations.}
We use notations $\fO(\,\cdot\,),~ \tilde \fO(\,\cdot\,),~ \Omega(\,\cdot\,),~ \asymp$ as follows: given two functions $p: \BR^+ \rightarrow \BR^+$ and $q: \BR^+ \rightarrow \BR^+$,  $p(x) = \fO(q(x))$ means $ \lim \sup_{x \rightarrow +\infty} {p(x)}/{q(x)} < +\infty $; $p(x) = \tilde \fO(p(x)) $ means there exists some positive integer $k \ge 0$ such that $p(x) = \fO(q(x) \log^k (q(x)))$, $p(x) = \Omega(q(x))$ means $ \lim \sup_{x \rightarrow + \infty} p(x) / q(x) >0$, and $p(x) \asymp q(x)$ means we both have $p(x) = \fO(q(x))$ and $p(x) = \Omega(q(x))$.
% Given two sequences of positive scalars $p_n$ and $q_n$, we denote $p_n = \fO(q_n)$ if there exists an absolute constant $C>0$ such that $p_n \le Cq_n$. 
% We denote $p_n= \tilde \fO(q_n)$ if there exists an absolute constant $\tilde C>0$ which hides logarithmic factors such that $p_n \le \tilde Cq_n$.
% We denote $p_n = \Omega(q_n)$ if there exists an absolute constant $c>0$ such that $p_n \ge cq_n$, and denote $p_n \asymp q_n$ if we both have $p_n = \fO(q_n)$ and $p_n = \Omega(q_n)$.
%We use $\otimes$ to denote the Kronecker product of two matrices.
We use $I_d \in \BR^{d\times d}$ to denote a $d$=dimensional identity matrix. For two symmetric matrices $A$ and $B$, we use $A \succeq B$ to indicate that $A-B$ is positive semidefinite.
We use $\sB(r)$ to denote the Euclidean ball centered at the origin and radius $r$.
For a function $h: \BR^d \rightarrow \BR$, we use $\nabla h \in \BR^d, \nabla^2 h \in \BR^{d \times d},  \nabla^3 h \in \BR^{d \times d \times d}$ to denote its gradient, Hessian, and third-order derivative, and use $h^*$ to denote the global minimum of $h(\,\cdot\,)$. We
denote $\Vert \cdot\Vert$ to be the operator norm of a tensor and more details about the notations of tensors can be found in Appendix \ref{apx:tensor}.

\section{Related Work}

% Non-asymptotic analysis of second-order methods for nonconvex-strongly-convex bilevel optimization problems has been established by many works, but finite-time convergence of first-order methods has not been studied until recently.

\paragraph{Second-Order Methods for Bilevel Optimization.
} 

Most existing second-order methods for nonconvex-strongly-convex bilevel optimization can be categorized into
the approximate implicit differentiation (AID) methods  and the iterative differentiation
(ITD) methods. The AID approach~\citep{liao2018reviving,ji2021bilevel,ghadimi2018approximation,lorraine2020optimizing} constructs hyper-gradients explicitly according to \Eqref{hyper-grad} that $\nabla \varphi(x) = \nabla_x f(x,y^*(x)) - \nabla_{xy}^2 g(x,y^*(x)) v^*$, where $v^* $ is the solution to the linear system $\nabla_{yy}^2 g(x,y^*(x)) v^* =\nabla_y f(x,y^*(x)) $. Then one can use
iterative algorithms such as fix point iteration or conjugate gradient method to solve this linear system, avoiding the computation of the Hessian inverse \citep{grazzi2020iteration,ji2021bilevel,hong2023two}. 
The ITD approach~\citep{maclaurin2015gradient,shaban2019truncated,domke2012generic,franceschi2017forward,franceschi2018bilevel}  takes advantage of the fact that backpropagation can be efficiently implemented via modern
automatic differentiation frameworks such as PyTorch~\citep{paszke2019pytorch}. These methods approximate hyper-gradients by $\partial f(x,y^K(x)) / \partial x $, where $y^K(x)$ is the output from $K$-steps of gradient descent on $g(x,\,\cdot\,)$. Although the ITD approach does not query second-order information explicitly, the analytical form of $\partial f(x,y^K(x)) / \partial x $ involves second-order derivatives \citep{ji2021bilevel}, so they also belong to second-order methods. Both AID and ITD methods require $\tilde \fO(\epsilon^{-2})$ first-order oracle calls from $f$ and $\tilde \fO(\epsilon^{-2})$ second-order oracle calls from $g$ to find an $\epsilon$-first-order stationary point of $\varphi(x)$. Recently, \citet{huang2022efficiently} proposed the perturbed AID to find an $\epsilon$-second-order stationary point with $\tilde \fO(\epsilon^{-2})$ complexity under additional smoothness conditions of both $f$ and $g$. 

\paragraph{First-Order Methods for Bilevel Optimization.}

The basic idea of first-order methods for bilevel optimization is to approximate the hyper-gradient in \Eqref{hyper-grad} using gradient information. \citet{sow2022convergence} proposed the 
Partial Zeroth-Order based Bilevel Optimizer (PZOBO) that
applies a zeroth-order-like estimator to approximate the response Jacobian matrix $\nabla y^*(x) $ in \Eqref{hyper-grad}, and hence 
the complexity has a polynomial 
dependency on the dimension of the problem like the standard results in zeroth-order optimization~\citep{duchi2015optimal,ghadimi2013stochastic,kornowski2023algorithm}.
\citet{liu2022bome} first observed \Eqref{eq:nabla-Lag} that $\nabla \fL_{\lambda}^*(x)$ only involves first-order information and proposed the method named Bilevel Optimization Made Easy (BOME), 
but \citet{liu2022bome} did not provide any convergence result of $\varphi(x)$. Remarkably, \citet{kwon2023fully}
established
the relationship between $\fL_{\lambda}^*(x)$  and $\varphi(x)$, and proposed the Fully First-order Bilevel Approximation (F${}^2$BA) that can provably find an $\epsilon$-first-order stationary point of $\varphi(x)$ within $\tilde \fO(\epsilon^{-3})$ oracle complexity. As a by-product of their analysis, one can also show that BOME~\citep{liu2022bome} also converges within $\tilde \fO(\epsilon^{-6})$ first-order oracle complexity. 
However, before our work, we did not know whether first-order methods could have comparable theoretical guarantees to second-order methods.

\paragraph{Distributed Bilevel Optimization.}

Distributed optimization uses 
several agents to
train a large model in parallel \citep{boyd2011convex}. 
Despite the popularity of distributed first-order methods~\citep{liconvergence,recht2011hogwild,lin2019don}, the design of distributed second-order methods~\citep{wang2018giant,shamir2014communication,liu2023communication} is more challenging since it is inefficient to communicate the Jacobian/Hessian matrix across the agents.
This also becomes the main obstacle for second-order methods in distributed bilevel optimization~\citep{yangdecentralized,huang2023achieving,tarzanagh2022fednest,lu2022decentralized}. 
Note that directly aggregating the local hyper-gradients\footnote{The local hyper-gradient on the $i$-th agent with its local variable $(x_i,y_i)$ refers to the quantity $\nabla_x f_i(x_i,y_i)  - \nabla_{xy}^2 g_i(x_i,y_i) [\nabla_{yy}^2 g_i(x_i,y_i)]^{-1} \nabla_y f_i(x_i,y_i)$.} in \Eqref{hyper-grad} on all agents cannot give the correct form of $\nabla \varphi(x)$. 
Therefore,  
all the existing distributed second-order methods for bilevel optimization are more complicated than their single-machine counterparts. \citet{tarzanagh2022fednest} and \citet{chen2022decentralized} both
discussed how one can use an additional distributed 
sub-solver to overcome this challenge 
for AID methods.
\citet{xiao2023communication} proposed a well-designed aggregated iterative differentiation scheme (AggITD) to conduct the hyper-gradient
estimation and the lower-level optimization simultaneously. 

\section{Preliminaries}

In this section, we introduce the three different setups studied in this work. We focus on stating the assumptions and definitions used later, while delaying a more comprehensive description to future sections where we state and describe our main results.

\paragraph{First-Order Stationary Points} 

We first discuss the assumptions for finding first-order stationary points of the hyper-objective $\varphi(x)$, detailed below.

\begin{asm} \label{asm:sc}
Suppose that
\begin{enumerate}[label=\alph*.]
    \item $g(x,y)$ is $\mu$-strongly convex in $y$;
     \item $g(x,y)$ is $L_g$-gradient Lipschitz;
    \item $g(x,y)$ is $\rho_g$-Hessian Lipschitz;
     \item $f(x,y)$ is $C_f$-Lipschitz in $y$;
    \item $f(x,y)$ is $L_f$-gradient Lipschitz;
    \item $f(x,y)$ is two-times continuous differentiable;
    \item $\varphi(x)$ is lower bounded, \emph{i.e.} $\inf_{x \in \BR^{d_x}} \varphi(x) > - \infty$;
\end{enumerate}
\end{asm}

The above assumptions are common and necessary for nonasymptotic analyses. According to Example 5.1 in \citep{chen2023bilevel},
the hyper-objective $\varphi(x)$ can be discontinuous in general. For this reason, existing non-asymptotic analyses for bilevel optimization commonly make the 
lower-level strong convexity assumption (Assumption \ref{asm:sc}a)
 to ensure the existence of the hyper-gradient $\nabla \varphi(x)$. Under this lower-level strong convexity assumption,  $\nabla \varphi(x)$ can be expressed jointly by $ \nabla_x f(x,y)$, $\nabla_y f(x,y)$, $\nabla_{xy}^2 g(x,y)$ and $ \nabla_{yy}^2 g(x,y)$ as \Eqref{hyper-grad}.
This expression indicates that we need the smoothness conditions for $f$ and $g$ (Assumption \ref{asm:sc}b - \ref{asm:sc}e) to guarantee the gradient  Lipschitz continuity of $ \varphi(x)$. 
Besides these, 
we adopt Assumption \ref{asm:sc}f to ensure that $ \fL_{\lambda}(x,y)$ (\Eqref{eq:L-lambda}) is two-times continuous differentiable, and
Assumption \ref{asm:sc}g to ensure the bilevel optimization problem (\Eqref{hyper-refo}) is well-defined. 
%All these assumptions are adopted in previous works
% ~\citep{dagreou2022framework,ji2021bilevel,yang2021provably,ji2022will,ji2021lower,chen2022single,hong2023two,huo2023new,khanduri2021near}.

\begin{dfn}  \label{dfn:kappa-sc}
Under Assumption \ref{asm:sc}, we define the largest smoothness constant $\ell: =\max \{ C_f,L_f,L_g,\rho_g  \}$ and the condition number 
$\kappa:= \ell / \mu$. 
\end{dfn}

We can derive 
from the above assumptions that $\nabla \varphi(x)$ is uniquely defined and Lipschitz continuous, formally stated as follows.

\begin{restatable}[Lemma 2.2 by \citet{ghadimi2018approximation}]{prop}{propFsmooth}
    \label{prop:F-smooth}
Under Assumption \ref{asm:sc},
the hyper-gradient $\nabla \varphi(x)$ is uniquely defined by \Eqref{hyper-grad}, and
the hyper-objective $\varphi(x)$ is $L_\varphi$-gradient Lipschitz, where $L_\varphi = \fO(\ell \kappa^3)$.
% \begin{align*}
%     L_F := \left( 1 + \frac{L_g}{\mu} \right) \left( L_f + \frac{C_f \rho_g}{\mu} + \frac{L_f L_g}{\mu} + \frac{C_f L_g \rho_g}{\mu^2} \right) = \fO(\kappa^3).
% \end{align*}
 \end{restatable}

As the above proposition ensures that the hyper-objective $\varphi(x)$ is differentiable, we can define the $\epsilon$-first-order stationary points as follows.
\begin{dfn} \label{dfn:sta-hyper}
Given a differentiable function $\varphi(x): \BR^d \rightarrow \BR $,
we call $\hat x$ an $\epsilon$-first-order stationary point of $\varphi(x)$ if $ \Vert \nabla \varphi(\hat x) \Vert \le \epsilon$.
\end{dfn}

\paragraph{Second-Order Stationary Points}

Algorithms that pursue first-order stationary points may get stuck in saddle points and have poor performances \citep{,dauphin2014identifying}.
For
single-level optimization problems, there have been many researchers studying how to escape saddle points  
~\citep{jin2017escape,ge2015escaping,fang2019sharp,tripuraneni2018stochastic,agarwal2017finding,lee2016gradient,allen2018neon2,carmon2017convex,zhou2020stochastic,allen2018natasha,xu2018first,zhang2021escape}.
A common assumption in these works is to suppose that the objective is Hessian Lipschitz. 
When generalizing to bilevel optimization, we also expect $\varphi(x)$ to be Hessian Lipschitz, which can be proved if 
we further assume the following higher-order smoothness condition of $f$ and $g$.

%In many scenarios, it is not enough to merely find a first-order stationary point, since a first-order stationary point may be a saddle point that has poor xxxxxxxxxx~\citep{goodfellow2014qualitatively,choromanska2015loss}. This motivates us to pursue higher-order stationary points.When pursuing second-order stationarity of bilevel optimization, we additionally require the following assumptions \citep{huang2022efficiently}. 

\begin{asm} \label{asm:third}
Suppose that
\begin{enumerate}[label=\alph*.]
    \item $f(x,y)$ is three-times continuous differentiable;
    \item $f(x,y)$ is $\rho_f$-Hessian Lipschitz;
    \item  $ g(x,y)$ is $\nu_g$-third-order derivative Lipschitz.
\end{enumerate}
\end{asm}

\begin{dfn} \label{dfn:kappa-third}
Under Assumption \ref{asm:sc} and \ref{asm:third}, we define the largest smoothness constant $\ell: =\max \{ C_f,L_f,L_g,\rho_g,\rho_f,\nu_g  \}$ and the condition number 
$\kappa:= \ell / \mu$. 
\end{dfn}

\begin{prop}[Lemma 3.4 by \citet{huang2022efficiently}] \label{prop:Hess-smooth}
    Under Assumption \ref{asm:sc} and \ref{asm:third}, the hyper-objective $\varphi(x)$ is two-times continuously differentiable and $\rho_\varphi$-Hessian Lipschitz, where $\rho_\varphi = \fO(\ell \kappa^5)$.
\end{prop}

We can then formally define the
 approximate second-order stationary point as follows.
\begin{dfn}[\citet{nesterov2006cubic}] \label{dfn:sta-hyper-second}
Given a two-times continuously differentiable function $\varphi(x): \BR^d \rightarrow \BR $ with $\rho$-Lipschitz Hessian,
we call $\hat x$ an $\epsilon$-second-order stationary point of $\varphi(x)$ if 
\begin{align*}
    \Vert \nabla \varphi(\hat x) \Vert \le \epsilon,\quad  \nabla^2 \varphi(\hat x)  \succeq -\sqrt{\rho \epsilon}  I_{d}.
\end{align*}
\end{dfn}
We will discuss finding second-order stationary points for bilevel problems later in Section~\ref{sec:second}.

\paragraph{Distributed Bilevel Problems} 

As we have discussed, first-order methods have superiority over second-order methods when the Hessian-vector product oracle is expensive.
One of these scenarios is distributed bilevel optimization, where both $f$ and $g$ are distributed on $m$ agents as \Eqref{prob:fn}. The direct extension of second-order methods is to aggregate the local hyper-gradients in \Eqref{hyper-grad} on all the agents. Unfortunately, this simple algorithm may not converge and additional techniques are required to address this challenge~\citep{xiao2023communication,tarzanagh2022fednest,chen2022decentralized}. In contrast, the first-order algorithm F${}^2$BA does not have this issue since 
aggregating the local gradients in \Eqref{eq:nabla-Lag} gives the exact 
form of $\nabla \fL_{\lambda}^*(x)$.

Below, we introduce the notations for the distributed F${}^2$BA. 
Under the distributed setting, each agent $i$ has its own local variables within the optimization algorithm: $X(i) \in \BR^{d_x}, Y(i) \in \BR^{d_y},Z(i) \in \BR^{d_y}$~\footnote{$z$ is an auxiliary variable in F${}^2$BA that will be introduced in the subsequent section.}.
For the convenience of presenting the algorithm, we aggregate all the local variables (denoted by row vectors) in one matrix and denote
\begin{align*}
X=\begin{bmatrix}
X(1) \\ \vdots \\ X(m)
\end{bmatrix}&\in\BR^{m\times d_x},
\quad
Y=\begin{bmatrix}
Y(1) \\ \vdots \\ Y(m)
\end{bmatrix}\in\BR^{m\times d_y}
\quad \text{and} \quad
Z=\begin{bmatrix}
Z(1) \\ \vdots \\ Z(m)
\end{bmatrix}\in\BR^{m\times d_y},
\end{align*}
We denote~$\vone = [1,\cdots,1]^\top \in \BR^m$ and use the lowercase with the bar to represent the mean vector, such as $\bar x = \frac{1}{m} \vone^\top X$.
Let $d = d_x+d_y$ and we similarly define the aggregated gradients 
\begin{align*}
\nabla \vf(X,Y)=\begin{bmatrix}
\nabla f_1(X(1),Y(1)) \\ \vdots \\ \nabla f_m(X(m),Y(m))
\end{bmatrix}
\in\BR^{m\times d}, ~\nabla \vg(X,Y)=\begin{bmatrix}
\nabla g_1(X(1),Y(1)) \\ \vdots \\ \nabla g_m(X(m),Y(m))
\end{bmatrix}
\in\BR^{m\times d}.    
\end{align*}
Then we can present the update of the distributed algorithm in matrix form. More details will be provided later in Section~\ref{sec:distributed}.

\section{Finding First-Order Stationarity} \label{sec:first}

\begin{algorithm*}[t]  
\caption{F${}^2$BA $\left(x_0,y_0, \eta,\tau,\alpha, \lambda, T,K\right)$} \label{alg:F2BA}
\begin{algorithmic}[1] 
\STATE $ z_0 = y_0$ \\[1mm]
\STATE \textbf{for} $ t =0,1,\cdots,T-1 $ \\[1mm]
\STATE \quad $ y_t^0 =  y_{t}, ~ z_t^0 = z_{t}$ \\[1mm]
\STATE \quad \textbf{for} $ k =0,1,\cdots,K-1$ \\[1mm]
\STATE \quad \quad $ z_t^{k+1} = z_t^{k}- \alpha \nabla_y g(x_t, z_t^k)   $ \\[1mm]
\STATE \quad \quad $ y_t^{k+1} = y_t^k - \tau \left( \nabla_y f(x_t,y_t^k) +  \lambda  \nabla_y g(x_t,y_t^k) \right)$ \\[1mm]
\STATE \quad \textbf{end for} \\[1mm]
% \STATE \quad $z_{t+1} = z_t^K,~ y_{t+1} = y_t^K $ \\[1mm]
\STATE \quad $ \hat \nabla \fL_{\lambda}^*(x_t)= \nabla_x f(x_t,y_{t}^K) + \lambda ( \nabla_x g(x_t,y_{t}^K) - \nabla_x g(x_t,z_{t}^K) )$ \\[1mm]
\STATE \quad  $x_{t+1} = x_t -  \eta  \hat \nabla \fL_{\lambda}^*(x_t)$ \\[1mm]
% \STATE \quad \textbf{Option II: } $x_{t+1}  = x_t - \eta \hat \nabla \fL_{\lambda}^*(x_t) / \Vert \hat \nabla \fL_{\lambda}^*(x_t) \Vert $ \quad // Normalized Gradient Descent 
\STATE \textbf{end for} \\[1mm]
\end{algorithmic}
\end{algorithm*}

In bilevel optimization, the hyper-objective $\varphi(x)$ is usually a nonconvex function. Since finding the global minimum of a nonconvex function in the worst case requires an exponential number of
queries, a common compromise is to find a local minimum~\citep{ge2016matrix}. First-order stationary points (Definition \ref{dfn:sta-hyper}) are the points that satisfy the first-order necessary condition of a local minimum, which turns out to be a valid optimality criterion for nonconvex optimization.

For the task of finding an $\epsilon$-first-order stationary point of the hyper-objective $\varphi(x)$, 
\citet{kwon2023fully} proposed the algorithm F${}^2$BA (see Algorithm \ref{alg:F2BA}) and proved a upper complexity bound of $\tilde \fO(\epsilon^{-3})$. They noted that this upper bound is worse than the $\tilde \fO(\epsilon^{-2})$ near-optimal rate achieved by second-order methods, and hoped future works to investigate the fundamental limits of first-order methods. However, we give a surprising result in this section: the F${}^2$BA is already near-optimal if carefully analyzed.

The detailed procedure of F${}^2$BA is presented in Algorithm \ref{alg:F2BA}.  
The algorithm introduces an auxiliary variable $z \in \BR^{d_y}$, and performs gradient descent jointly in $x,y,z$ to solve the following optimization problem:
\begin{align}  \label{eq:opt-xyz}
     \min_{x \in \BR^{d_x}, y \in \BR^{d_y}} \left\{f(x,y) + \lambda \left(g(x,y) - \min_{z \in \BR^{d_y}} g(x,z)\right)\right\} \overset{\rm Eq. \ref{eq:L-lambda}}{=} \min_{x \in \BR^{d_x}} \fL_{\lambda}^*(x).
\end{align}
The intuition behind the algorithm is that optimizing $\fL_{\lambda}^*(x)$ is almost equivalent to optimizing $\varphi(x)$ when $\lambda$ is large.
Therefore, to analyze the convergence of the algorithm, we first characterize the relationship between 
$\fL_{\lambda}^*(x)$ and $\varphi(x)$ in the following lemmas.
 
% For a fixed $x$, the tasks of finding the best response variables $y$ and $z$ in \Eqref{eq:xyz} are two independent sub-problems that can be solved by gradient descent (Line 5-6 of Algorithm \ref{alg:F2BA}). After $K$-steps of gradient descent on variable $y$ and $z$, one can expect that it holds $y_t^K \approx y_{\lambda}^*(x_t)$ and $z_t^K \approx y^*(x_t)$. 
% Then, an estimator of $ \nabla \fL_{\lambda}^*(x_t)$ is constructed according to \Eqref{eq:nabla-Lag} (Line 9 of Algorithm \ref{alg:F2BA}).
% Next, the algorithm performs a gradient step on $x$ via this estimator (Line 10).
% Finally, after $T= \fO(\epsilon^{-2})$ steps of gradient descent on $x$, 
% this algorithm can output an $\epsilon$-stationary point of $\fL_{\lambda}^*(x)$.
 

\begin{lem} \label{lem:1st}
Suppose Assumption \ref{asm:sc} holds. Define $\ell$,  $\kappa$ according to Definition \ref{dfn:kappa-sc}, and $\fL_{\lambda}^*(x)$ according to \Eqref{eq:L-lambda}. Set $\lambda \ge 2L_f /\mu $, then it holds that 
\begin{enumerate}[label=\alph*.]
    \item $\Vert \nabla \fL_{\lambda}^*(x) - \nabla \varphi(x) \Vert = \fO(\ell \kappa^3 /\lambda),~\forall x \in \BR^{d_x}$ {\rm (Lemma \ref{lem:key-fully})}.
    \item $\vert \fL_{\lambda}^*(x) - \varphi(x) \vert =\fO(\ell \kappa^2/ \lambda),~\forall x \in \BR^{d_x}$ {\rm (Lemma \ref{lem:F-Lag-close})}.
    \item $\fL_{\lambda}^*(x) $ is $\fO(\ell \kappa^3)$-gradient Lipschitz {\rm (Lemma \ref{lem:nabla2-bound})}.
\end{enumerate}
\end{lem}

All the formal versions of these lemmas and the corresponding proofs can be found in Appendix \ref{apx:lem-1st}.
Lemma \ref{lem:1st}a is a restatement of Lemma 3.1 by \citet{kwon2023fully}, which demonstrates that when $\lambda \asymp \epsilon^{-1}$, an $\epsilon$-first-order stationary point of $\fL_{\lambda}^*(x)$ is also an $\fO(\epsilon)$-first-order stationary of $\varphi(x)$.
Lemma \ref{lem:1st}c is a new result proved in this paper. It means although $\fL_{\lambda}^*(x)$ depends on $\lambda$, when $\lambda$ exceeds a certain threshold, the gradient Lipschitz coefficient of $\fL_{\lambda}^*(x)$ only depends on that of $\varphi(x)$ and does not depend on $\lambda$.  

Since the convergence rate of gradient descent depends on the gradient Lipschitz coefficient of the objective, Lemma \ref{lem:1st}c indicates that optimizing $\fL_{\lambda}^*(x)$ is as easy as optimizing $\varphi(x)$. 
Note that
$\nabla \fL_{\lambda}^*(x)$ only involves first-order information (\Eqref{eq:nabla-Lag}), Lemma \ref{lem:1st}c then tells that first-order methods can have the same convergence rate as second-order methods, as stated in the following theorem.

% Mathematically speaking, it holds true that $ \Vert \nabla^2 \fL_{\lambda}^*(x) \Vert \asymp \Vert \nabla^2 \varphi(x) \Vert \asymp \ell \kappa^3$.
% % Therefore, the gradient descent on $\fL_{\lambda}^*(x)$ has the same convergence rate as the gradient descent on $\varphi(x)$. 
% By standard analysis of gradient descent, one can show that
% if $\Vert \hat \nabla \fL_{\lambda}^*(x_t) - \nabla \fL_{\lambda}^*(x_t) \Vert = \fO(\epsilon)$ holds for all $t$, then the algorithm can find an $\epsilon$-first-order stationary point of $\varphi(x)$ with no more than $T = \fO(\ell \kappa^3 \epsilon^{-2})$ outer loop iterations.

% % The good news is that $\nabla \fL_{\lambda}^*(x)$ can be expressed only using first-order information (\Eqref{eq:nabla-Lag}), while $\nabla \varphi(x)$ involves second-order information. 

% Now it remains to analyze the complexity to ensure $\Vert \hat \nabla \fL_{\lambda}^*(x_t) - \nabla \fL_{\lambda}^*(x_t) \Vert = \fO(\epsilon)$. Recalling the definition of $y_{\lambda}^*(x)$ in \Eqref{eq:L-lambda} and the expression of $\nabla \fL_{\lambda}^*(x)$ in
% \Eqref{eq:nabla-Lag}, it suffices to maintain $ \max \{ \Vert y_t^K - y_{\lambda}^*(x_t) \Vert,~\Vert z_t^K  - y^*(x_t) \Vert \} = \fO(\nicefrac{\epsilon^2}{\lambda}  )  $ for all $t$.
% Since the objective in \Eqref{eq:opt-xyz} is strongly convex in both $y$ and $z$, the complexity in the inner loop can be bounded by $K = \tilde \fO(\kappa)$ first-order oracle calls.

% %Note that both the evaluation of $y_t^K \approx y_{\lambda}^*(x_t)$ and $z_t^K \approx y^*(x_t)$ can be done efficiently since the objective in \Eqref{eq:opt-xyz} is strongly convex in both $y$ and $z$.
% According to the above analysis, we
% can formally prove the theoretical guarantee for Algorithm \ref{alg:F2BA} as follows.

% These lemmas allow us to prove the following theorem, which improves the upper complexity bound of F${}^2$BA from previous $\tilde \fO(\lambda \epsilon^{-2})$ to $ \tilde \fO( \log(\lambda) \epsilon^{-2})$.
% As we have discussed,
% the key to our improvement comes from 
% Lemma \ref{lem:1st}c.
\begin{restatable}{thm}{thmFBA}\label{thm:F2BA}
Suppose Assumption \ref{asm:sc} holds. Define
$\Delta:=\varphi(x_0) - \inf_{x \in \BR^{d_x}} \varphi(x)$ and $R := \Vert y_0 - y^*(x_0) \Vert^2 $.
Let $\eta \asymp \ell^{-1}\kappa^{-3}$, $\lambda \asymp \max\left\{\kappa/R,~ \ell \kappa^2/ \Delta, ~ \ell \kappa^3 / \epsilon \right\} $
and set other parameters 
in Algorithm \ref{alg:F2BA} as
\begin{align*}
     \alpha = \frac{1}{L_g}, ~ \tau = \frac{1}{2 \lambda L_g},~ K = \fO\left( \frac{L_g }{\mu} \log \left( \frac{\lambda L_g}{\mu} \right) \right),
\end{align*}
then it
can find an $\epsilon$-first-order stationary point of $\varphi(x)$ within $ T = \fO( \ell \kappa^4 \epsilon^{-2} \log(\nicefrac{\ell \kappa}{\epsilon})) $ first-order oracle calls, where $\ell$,  $\kappa$ are defined in Definition \ref{dfn:kappa-sc}.
\end{restatable}

% xxxxx.

% \begin{restatable}{thm}{thmNFBA}\label{thm:NF2BA}
% Define $\Delta:=\varphi(x_0) - \inf_{x \in \BR^{d_x}} \varphi(x)$, $R := \Vert y_0 - y^*(x_0) \Vert^2 $.
% Let $\eta \asymp \epsilon \ell^{-1}\kappa^{-3}$, $\lambda \asymp \max\left\{\kappa/R,~ \ell \kappa^2/ \Delta, ~ \ell \kappa^3 / \epsilon \right\} $
% and set other parameters 
% in Algorithm \ref{alg:F2BA} as
% \begin{align*}
%      \alpha = \frac{1}{L_g}, ~ \tau = \frac{1}{2 \lambda L_g},~ K_0 = \fO \left( \frac{L_g}{\mu} \log  \right)  = \fO\left( \frac{L_g }{\mu} \log \left( \frac{\lambda L_g}{\mu} \right) \right).
% \end{align*}
% Then Algorithm \ref{alg:F2BA} with Option II
% can find an $\epsilon$-first-order stationary point of $\varphi(x)$ within $ T = \fO( \ell \kappa^4 \epsilon^{-2} \log(\nicefrac{\ell \kappa}{\epsilon})) $ first-order oracle calls.
% \end{restatable}

\begin{rmk} 
When the upper-level function only depends on $x$, \textit{i.e.}
we have $f(x,y) \equiv h(x)$ for some function $h(\,\cdot\,)$, the bilevel problem reduces to a single-level problem, for which \citet{carmon2020lower} proved a lower complexity bound of $\Omega(\epsilon^{-2})$. Therefore, we can conclude that the
first-order oracle complexity of {F${}^2$BA} we proved is near-optimal. 
\end{rmk}

We defer the proof of Theorem \ref{thm:F2BA} to Appendix \ref{apx:1st}. The complexity of F${}^2$BA in Theorem \ref{thm:F2BA}
achieves the near-optimal rate in the dependency on $\epsilon$, and matches the state-of-the-art second-order methods AID and ITD ~\citep{ji2021bilevel} in the dependency of $\kappa$. Therefore,
our result, for the first time, closes the gap between first-order and second-order methods for nonconvex-strongly-convex bilevel optimization.

\begin{rmk}
The complexity of {F${}^2$BA} may be further improved to $\tilde \fO(\ell \kappa^{3.5} \epsilon^{-2})$ if we use accelerated gradient descent~\citep{nesterov2018lectures} to replace gradient descent in the inner loop. We leave that in future works.
\end{rmk}

\section{Finding Second-Order Stationarity} \label{sec:second}

\begin{algorithm*}[t]  
\caption{Perturbed F${}^2$BA$\left(x_0,y_0, \eta,\tau,\alpha, \lambda, T,\{K_t\}_{t=0}^{T-1}, \epsilon,r ,\fT \right)$} \label{alg:PF2BA}
\begin{algorithmic}[1] 
\STATE $ z_0 = y_0$ \\[1mm]
\STATE \textbf{for} $ t =0,1,\cdots,T-1 $ \\[1mm]
\STATE \quad $ y_t^0 =  y_{t}, ~ z_t^0 = z_{t}$ \\[1mm]
\STATE \quad \textbf{for} $ k =0,1,\cdots,K_t-1$ \\[1mm]
\STATE \quad \quad $ z_t^{k+1} = z_t^{k}- \alpha \nabla_y g(x_t, z_t^k)   $ \\[1mm]
\STATE \quad \quad $ y_t^{k+1} = y_t^k - \tau \left( \nabla_y f(x_t,y_t^k) +  \lambda  \nabla_y g(x_t,y_t^k) \right)$ \\[1mm]
\STATE \quad \textbf{end for} \\[1mm]
% \STATE \quad $z_{t+1} = z_t^K,~ y_{t+1} = y_t^K $ \\[1mm]
\STATE \quad $ \hat \nabla \fL_{\lambda}^*(x_t)= \nabla_x f(x_t,y_{t}^{K_t}) + \lambda ( \nabla_x g(x_t,y_{t}^{K_t}) - \nabla_x g(x_t,z_{t}^{K_t}) )$ \\[1mm]
\STATE \quad \textbf{if} $ \Vert \hat \nabla \fL_{\lambda}^*(x_t) \Vert \le \frac{4}{5} \epsilon$ \textbf{and} no perturbation added in the last $\fT  $ steps \\[1mm]
\STATE \quad \quad $x_{t} = x_t - \eta \xi_t $, where $\xi_t \sim \sB(r)$  \\[1mm]
\STATE \quad \textbf{end if} \\[1mm]
\STATE \quad $x_{t+1} = x_t -  \eta  \hat \nabla \fL_{\lambda}^*(x_t)$ \\[1mm]
\STATE \textbf{end for} \\[1mm]
\end{algorithmic}
\end{algorithm*}

We have shown in the previous section that the F${}^2$BA is near-optimal for finding first-order stationary points. However, a first-order stationary point may be a saddle point or a local maximizer, which needs to be escaped from for an effective optimizer.
For this reason, many works aim to find a second-order stationary point (Definition \ref{dfn:sta-hyper-second}).
In this section, we propose a simple variant of F${}^2$BA (Algorithm \ref{alg:PF2BA}) that can achieve this higher goal. 
The only difference to Algorithm \ref{alg:F2BA} is the additional Line 9-10 in Algorithm \ref{alg:PF2BA}, which is motivated by the perturbed strategy for escaping saddle points~\citep{jin2017escape}.

% Recent advances in bilevel optimization~\citep{huang2022efficiently} show that the AID method can provably find a second-order stationary point for nonconvex-strongly-convex bilevel optimization. 
% It is natural to ask whether first-order methods can also achieve this goal as second-order methods. In this section, we give an affirmative answer to this question by proposing the perturbed F${}^2$BA in Algorithm \ref{alg:PF2BA}, which uses the perturbed strategy proposed by \citet{jin2017escape} to avoid getting stuck in saddle points.

To prove the desired conclusion, we need to extend the analysis in Lemma \ref{lem:1st}
to higher-order derivatives. Below, we show that once $\lambda$ is sufficiently large, $\fL_{\lambda}^*(x)$ and $\varphi(x)$ have not only very close gradients (Lemma \ref{lem:1st}a) but also very close Hessian matrices. 

% In this section, we show that under the additional smoothness conditions (Assumption \ref{asm:third}), the perturbed F${}^2$BA (Algorithm \ref{alg:PF2BA}) can provably converge to an approximate second-order stationary point of $\varphi(x)$. 
% Our analysis relies on the following lemmas, which characterize the higher-order connections between  $\fL_{\lambda}^*(x)$ and $\varphi(x)$ than Lemma \ref{lem:1st}.

\begin{lem} \label{lem:2rd}
Suppose both Assumption \ref{asm:sc} and \ref{asm:third} hold. Define $\ell$,  $\kappa$ according to Definition \ref{dfn:kappa-third}, and $\fL_{\lambda}^*(x)$ according to \Eqref{eq:L-lambda}. Set $\lambda \ge 2L_f /\mu $, then it holds that 
\begin{enumerate}[label=\alph*.]
    \item $\Vert \nabla^2 \fL_{\lambda}^*(x) - \nabla^2 \varphi(x) \Vert = \fO(\ell \kappa^6 /\lambda),~\forall x \in \BR^{d_x}$. {\rm (Lemma \ref{lem:nabla2-phi})}
    \item $\fL_{\lambda}^*(x) $ is $\fO(\ell \kappa^5)$-Hessian Lipschitz. {\rm (Lemma \ref{lem:nabla3-bound})}
\end{enumerate}
\end{lem}

All the formal versions of these lemmas and the corresponding proof can be found in Appendix \ref{apx:lem-2nd}. Based on these lemmas, we can prove the convergence of perturbed F${}^2$BA in the following theorem. Due to some technical reasons for controlling the approximation error of $\Vert y_t^0 - y_{\lambda}^*(x_t) \Vert^2 + \Vert z_t^0 - y^*(x_t) \Vert^2$, we use a dynamic sequence $\{K_t\}_{t=0}^{T-1}$ that depends on \Eqref{eq:delta-t}, instead of the constant $K_t \equiv K$ used in Algorithm \ref{alg:F2BA}.

% \citet{jin2017escape} showed that 
% the perturbed gradient descent can escape saddle points efficiently for a  Hessian Lipschitz function.
% It
% motivates us to propose Algorithm \ref{alg:PF2BA}, which applies 
% the inexact perturbed gradient descent on $\fL_{\lambda}^*(x)$, and gives the following guarantee~\citep{huang2022efficiently}.

% \begin{restatable}{prop}{propPGD} \label{prop:iPGD}
% Given $\lambda >0$.
% Suppose $\fL_{\lambda}^*(x)$ is $L$-gradient Lipschitz and $\rho$-Hessian Lipschitz. Set the parameters in Algorithm \ref{alg:PF2BA} as 
% \begin{align} \label{cond:iota}
%     \eta = \frac{1}{L}, \quad r = \frac{\epsilon}{400 \iota^3}, \quad \fT = \frac{L}{\sqrt{\rho \epsilon}} \cdot \iota ~~{\rm with }~~ \iota \ge 1 ~~ {\rm and} ~~ \delta \ge  \frac{L \sqrt{d}}{\sqrt{\rho \epsilon}} \iota^2 2^{8-\iota}.
% \end{align}
% Once the following inequality holds in each
% iteration:
% \begin{align} \label{cond:PF2BA}
%     \Vert \hat \nabla \fL_{\lambda}^*(x_t) - \nabla \fL_{\lambda}^*(x_t) \Vert \le \min \left\{ 
%     \frac{1}{20 \iota^2},  \frac{1}{16 \iota^2 2^{\iota}} \right\} \epsilon,
% \end{align}
% then Algorithm \ref{alg:PF2BA} can find a point $x_t$ satisfying
% \begin{align*}
%     \Vert \nabla \fL_{\lambda}^*(x_t) \Vert \le \epsilon, \quad  \nabla^2 \fL_{\lambda}^*(x_t)  \succeq - \sqrt{\rho \epsilon}~ I_{d_x}
% \end{align*}
% within $T= \fO \left( {\iota^4 \tilde \Delta L}{\epsilon^{-2}}   \right)$ iterations with probability at least $1-\delta$, where
% $\tilde \Delta = \fL_{\lambda}^*(x_0) - \inf_{x \in \BR^{d_x}} \fL_{\lambda}^*(x)$.    
% \end{restatable} 

% Note that (\ref{cond:iota}) can be satisfied by letting
% $\iota = \fO\left( \log \left(\nicefrac{L\sqrt{d}}{\sqrt{\rho \epsilon}\delta}\right)\right)$ and 
% (\ref{cond:PF2BA}) can be fulfilled within $\fO(\log(\nicefrac{\kappa}{\epsilon}))$ oracle complexity by Proposition \ref{prop:error}.
% Therefore, Algorithm \ref{alg:PF2BA} can provably find an $\epsilon$-second-order stationary point of $\fL_{\lambda}^*(x)$ within $\tilde \fO\left(\tilde \Delta \kappa L \epsilon^{-2} \right)$ first-order oracle complexity.
% According to Lemma \ref{lem:1st} and Lemma \ref{lem:2rd}, for a sufficiently large $\lambda$, we have $\tilde \Delta \asymp \Delta$, $L= \fO(\kappa^3)$, $\rho = \fO(\kappa^5) $ and the output is also an $\epsilon$-second-order stationary point of $\varphi(x)$, formally stated as follows. 

\begin{restatable}{thm}{thmPFBA}\label{thm:PF2BA}
Suppose both Assumption \ref{asm:sc} and \ref{asm:third} hold.
Define $\Delta:=\varphi(x_0) - \inf_{x \in \BR^{d_x}} \varphi(x)$ and $R := \Vert y_0 - y^*(x_0) \Vert^2 $.
Let $\eta \asymp \ell^{-1} \kappa^{-3} $, $\lambda \asymp \max\left\{ \kappa/R,~\ell \kappa^2/ \Delta, ~ \ell \kappa^3 / \epsilon,~\ell \kappa^6 / \sqrt{\epsilon} \right\} $ and set other parameters 
in Algorithm \ref{alg:F2BA} as 
\begin{align*}
    \alpha = \frac{1}{L_g}, ~ \tau = \frac{1}{2 \lambda L_g},~ K_t = \tilde \fO\left(\frac{L_g  \log \delta_t}{\mu}   \right),
\end{align*}
where $\delta_t$ is defined via the recursion
\begin{align} \label{eq:delta-t}
    \delta_{t+1} = \frac{1}{2}\delta_t + \frac{34 L_g^2}{\mu^2}  \Vert x_{t+1} - x_t \Vert^2, \quad \delta_0 = \fO(R),
\end{align}
then 
it can find an $\epsilon$-second-order stationary point of $\varphi(x)$
within $ T = \tilde \fO\left(\ell \kappa^4 \epsilon^{-2} \right) $ first-order oracle calls, where $\ell$,  $\kappa$ are defined in Definition \ref{dfn:kappa-third} and the notation $\tilde \fO(\,\cdot\,)$ hides logarithmic factors of $d_x,\kappa,\ell$, and $\epsilon$.
\end{restatable}
We defer the proof to Appendix \ref{apx:2nd}.
The above complexity 
for finding $\epsilon$-second-order stationary points matches that  
for finding $\epsilon$-first-order stationary points (Theorem~\ref{thm:F2BA}), up to logarithmic factors. Therefore, we conclude that F${}^2$BA 
can escape saddle points almost for free by simply adding some small perturbation in each step.

\begin{rmk}
Although the $\tilde \fO(\epsilon^{-2})$ complexity of {F${}^2$BA} is near-optimal for finding $\epsilon$-first-order stationary points under Assumption \ref{asm:sc},  the $\tilde \fO(\epsilon^{-2})$ complexity of perturbed {F${}^2$BA} may be further improved to $\tilde \fO(\epsilon^{-1.75})$ by applying the restarted accelerated gradient descent ~\citep{li2022restarted,yang2023accelerating} in 
the outer loop.
The acceleration is possible because the additional Assumption \ref{asm:third} implies the Hessian smoothness of $\varphi(x)$. And we also leave that in future works.
\end{rmk}


\section{Extension to the Distributed Scenario} \label{sec:distributed}

\begin{algorithm*}[t]  
\caption{Distributed F${}^2$BA $\left(\bar x_0,\bar y_0, \eta,\tau,\alpha, \lambda, T,K\right)$} \label{alg:CF2BA}
\begin{algorithmic}[1] 
\STATE $ X_0 = \vone  \bar x_0, ~ Y_0 = \vone  \bar y_0,~ Z_0 = \vone  \bar y_0$ \\[1mm]
\STATE \textbf{for} $ t =0,1,\cdots,T-1 $ \\[1mm]
\STATE \quad $ Y_t^0 =  Y_{t}, ~ Z_t^0(i) = Z_{t}$ \\[1mm]
\STATE \quad \textbf{for} $ k =0,1,\cdots,K-1$ \\[1mm]
\STATE \quad \quad  $ V_t^{k} = \nabla_y \vg(X_t, Z_t^k)$, \quad  $U_t^{k} =  \nabla_y \vf(X_t,Y_t^k) +  \lambda  \nabla_y \vg(X_t,Y_t^k) $ \\[1mm]
\STATE \quad \quad Aggregate and broadcast $\bar v_t^{k} = \frac{1}{m} \vone \vone^\top V_t^k, $~ $\bar u_t^{k} = \frac{1}{m} \vone   \vone^\top U_t^k$ \\[1mm]
\STATE \quad \quad  $ Z_t^{k+1} = Z_t^k - \alpha \vone   \bar v_t^k   $ ,\quad  $ Y_t^{k+1} = Y_t^k - \tau \vone \bar u_t^k   $ \\[1mm]
\STATE \quad \textbf{end for} \\[1mm]
% \STATE \quad $Z_{t+1} = Z_t^K,~ Y_{t+1} = Y_t^K $ \\[1mm]
\STATE \quad $ H_t= \nabla_x \vf(X_t,Y_{t}^K) + \lambda ( \nabla_x \vg(X_t,Y_{t}^K) - \nabla_x \vg(X_t,Z_{t}^K) )$ \\[1mm]
\STATE \quad Aggregate and broadcast $\bar h_t = \frac{1}{m} \vone   \vone^\top H_t$ \\[1mm]
\STATE \quad $X_{t+1} = X_t -  \eta  \vone  \bar h_t$ \\[1mm]
\STATE \textbf{end for} \\[1mm]
\end{algorithmic}
\end{algorithm*}


Some applications of bilevel optimization like neural architecture search suffer from a heavy computational burden, so the experiments are usually conducted in a distributed way~\citep{chen2019progressive}.
In this section, we extend
F${}^2$BA to the distributed setting, where both the upper and lower-level functions are distributed on $m$ agents as \Eqref{prob:fn}. The detailed procedure is presented in Algorithm \ref{alg:CF2BA}.

The classic paradigm for distributed algorithms is to compute the local gradients on each agent in parallel, and then aggregate them on the server. 
However, challenges arise when extending existing second-order methods for bilevel optimization from the single-machine setting to the distributed setting. Given a variable $\bar x$, one may want to calculate its hyper-gradient (\Eqref{hyper-grad}) according to this paradigm by:
\begin{align*}
    \hat \nabla \varphi(\bar x) = {\bf \Pi} \left( \nabla_x \vf(\vone x, \vone y^*(\bar x)) + \nabla_{xy}^2 \vf(\vone x,\vone y^*(\bar x)) [\nabla_{yy}^2 \vg(\vone \bar x, \vone y^*(\bar x))]^{-1} \nabla_y \vf(\vone \bar x, \vone y^*(\bar x))  \right),
\end{align*}
where $y^*(x) := \arg \min_{y \in \BR^{d_y}} g(x,y)$ and ${\bf \Pi} := \frac{1}{m} \vone \vone^\top$ denotes the aggregation operator on the server. 
But the nested structure of the hyper-gradient indicates that $ \hat \nabla \varphi(\bar x) \ne \nabla \varphi(\bar x)$. As a consequence, 
researchers need to pay extra effort to make second-order methods work in distributed bilevel problems~\citep{tarzanagh2022fednest,chen2022decentralized}.
It is a non-trivial issue to address, especially when the operator $ {\bf \Pi} \nabla_{yy}^2 \vg(\vone \bar x,\vone y^*(\bar x)) $ is forbidden to avoid an unacceptable $\fO(d_y^2)$ communication complexity.

In contrast, the F${}^2$BA naturally avoid this challenge since
\begin{align*}
    \nabla \fL_{\lambda}^*(\bar x) = {\bf \Pi} \left( \nabla_x \vf(\vone x,\vone y_{\lambda}^*(\bar x)) + \lambda (\nabla_x \vg(\vone \bar x, \vone y_{\lambda}^*(\bar x) ) - \nabla_x \vg (\vone \bar x, \vone y^*(\bar x)) ) \right),
\end{align*}
where $y_{\lambda}^*(x) := \arg \min_{y \in \BR^{d_y}} f(x,y) + \lambda g(x,y)$. It means that the previously mentioned classic aggregation paradigm for distributed optimization directly works for F${}^2$BA without any additional modification in the algorithm, as stated below.

% The most natural idea is to aggregate the local gradients in \Eqref{eq:nabla-Lag} on all agents. This procedure is presented in 
% Algorithm \ref{alg:CF2BA}.
% Thanks to the form of \Eqref{eq:nabla-Lag}, one can easily verify that the mean variables in Algorithm \ref{alg:CF2BA} follow the same trajectories as in Algorithm \ref{alg:F2BA}.
\begin{prop} 
Running Algorithm \ref{alg:CF2BA} is equivalent to running  Algorithm \ref{alg:F2BA}  on the mean variables
\begin{align*}
    \bar x = \frac{1}{m} \vone^\top X,~ \bar y = \frac{1}{m} \vone^\top Y ~{\rm and}~ \bar z = \frac{1}{m} \vone^\top Z.
\end{align*}
\end{prop}

Then the convergence of Algorithm \ref{alg:CF2BA} directly follows Theorem \ref{thm:F2BA}.
%that for F${}^2$BA distributed bilevel problems are as easy as the single-machine setup.  As we have discussed, such an advantage does not hold for second-order methods.
\begin{cor} \label{cor:CF2BA}
Suppose Assumption \ref{asm:sc} holds.
Algorithm \ref{alg:CF2BA} with the same parameters in Theorem \ref{thm:F2BA} can find  $X_t$ satisfying $ \Vert \nabla \varphi(\bar x_t) \Vert \le \epsilon $ within $\fO(\ell \kappa^4 \epsilon^{-2} \log(\nicefrac{\ell \kappa}{\epsilon})) $ first-order oracle calls and  $\fO( \ell \kappa^4 \epsilon^{-2} \log(\nicefrac{\ell \kappa}{\epsilon})) $ communication rounds.
\end{cor}

Algorithm \ref{alg:CF2BA} is a near-optimal distributed algorithm since both the computation and communication complexity match the $\Omega(\epsilon^{-2})$ lower bound (Theorem 1 by \citet{lu2021optimal}), up to logarithmic factors.
Compared with the second-order methods, the distributed F${}^2$BA is more
practical since it
neither requires a $\fO(d_y^2)$ communication complexity per iteration~\citep{yangdecentralized} nor an additional distributed sub-solver for matrix inverse~\citep{tarzanagh2022fednest,chen2022decentralized}. 

\begin{rmk}
Similarly, we can also implement perturbed {F${}^2$BA} (Algorithm \ref{alg:PF2BA}) under the distributed setting. This would lead to a distributed first-order algorithm that can provably find an $\epsilon$-second-order stationary point with $\tilde \fO(\epsilon^{-2})$ complexity.
\end{rmk}


% \begin{algorithm*}[t]  
% \caption{Decentralized F${}^2$BA $\left(\bar x_0, \bar y_0, \eta,\tau,\alpha, \lambda, T,\{K_t\}_{t=0}^{T-1}, S_{\rm in}, S_{\rm out} \right)$} \label{alg:DF2BA}
% \begin{algorithmic}[1] 
% \STATE $ X_0 = \vone \bar x_0,~ Y_0 = \vone \bar y_0,~ Z_0 = \vone \bar y_0$\\[1mm]
% \STATE $H_{-1} = \vzero, ~ \nabla \fL_{\lambda}^*(X_{-1})= \vzero$ \\[1mm]
% \STATE \textbf{for} $ t =0,1,\cdots,T-1 $ \\[1mm]
% \STATE \quad $ Z_t^0 =  Z_{t},~ V_t^{-1}= \vzero,~ \nabla_y g(X_t,Z_t^{-1}) = \vzero $ \\[1mm]
% \STATE \quad $ Y_t^0 =  Y_{t},~ U_t^{-1}= \vzero,~ \nabla_y g(X_t,Y_t^{-1}) = \vzero $
% \STATE \quad \textbf{for} $ k =0,1,\cdots,K_t-1$ \\[1mm]
% \STATE \quad \quad $V_t^{k} =
% W^{S_{\rm in}} \cdot \left(
% V_t^{k-1} + \nabla_y g(X_t,Z_t^{k}) - \nabla_y g(X_t,Z_t^{k-1}) \right)$ \\[1mm]
% \STATE \quad \quad $ U_t^{k} ~=
% W^{S_{\rm in}} \cdot  \left(U_t^{k-1} + \nabla_y f(X_t,Y_{t}^{k}) - \nabla_y f(X_t, Y_t^{k-1}) + \lambda \left(  \nabla_y g(X_t,Y_{t}^{k}) - \nabla_y g(X_t, Y_t^{k-1}) \right) \right)$ \\[1mm]
% \STATE \quad \quad $ Z_t^{k+1} = W^{S_{\rm in}} \cdot \left(Z_t^{k}- \alpha V_t^k \right)  $ \\[1mm]
% \STATE \quad \quad $Y_{t}^{k+1} = W^{S_{\rm in}} \cdot \left(Y_t^k - \tau U_t^k\right)$ \\[1mm]
% \STATE \quad \textbf{end for} \\[1mm]
% \STATE \quad $Z_{t+1} = Z_t^K,~ Y_{t+1} = Y_t^K $ \\[1mm]
% \STATE \quad $ \hat \nabla \fL_{\lambda}^*(X_t)= \nabla_x f(X_t,Y_{t+1}) + \lambda ( \nabla_x g(X_t,Y_{t+1}) - \nabla_x g(X_t,Z_{t+1}) )$ \\[1mm]
% \STATE \quad $H_t = W^{S_{\rm out}} \cdot \left(H_{t-1} + \hat \fL_{\lambda}^*(X_t) - \hat \fL_{\lambda}^*(X_{t-1})  \right)$ \\[1mm]
% \STATE \quad $X_{t+1} = 
% W^{S_{\rm out}} \cdot \left(
% X_t -  \eta  H_t\right)$ \\[1mm]
% \STATE \textbf{end for} \\[1mm]
% \end{algorithmic}
% \end{algorithm*}


% When the central server is absent, we apply the standard gradient tracking technique~\citep{qu2017harnessing,xu2015augmented,di2016next,nedic2017achieving}
% to achieve the same convergence rate as the centralized setting. The proposed decentralized F${}^2$BA are (Algorithm \ref{alg:DF2BA}) can be viewed as performing an inexact gradient descent with gradient tracking on $\fL_{\lambda}^*(x)$, which has the following convergence result.

% \begin{restatable}{prop}{propGTGD}
%     Given $\lambda >0$. Suppose $\fL_{\lambda}^*(x)$ is $L$-gradient Lipschitz.  Set $\eta = 1/(2L)$. As long as the following inequality holds in each
% iteration:
% \begin{align} \label{cond:DF2BA}
%     \Vert \hat \nabla \fL_{\lambda}^*(X_t) - \nabla \fL_{\lambda}^*(X_t) \Vert_F \le xxx,
% \end{align} 
% Algorithm \ref{alg:DF2BA} can find a point 
% $x_t$
% satisfying $ \Vert \nabla \fL_{\lambda}^*(\bar x_t) \Vert \le \epsilon$ within $T = \fO(\Delta L \epsilon^{-2})$ iterations, where $\Delta = \fL_{\lambda}^*(x_0) - \inf_{x \in \BR^{d_x}} \fL_{\lambda}^*(x)$.
% \end{restatable}

% Analogous to Proposition \ref{prop:error}, Condition (\ref{cond:DF2BA}) can also be achieved within $\fO(\nicefrac{1}{\epsilon})$ complexity.

% \begin{restatable}{prop}{propDerror} \label{prop:Derror}
% Given $\lambda \ge 2L_f / \mu$. Suppose Assumption \ref{asm:sc} and \ref{asm:W} holds, running Algorithm \ref{alg:DF2BA} with a given outer-loop stepsize $\eta$ and 
% \begin{align*}
%     xxx, ~{\rm and} \\
% R_t =
% \begin{cases}
%     \Vert y_0 - y^*(x_0) \Vert + \dfrac{C_f}{\lambda \mu}, & t=0; \\
%     \dfrac{\zeta}{2 \lambda L_g} + \dfrac{4 L_g \eta }{\mu} \left( \zeta +   \left( 1 + \dfrac{L_g}{\mu}\right) C_f\right), & t \ge 1;
% \end{cases}
% \end{align*}
% fulfills the goal of $ \Vert \hat \nabla \fL_{\lambda}^*(x_t) - \nabla \fL_{\lambda}^*(x_t) \Vert \le \zeta $ for all $t$.
% \end{restatable}

% By combing Lemma \ref{lem:1st} and the above two propositions, we have the following result.

% \begin{cor}
% Algorithm \ref{alg:DF2BA} with appropriate parameters can find  $X_t$ satisfying $ \Vert \nabla F(\bar x_t) \Vert \le \epsilon $ within $\tilde \fO( \kappa^4 \epsilon^{-2}) $ first-order oracle complexity and  $\tilde \fO( \chi \kappa^4 \epsilon^{-2}) $ communication complexity.
% \end{cor}

% We remark that the communication complexity can be further improved by applying the Chebyshev acceleration~\citep{scaman2017optimal,alghunaim2020decentralized}, which constructs $p(W)$ that is a polynomial of $W$ with degree $\fO(\sqrt{\chi})$. The constructed $p(W)$ has $\fO(1)$ characteristic number and thus replacing it with the original $W$ in Algorithm \ref{alg:DF2BA} leads to the optimal $\fO(\sqrt{\chi} \epsilon^{-2})$ communication complexity~\citep{lu2021optimal}.
 
\section{Numerical Experiments}

We conduct experiments to compare the proposed first-order method F${}^2$BA with existing second-order methods under both the single-machine and distributed settings.
We implement all the algorithms with PyTorch~\citep{paszke2019pytorch} and conduct the experiments on a single NVIDIA A40 GPU.

\subsection{Single-Machine Data Hyper-Cleaning} \label{sec:exp-clean}

The task of data hyper-cleaning~\citep{franceschi2018bilevel,shaban2019truncated,zhou2022model} considers finding the optimal weights of samples on a corrupted training set $\fD^{\rm tr}$, such that models trained on the weighted training set have good performances on the validation set $\fD^{\rm val}$. 
Let $(a_i,b_i)$ be the $i$-th sample in the dataset, where $a_i$ denotes the feature and $b_i$ denotes its corresponding label. And let $\sigma(\,\cdot\,)$ be the function that clips a scalar to the interval $[0,1]$.
The problem can be formulated as the following bilevel problem, with upper and lower-level functions given by:
\begin{align*}
    f(x,y) &:= \frac{1}{\vert \fD^{\rm val} \vert} \sum_{(a_i,b_i) \in \fD^{\rm val}} \ell(\, \langle a_i,y \rangle\,,b_i), \\ g(x,y) &:= \frac{1}{\vert \fD^{\rm tr} \vert} \sum_{(a_i,b_i) \in \fD^{\rm tr}} 
\sigma(x_i) ~\ell \left( \,\langle a_i,y \rangle\,,b_i \right) + c \Vert y \Vert^2,
\end{align*}
where $\ell(\,\cdot\,, \,\cdot\,)$ is the cross-entropy loss. We set  $c = 10^{-3}$ by following \citep{shaban2019truncated,ji2021bilevel,liu2022bome}. We corrupt $50\%$ of the training set by setting their labels as random classes by following \citep{liu2022bome}. We compare F${}^2$BA (Algorithm \ref{alg:F2BA}) with ITD~\citep{ji2021bilevel} and AID with conjugate gradient~\citep{maclaurin2015gradient} on datasets MNIST and Fashion-MNIST. 
%Since the algorithms can access the testing set, we remain an independent validation set to evaluate their performance. 
For all the algorithms, we set the number of inner loops as 10 and tune the learning rates in $ \{10^3,10^2,10^1,10^0,10^{-1},10^{-2},10^{-3} \}$. For F${}^2$BA, we additionally tune the multiplier $\lambda$ in $\{10^1,10^2,10^3\}$. 
Figure \ref{fig:cleaning} shows the comparison on
the computational time (seconds) against
the suboptimality gap in the logarithmic scale, where
the suboptimality gap is defined as the difference between the testing loss and the lowest achievable loss by running all the algorithms for 50,000 epochs.

\begin{figure*}[t]
    \centering
    \includegraphics[scale= 0.3]{fig/data_cleaning.eps} 
    \caption{
    Experiments of data hyper-cleaning on MNIST (left) and Fashion-MNIST (right).} 
    
    \label{fig:cleaning}
\end{figure*}

\begin{figure*}[t]
    \centering
    \includegraphics[scale= 0.3]{fig/l2reg.eps} 
    \caption{ Experiments of distributed learnable regularization on the 20news dataset.
    }
    \label{fig:l2reg}
\end{figure*}

\subsection{Distributed Learnable Regularization}

We also use bilevel optimization to find the optimal regularization coefficient on the 20news dataset as~\citep{ji2021bilevel,liu2022bome}. In this formulation, the upper and lower-level functions take the form of:
\begin{align*}
    f(x,y) := \frac{1}{\vert \fD^{\rm val} \vert} \sum_{(a_i,b_i) \in \fD^{\rm val}} \ell(\, \langle a_i,y \rangle\,,b_i), \quad g(x,y) := \frac{1}{\vert \fD^{\rm tr} \vert} \sum_{(a_i,b_i) \in \fD^{\rm tr}} 
\ell \left( \,\langle a_i,y \rangle\,,b_i \right) + y^\top \Sigma(x) \,y,
\end{align*}
where $\ell(\,\cdot\,, \,\cdot\,)$ is the cross-entropy loss and $\Sigma(x) := {\rm diag}(\exp(x))$ determines the weights of features. We evenly distribute the datasets 
on four agents and use the Gloo backend to implement parallel computing. We compare the distributed F${}^2$BA (Algorithm \ref{alg:CF2BA}) with AggITD~\citep{xiao2023communication} and AID with HIGD oracle~\citep{chen2022decentralized}. 
The hyperparameter selection strategy follows the same way as Section \ref{sec:exp-clean}. 
%For all the algorithms, we set the number of inner loops as  10 and tune the learning rates in $ \{1000,100,10,1,0.1,0.01,0.001 \}$.  For F${}^2$BA, we additionally tune the multiplier $\lambda$ in $\{10,100,1000\}$.
We run all the algorithms for 5,000 epochs and plot the testing loss in Figure \ref{fig:l2reg}. It can be seen that the F${}^2$BA is more efficient than second-order methods in both computation and communication.


\section{Conclusions and Future Directions}

This paper shows how to achieve the near-optimal complexity for first-order methods in nonconvex-strongly-convex bilevel optimization.
Our results close the gap between first-order and second-order methods for bilevel optimization under various setups. We conclude this paper with 
several potential directions for future research.
\paragraph{Acceleration.} As discussed in our remarks, it is possible to apply acceleration techniques to improve the complexity dependency on $\kappa$ for finding both first-order and second-order stationary points, and the dependency on $\epsilon$ for finding second-order stationary points.
\paragraph{Logarithmic factors.}
The complexity of F${}^2$BA has an additional $\log(\nicefrac{1}{\epsilon})$ factor compared with the lower bound for finding first-order stationary points~\citep{carmon2020lower}. 
Although this factor can be shaved for second-order methods using tighter analysis~\citep{ji2021bilevel},
it remains open whether it is avoidable or unavoidable for first-order methods.
\paragraph{Single-Loop Methods.} 
Our method adopts a double-loop structure. Designing single-loop methods could be more efficient in certain scenarios, but the analysis would also be more challenging~\citep{lin2020gradient,hong2023two,khanduri2021near,chen2022single,kwon2023fully,chen2021closing,chen2023optimal}.
\paragraph{Stochastic Bilevel Problems.} This paper only considers the deterministic case.  It is of great significance to study first-order methods using stochastic oracles~\citep{khanduri2021near,ji2021bilevel,kwon2023fully,huo2023new,yang2021provably} in future works.
\paragraph{Distributed Bilevel Problems.}
We believe that the distributed first-order methods would have broader application prospects than second-order methods and anticipate our algorithm could be applied to more distributed scenarios~\citep{lian2017can,alistarh2018byzantine,bernstein2018signsgd,yuan2022decentralized,wang2022fine,gandikota2021vqsgd,alistarh2018convergence,ye2018communication,scaman2017optimal,mishchenko2022proxskip}.
\paragraph{More Lower-Level Problems.}
Assumption \ref{asm:sc} is somewhat restricted. It would be important to 
study the convergence of first-order algorithms for bilevel problems under more relaxed assumptions on the lower-level problems, such as nonsmooth lower-level problems~\citep{lu2023first}, constrained lower-level problems~\citep{liu2021towards,khanduri2023linearly,tsaknakis2022implicit,xiao2023alternating}, and bilevel problems without lower-level strong convexity~\citep{sow2022constrained,liu2021towards,liu2021value,liu2020generic,chen2023bilevel,shen2023penalty,xiao2023generalized}.

\bibliographystyle{plainnat}
\bibliography{sample}

\appendix

\newpage

\section{Notations for Tensors and Derivatives} \label{apx:tensor}
We follow the notations of tensors used in \cite{kolda2009tensor}.
For a three-way tensor $\fX \in \BR^{d_1 \times d_2 \times d_3}$,
we use $ [\fX]_{i_1,i_2,i_3}$ to represent its $(i_1,i_2,i_3)$-th element.
The inner product of two three-way tensors $\fX,\fY$ is defined by $\langle \fX, \fY \rangle := \sum_{i_1,i_2,i_3} [\fX]_{i_1,i_2,i_3} \cdot [\fY]_{i_1,i_2,i_3}$. The operator norm of three-way tensor $\fX \in \BR^{d_1 \times d_2 \times d_3}$ is defined by $\Vert \fX \Vert:= \sup_{\Vert x_i \Vert =1} \langle \fX, x_1 \circ x_2 \circ x_3 \rangle $, where the elements of  $x_1 \circ x_2 \circ x_3 \in \BR^{d_1 \times d_2 \times d_3}$ is $ [x_1 \circ x_2 \circ x_3]_{i_1,i_2,i_3} := [x_1]_{i_1} \cdot  [x_2]_{i_2} \cdot [x_3]_{i_3} $. 
It can be verified that such a definition generalizes the Euclidean norm of a vector and the spectral norm of a matrix to tensors.
For a three-way tensor  $ \fX \in \BR^{d_1 \times d_2 \times d_3}$ and a vector $v \in \BR^{d_1}$, their mode-1 product, denoted by $\fX \bar \times_1 v$, gives a matrix in $\BR^{d_2 \times d_3}$ with elements $ [\fX \bar \times_1 v]_{i_2,i_3} := \sum_{i_1} [\fX]_{i_1,i_2,i_3} \cdot [v]_{i_1} $.
We define $\bar \times_2$ and $\bar \times_3$ in a similar way, and it can  be verified that 
$ \Vert \fX \bar \times_i v \Vert \le \Vert \fX \Vert \Vert v \Vert$.
For a three-way tensor $\fX \in \BR^{d_1 \times d_2 \times d_3}$ and a matrix $A \in \BR^{d_1' \times d_1}$, their mode-1 product, denoted by $ \fX \times_1 A$, gives a tensor in $\BR^{d_1' \times d_2 \times d_3}$ with elements  $ [\fX \times_1 A]_{i_1',i_2,i_3} := \sum_{i_1} [\fX]_{i_1,i_2,i_3} \cdot [A]_{i_1',i_1} $. We define $ \times_2$ and $ \times_3$ in a similar way, and it can also be verified that 
$ \Vert \fX \times_i A \Vert \le \Vert \fX \Vert \Vert A \Vert$.

For a function $h(x,y) : \BR^{d_x} \times \BR^{d_y} \rightarrow \BR$,
we denote $\nabla_x h(x,y) \in \BR^{d_x}$ to be the partial gradient with respect to $x$, with elements given by $ [\nabla_x h(x,y) ]_{i} := \partial f(x,y) / \partial [x]_{i}$. And we define $\nabla_y h(x,y) \in \BR^{d_y}$ in a similar way.  We denote $ \nabla_{xx}^2 h(x,y) \in \BR^{d_x} \times \BR^{d_x}$ to be the partial Hessian with respect to $x$, with elements given by $ [\nabla_{xx}^2 h(x,y)]_{i_1,i_2} := \partial^2 h(x,y) / (\partial [x]_{i_1} \partial [x]_{i_2})$. And we define $ \nabla_{xy}^2 h(x,y) \in \BR^{d_x} \times \BR^{d_y} $, $\nabla_{yx}^2 h(x,y) \in \BR^{d_y} \times \BR^{d_x}$ and  $ \nabla_{yy}^2 h(x,y) \in \BR^{d_y} \times \BR^{d_y}$ in a similar way.  We denote $\nabla_{xxx}^3 h(x,y) \in \BR^{d_x} \times \BR^{d_x} \times \BR^{d_x}$ to be the partial third-order derivative with respect to $x$, with elements given by 
$ [\nabla_{xxx}^3 h(x,y)]_{i_1,i_2,i_3} := \partial^3 h(x,y) / (\partial [x]_{i_1} \partial [x]_{i_2} \partial [x]_{i_3})$. And we define $\nabla_{xxy}^3 h(x,y) \in \BR^{d_x} \times \BR^{d_x} \times \BR^{d_y}$, $\nabla_{xyx}^3 h(x,y) \in \BR^{d_x} \times \BR^{d_y} \times \BR^{d_x}$,
$\nabla_{xyy}^3 h(x,y) \in \BR^{d_x} \times \BR^{d_y} \times \BR^{d_y}$,
$\nabla_{yyy}^3 h(x,y) \in \BR^{d_y} \times \BR^{d_y} \times \BR^{d_y}$,
$\nabla_{yyx}^3 h(x,y) \in \BR^{d_y} \times \BR^{d_y} \times \BR^{d_x}$,
$\nabla_{yxy}^3 h(x,y) \in \BR^{d_y} \times \BR^{d_x} \times \BR^{d_y}$
and$\nabla_{yxx}^3 h(x,y) \in \BR^{d_y} \times \BR^{d_x} \times \BR^{d_x}$
in a similar way.
We denote $\nabla y^*(x) \in \BR^{d_x} \times \BR^{d_y} $ to be the matrix with elements given by $[\nabla y^*(x) ]_{i_1,i_2} :=  \partial [y^*(x)]_{i_2}/\partial [x]_{i_1}$. We denote $\nabla^2 y^*(x) \in \BR^{d_x} \times \BR^{d_x} \times \BR^{d_y} $ to be the three-way tensor with elements given by $[\nabla^2 y^*(x)]_{i_1,i_2,i_3}: = \partial^2 [y^*(x)]_{i_3} / (\partial [x]_{i_1} \partial [x]_{i_2})$. And we define $\nabla y_{\lambda}^*(x)$ and $\nabla^2 y_{\lambda}^*(x)$ in a similar way.

\section{Lemmas for Finding First-Order Stationarity} \label{apx:lem-1st}

\begin{lem}[From \citet{kwon2023fully}] \label{lem:Lag-sc}
Under Assumption \ref{asm:sc}, for $\lambda \ge 2L_f /\mu$,   $\fL_{\lambda}(x,\,\cdot\,)$ is  $(\lambda \mu/2)$-strongly convex.  
\end{lem}
% \begin{proof}
% It can be verified that
% \begin{align*}
%     \nabla_{yy}^2 \fL(x,y) = \nabla_{yy}^2 f(x,y) +  \lambda \nabla_{yy}^2 g(x,y) \succeq \frac{\lambda \mu}{2} I_{d_y}.
% \end{align*}
% % Since  $g^*(x)$ is $(1+ L_g / \mu)L_g$-smooth, then the smoothness constant of $\fL_{\lambda}(x,y)$ can be bounded by:
% % \begin{align*}
% %     L_f + \lambda \left(2+ \frac{L_g}{\mu} \right)L_g \le \frac{4 \lambda L_g^2}{\mu^2}.
% % \end{align*}
% \end{proof}


% Then it is shown in \cite{kwon2023fully} that the following results hold. We provide the proof here for completeness.

\begin{lem} \label{lem:yy-lambda}
    Under Assumption \ref{asm:sc}, for $\lambda \ge 2L_f/ \mu$, it holds that 
    \begin{align*}
        \Vert y_{\lambda}^*(x) - y^*(x) \Vert \le \frac{C_f}{\lambda \mu}
    \end{align*}
\end{lem}

\begin{proof}
    % It is a special case of Lemma 3.2 in \cite{kwon2023fully} with $x_1 = x_2, \lambda_1 = \lambda $ and $\lambda_2 = + \infty $.
By the first-order optimality condition, we know that
\begin{align*}
    \nabla_y f(x, y_{\lambda}^*(x)) + \lambda \nabla_y g(x,y_{\lambda}^*(x))  = 0.
\end{align*}
Then we have
\begin{align*}
    \Vert y_{\lambda}^*(x) - y^*(x) \Vert 
    \le \frac{1}{\mu} \Vert \nabla_y g(x,y_{\lambda}^*(x)) \Vert = \frac{1}{\lambda \mu} \Vert \nabla_y f(x,y_{\lambda}^*(x)) \Vert \le \frac{C_f}{\lambda \mu}
\end{align*}
\end{proof}

As a direct consequence, we can show that $\fL_{x}^*$ and $\varphi(x)$ are close.
\begin{lem} \label{lem:F-Lag-close}
 Under Assumption \ref{asm:sc}, for $\lambda \ge 2L_f / \mu$, it holds that $\vert \fL_{\lambda}^*(x) - \varphi(x) \vert \le D_0/ \lambda$, where
 \begin{align*}
     D_0 : = \left(C_f + \frac{C_f L_g}{2 \mu} \right) \frac{C_f }{\mu} = \fO(\ell \kappa^2).
 \end{align*}
\end{lem}

\begin{proof}
A simple calculus shows that
    \begin{align*}
    &\quad \vert \fL_{\lambda}^*(x) - \varphi(x) \vert \\
    &\le   \vert f(x,y_{\lambda}^*(x)) - f(x,y^*(x)) \vert + \lambda \vert g(x,y_{\lambda}^*(x)) - g(x,y^*(x)) \vert \\
    &\le C_f  \Vert y_{\lambda}^*(x) - y^*(x) \Vert + \frac{\lambda L_g}{2} \Vert y_{\lambda}^*(x) - y^*(x) \Vert^2 \\
    &\le \left(C_f + \frac{C_f L_g}{2 \mu} \right) \Vert y_{\lambda}^*(x) - y^*(x) \Vert \\
    &\le \left(C_f + \frac{C_f L_g}{2 \mu} \right) \frac{C_f }{\lambda \mu}.
    \end{align*}
\end{proof}
The following result is the key to designing fully first-order methods for bilevel optimization, first proved in Lemma 3.1 in \cite{kwon2023fully}. We provide the proof here for completeness.
\begin{lem} \label{lem:key-fully}
Under Assumption \ref{asm:sc}, for $\lambda \ge 2L_f/ \mu$,  it holds that $\Vert \nabla \fL_{\lambda}^*(x) - \nabla \varphi(x) \Vert \le D_1 / \lambda$, where
\begin{align} \label{eq:D-lambda}
    D_1 :=  \left( L_f 
   + \frac{\rho_g L_g}{\mu} + \frac{C_f L_g \rho_g}{2\mu^2} + \frac{C_f \rho_g}{2\mu} 
    \right) \frac{C_f}{\mu} = \fO(\ell \kappa^3).
\end{align}
\end{lem}

\begin{proof}
Taking total derivative on $\varphi(x) = f(x,y^*(x))$, we obtain the following result:
\begin{align} \label{eq:hyper-grad}
    \nabla \varphi(x) = \nabla_x f(x,y^*(x)) - \nabla_{xy}^2 g(x,y^*(x)) [ \nabla_{yy} g(x,y^*(x))]^{-1} \nabla_y f(x,y^*(x)). 
\end{align}
Also, we know that
\begin{align*}
    \nabla \fL_{\lambda}^*(x) = \nabla_x f(x,y_{\lambda}^*(x)) + \lambda (\nabla_x g(x,y_{\lambda}^*(x)) - \nabla_x g(x,y^*(x))).
\end{align*}
By simple calculus, we have
\begin{align*}
    &\quad \nabla \fL_{\lambda}^*(x) - \nabla \varphi(x) \\
    &= \nabla_x f(x,y_{\lambda}^*(x)) - \nabla_x f(x,y^*(x)) \\
    &\quad + \nabla_{xy}^2 g(x,y^*(x)) [ \nabla_{yy} g(x,y^*(x))]^{-1} (\nabla_y f(x,y^*(x)) - \nabla_y f(x,y_{\lambda}^*(x)) \\
    &\quad + \lambda \nabla_{xy}^2 g(x,y^*(x)) [ \nabla_{yy} g(x,y^*(x))]^{-1} \times \\
    &~~~~~~~~~~~~~~~( \nabla_{yy}^2 g(x,y^*(x)) (y_{\lambda}^*(x) - y^*(x)) + \nabla_y g(x,y^*(x)) - \nabla_y g(x,y_{\lambda}^*(x)) ) \\
    &\quad + \lambda ( \nabla_x g(x,y_{\lambda}^*(x)) - \nabla_x  g(x,y^*(x)) - \nabla_{xy}^2 g(x,y^*(x)) (y_{\lambda}^*(x) - y^*(x)) ).
\end{align*}
Taking norm on both sides,
\begin{align*}
    \Vert \nabla \fL_{\lambda}^*(x) - \nabla \varphi(x) \Vert &\le L_f \Vert y_{\lambda}^*(x) - y^*(x) \Vert + \frac{\rho_g L_g}{\mu} \Vert y_{\lambda}^*(x) - y^*(x) \Vert  \\
    &\quad + \frac{\lambda L_g \rho_g}{2\mu} \Vert y_{\lambda}^*(x) - y^*(x) \Vert^2  + \frac{\lambda \rho_g}{2} \Vert y_{\lambda}^*(x) - y^*(x) \Vert^2 \\
    &\le \left( L_f 
   + \frac{\rho_g L_g}{\mu} + \frac{C_f L_g \rho_g}{2\mu^2} + \frac{C_f \rho_g}{2\mu} 
    \right) \Vert y_{\lambda}^*(x) - y^*(x) \Vert \\
    &\le \left( L_f 
   + \frac{\rho_g L_g}{\mu} + \frac{C_f L_g \rho_g}{2\mu^2} + \frac{C_f \rho_g}{2\mu} 
    \right) \frac{C_f}{\lambda \mu}
\end{align*}


\end{proof}

\begin{lem} \label{lem:yxyx-lambda}
Under Assumption \ref{asm:sc}, for $\lambda \ge 2 L_f / \mu$, it holds that $ \Vert \nabla y^*(x) - \nabla y_{\lambda}^*(x) \Vert \le D_2/\lambda $, where
\begin{align*}
    D_2 := \left( \frac{1}{\mu} + \frac{2L_g}{\mu^2} \right) \left( L_f + \frac{C_f \rho_g}{\mu} \right) = \fO(\kappa^3).
\end{align*}

\end{lem}


\begin{proof}
 Taking derivative on both sides of $\nabla_y g (x,y^*(x)) = 0$ yields
 \begin{align} \label{eq:yxyy}
     \nabla_{xy}^2 g(x,y^*(x)) +  \nabla y^*(x)\nabla_{yy}^2 g(x,y^*(x)) = 0.
 \end{align}
 Rearranging, we can obtain
\begin{align} \label{eq:y-x}
  \nabla y^*(x) = - \nabla_{xy}^2 g(x,y^*(x)) [\nabla_{yy}^2 g(x,y^*(x)) ]^{-1} . 
  \end{align}
Similarly, we also have
\begin{align} \label{eq:y-lambda-x}
    \nabla y_{\lambda}^*(x) = -  \nabla_{xy}^2 \fL_{\lambda}(x,y_{\lambda}^*(x)) [\nabla_{yy}^2 \fL_{\lambda}(x,y_{\lambda}^*(x)) ]^{-1}.
\end{align}
Using the matrix identity $ A^{-1} - B^{-1} = A^{-1} (B -A) B^{-1}$, we have
\begin{align} \label{eq:diff-inv}
    \begin{split}
    &\quad  \left \Vert [\nabla_{yy}^2 g(x,y^*(x))]^{-1} - \left [  \frac{\nabla_{yy}^2 \fL_{\lambda}(x,y_{\lambda}^*(x))}{\lambda} \right]^{-1}  \right \Vert \\
    &\le  \left \Vert [\nabla_{yy}^2 g(x,y_{\lambda}^*(x))]^{-1} \right \Vert  
    \left \Vert  \frac{\nabla_{yy}^2 \fL_{\lambda}(x,y_{\lambda}^*(x))}{\lambda} -   \nabla_{yy}^2 g(x,y^*(x))  \right \Vert \left \Vert \left[\frac{\nabla_{yy}^2 \fL_{\lambda} (x,y_{\lambda}^*(x))}{\lambda}\right]^{-1}  \right \Vert \\
    &\le \frac{2}{\mu^2} \left \Vert 
    \frac{\nabla_{yy}^2 f(x,y_{\lambda}^*(x))}{\lambda}
     + \nabla_{yy}^2 g(x,y_{\lambda}^*(x)) - \nabla_{yy}^2 g(x,y^*(x)) \right \Vert \\
     &\le \frac{2}{\mu^2} \left( \frac{L_f}{\lambda} + \rho_g \Vert y_{\lambda}^*(x) - y^*(x) \Vert  \right)\\
     &\le \frac{2}{\lambda \mu^2} \left(L_f + 
     \frac{C_f \rho_g}{\mu}
     \right). 
    \end{split}
\end{align}
Note that the setting of $\lambda \ge 2L_f / \mu  $ implies $\Vert \nabla_{xy}^2 \fL(\,\cdot\,,\,\cdot\,) \Vert \le 2 \lambda L_g$, then we further have
\begin{align*}
    &\quad \Vert \nabla y_{\lambda}^*(x) - \nabla y^*(x) \Vert \\
    & \le\left  \Vert \nabla_{xy}^2 g(x,y^*(x)) - \frac{\nabla_{xy}^2 \fL_{\lambda}(x,y_{\lambda}^*(x))}{ \lambda} \right \Vert \left \Vert [\nabla_{yy}^2 g(x,y^*(x))]^{-1}  \right \Vert  \\
    &\quad + \left \Vert \frac{\nabla_{xy}^2 \fL_{\lambda}(x,y_{\lambda}^*(x))}{\lambda}  \right \Vert  \left \Vert [\nabla_{yy}^2 g(x,y^*(x))]^{-1} - \left [  \frac{\nabla_{yy}^2 \fL_{\lambda}(x,y_{\lambda}^*(x))}{\lambda} \right]^{-1}  \right \Vert  \\
    &\le \frac{1}{\mu} \left \Vert \nabla_{xy}^2 g(x,y^*(x)) - \nabla_{xy}^2 g(x,y_{\lambda}^*(x)) - \frac{\nabla_{xy}^2 f(x,y_{\lambda}^*(x))}{\lambda}  \right \Vert  +  \frac{2L_g}{\lambda \mu^2} \left(L_f + 
     \frac{C_f \rho_g}{\mu}
     \right)  \\
    &\le \left( \frac{1}{\lambda \mu} + \frac{2L_g}{\lambda \mu^2} \right) \left( L_f + \frac{C_f \rho_g}{\mu} \right).
    \end{align*}
\end{proof}

It is clear that $\Vert \nabla y^*(x) \Vert \le L_g/\mu$, therefore $y^*(x)$ is $(L_g/\mu)$-Lipschitz. Below, we show a similar result also holds for $y_{\lambda}^*(x)$.

\begin{lem} \label{lem:nabla-y-lambda-bound}
Under Assumption \ref{asm:sc}, for $\lambda \ge 2L_f /\mu$, it holds that $\Vert \nabla y_{\lambda}^*(x) \Vert \le 4L_g/\mu$.    
\end{lem}


\begin{proof}
    Recall \Eqref{eq:y-lambda-x} that 
\begin{align*}
    \nabla y_{\lambda}^*(x) = -  \nabla_{xy}^2 \fL_{\lambda}(x,y_{\lambda}^*(x)) [\nabla_{yy}^2 \fL_{\lambda}(x,y_{\lambda}^*(x)) ]^{-1}.
\end{align*}
We arrive at the conclusion by noting 
that $  \nabla_{xy}^2 \fL_{\lambda}(\,\cdot\,,\,\cdot\,) \preceq 2 \lambda L_g$ and  $ \nabla_{yy}^2 \fL_{\lambda}(\,\cdot\, , \,\cdot\,) \succeq \lambda \mu /2 $ by Lemma \ref{lem:Lag-sc}.

\end{proof}
This implies that $y_{\lambda}^*(x)$ is $(4L_g/\mu)$-Lipschitz.

\begin{lem} \label{lem:nabla2-bound}
Under Assumption \ref{asm:sc}, for $\lambda \ge 2L_f /\mu$, $ \nabla \fL_{\lambda}^*(x)$ is $D_3$-Lipschitz, where
\begin{align*}
    D_3 &:= L_f + \frac{4L_f L_g}{\mu} +\frac{C_f \rho_g}{\mu} +  \frac{C_f L_g \rho_g}{\mu^2} + L_g D_2 = \fO(\kappa^3).
\end{align*}
\end{lem}

\begin{proof}
Note that
\begin{align} \label{eq:nabla-AB}
    \nabla \fL_{\lambda}^*(x) = \underbrace{\nabla_x f(x,y_{\lambda}^*(x))}_{A(x)} + \underbrace{\lambda ( \nabla_x g(x,y_{\lambda}^*(x)) - \nabla_x g(x,y^*(x)))}_{B(x)}. 
\end{align}
where $A(x)$ and $B(x)$ are both mappings  $\BR^{d_x} \rightarrow \BR^{d_x}$. 
From Lemma \ref{lem:nabla-y-lambda-bound} we know that $y_{\lambda}^*(x)$ is $(4L_g/\mu)$-Lipschitz. This further implies that $A(x)$ is  $(1+4L_g/\mu)L_f$-Lipschitz. Next, we bound the Lipschitz  coefficient of $B(x)$ via its derivative $\nabla B(x): \BR^{d_x} \rightarrow \BR^{d_x} \times \BR^{d_x} $, which has the following form by 
taking total derivative on $B(x)$:
\begin{align} \label{eq:nabla-Bx}
\begin{split}
    \nabla B(x) &= 
    \lambda (\nabla_{xx}^2 g(x,y_{\lambda}^*(x))-\nabla_{xx}^2 g(x,y^*(x))) \\
    &\quad + \lambda \left(\nabla y_{\lambda}^*(x)  \nabla_{yx}^2 g(x,y_{\lambda}^*(x)) - \nabla y^*(x) \nabla_{yx}^2 g(x,y^*(x)) \right).
\end{split}
\end{align}
And we can bound the operator norm of $\nabla B(x)$ by:
\begin{align*}
    \Vert \nabla B(x) \Vert &\le  \lambda \Vert \nabla_{xx}^2 g(x,y_{\lambda}^*(x)) - \nabla_{xx}^2 g(x,y^*(x)) \Vert \\
    &\quad + \lambda \Vert \nabla y^*(x) \Vert \Vert \nabla_{yx}^2 g(x,y_{\lambda}^*(x)) - \nabla_{yx}^2 g(x,y^*(x)) \Vert \\
    &\quad  + \lambda \Vert \nabla y_{\lambda}^*(x) - \nabla y^*(x) \Vert \Vert \nabla_{yx}^2 g(x,y_{\lambda}^*(x)) \Vert\\
    &\le  \lambda \rho_g \left(1+ \frac{L_g}{\mu} \right) \Vert y_{\lambda}^*(x) - y^*(x) \Vert +\lambda L_g \Vert \nabla y_{\lambda}^*(x) - \nabla y^*(x) \Vert  \\
    &\le  \left(1+ \frac{L_g}{\mu} \right) \frac{C_f \rho_g}{\mu} +L_g D_2,
\end{align*}
where we use Lemma \ref{lem:nabla-y-lambda-bound}
in the second inequality; Lemma \ref{lem:yy-lambda} and \ref{lem:yxyx-lambda} in the third one.
\end{proof}


\section{Lemmas for Finding Second-Order Stationarity} \label{apx:lem-2nd}


\begin{lem} \label{lem:nabla2-yx}
Under Assumption \ref{asm:sc} and \ref{asm:third}, 
for $\lambda \ge  2L_f / \mu$,
we have $ \Vert \nabla^2 y^*(x) -\nabla^2 y_{\lambda}^*(x)  \Vert \le D_4 /\lambda$, where 
\begin{align*}
    D_4 := \frac{2\rho_g}{\lambda \mu^2}  \left( 1 + \frac{L_g}{\mu}\right)^2  \left(L_f + 
     \frac{C_f \rho_g}{\mu}
     \right) +   \frac{14 L_g \rho_g D_2}{\lambda \mu^2} + \frac{50 L_g^2 }{\lambda \mu^3} \left( \frac{C_f \nu_g}{\mu} + \rho_f \right) = \fO(\kappa^5).
\end{align*}
\end{lem}

\begin{proof}
    By taking the derivative with respect to $x$ on
    \begin{align*}
        \nabla_{xy}^2 g(x,y^*(x)) + \nabla y^*(x) \nabla_{yy}^2 g(x,y^*(x)) =0,
    \end{align*}
    we obtain
    \begin{align*}
        & \nabla_{xxy}^3 g(x,y^*(x)) + \nabla_{yxy}^3 g(x,y^*(x)) \times_{1}  \nabla y^*(x) + \nabla^2 y^*(x) \times_3 \nabla_{yy}^2 g(x,y^*(x)) \\
        &\quad + \nabla_{xyy}^3 g(x,y^*(x)) \times_2 \nabla y^*(x)  +\nabla_{yyy}^3 g(x,y^*(x)) \times_1 \nabla y^*(x) \times_2 \nabla y^*(x) = 0.
    \end{align*}
    Rearranging to get
    {\begin{align} \label{eq:nabla2-yx}
    \begin{split}
         \nabla^2 y^*(x) 
        &= -\left( \nabla_{xxy}^3 g(x,y^*(x)) + \nabla_{yxy}^3 g(x,y^*(x)) \times_1 \nabla y^*(x) 
        \right)  \times_3 [\nabla_{yy}^2 g(x,y^*(x))]^{-1} \\
        &\quad -
        \nabla_{xyy}^3 g(x,y^*(x)) \times_2 \nabla y^*(x) \times_3 [\nabla_{yy}^2 g(x,y^*(x))]^{-1} \\
        &\quad-  \nabla_{yyy}^3 g(x,y^*(x)) \times_1 \nabla y^*(x) \times_2 \nabla y^*(x)    \times_3 [\nabla_{yy}^2 g(x,y^*(x))]^{-1}.
    \end{split}
    \end{align}}
    Similarly,
    {
    \begin{align} \label{eq:nabla2-yx-lambda}
    \begin{split}
        \nabla^2 y_{\lambda}^*(x)  &= -\left( \nabla_{xxy}^3 \fL_{\lambda} (x, y_{\lambda}^*(x)) + \nabla_{yxy}^3 \fL_\lambda(x, y_{\lambda}^*(x)) \times_1 \nabla y_{\lambda}^*(x)  \right) \times_3
        [ \nabla_{yy}^2 \fL_{\lambda}(x,y_{\lambda}^*(x))  ]^{-1} \\
        &\quad - 
        \nabla_{xyy}^3 \fL_{\lambda}(x,y_{\lambda}^*(x))
        \times_2 \times_3 [ \nabla_{yy}^2 \fL_{\lambda}(x,y_{\lambda}^*(x))]^{-1} \\
        &\quad - 
        \nabla y_{\lambda}^*(x)   +  
         \nabla_{yyy}^3 \fL_{\lambda} (x,y_{\lambda}^*(x))
        \times_1 \nabla y_{\lambda}^*(x)
        \times_2 \nabla y_{\lambda}^*(x) 
         \times_3 [ \nabla_{yy}^2 \fL_{\lambda}(x,y_{\lambda}^*(x))]^{-1}.
    \end{split}
    \end{align}}
Note that 
\begin{align*}
    \left \Vert  \nabla_{xxy}^3  g(x,y^*(x)) - \frac{\nabla_{xxy}^3 \fL_{\lambda} (x,y_{\lambda}^*(x))  }{\lambda}   \right \Vert \le \nu_g \Vert y_{\lambda}^*(x) - y^*(x) \Vert + \frac{\rho_f}{\lambda} \le \frac{1}{\lambda} \left(\frac{C_f \nu_g}{\mu} +\rho_f \right),
\end{align*}
and
\begin{align*}
    &\quad \left \Vert 
    \nabla_{yxy}^3 g(x,y^*(x)) \times_1 \nabla y^*(x)  - \frac{\nabla_{yxy}^3 \fL_{\lambda} (x,y_{\lambda}^*(x)) \times_1 \nabla y_{\lambda}^*(x) }{\lambda}
    \right \Vert \\
    &\le \Vert \nabla y^*(x) - \nabla y_{\lambda}^*(x) \Vert \Vert \nabla_{yxy}^3 g(x,y^*(x)) \Vert + \Vert \nabla y_{\lambda}^*(x) \Vert \left\Vert \nabla_{yxy}^3 g(x,y^*(x)) - \frac{\nabla_{yxy}^3 \fL_{\lambda}(x, y_{\lambda}^*(x))}{\lambda} \right \Vert \\
    &\le \frac{\rho_g D_2}{\lambda} + \frac{4L_g}{\lambda \mu} \left( \frac{C_f \nu_g}{\mu} + \rho_f \right),
\end{align*}
and
{ 
\begin{align*}
    &\quad  \left \Vert
   \nabla_{yyy}^3 g(x,y^*(x)) \times_1  \nabla y^*(x)  \times_2 \nabla y^*(x)  - \frac{\nabla_{yyy}^3 \fL_{\lambda} (x,y_{\lambda}^*(x)) \times_1 \nabla y_{\lambda}^*(x) \times_2  \nabla y_{\lambda}^*(x)} { \lambda}
    \right \Vert \\
    &\le \Vert \nabla y^*(x) \Vert \Vert \nabla_{yyy}^3 g(x,y^*(x))\Vert \Vert \nabla y^*(x) - \nabla y_{\lambda}^*(x) \Vert  \\
    &\quad + \Vert \nabla y_{\lambda}^*(x) \Vert \Vert \nabla_{yyy}^3 g(x,y^*(x))\Vert \Vert \nabla y^*(x) - \nabla y_{\lambda}^*(x) \Vert  \\
    &\quad +  \Vert \nabla y_{\lambda}^*(x) \Vert^2  \left \Vert  \nabla_{xxy}^3  g(x,y^*(x)) - \frac{\nabla_{xxy}^3 \fL_{\lambda} (x,y_{\lambda}^*(x))  }{\lambda}   \right \Vert \\
    &\le  \frac{5L_g \rho_g D_2}{\lambda \mu} + \frac{16L_g^2 }{\lambda \mu^2} \left( \frac{C_f \nu_g}{\mu} + \rho_f \right),
\end{align*}}
we can obtain that
{
\begin{align*}
    &\quad \Vert \nabla^2 y^*(x) - \nabla^2 y_{\lambda}^*(x) \Vert  \\
    &\le \rho_g \left( 1 + \frac{L_g}{\mu}\right)^2 \left \Vert [\nabla_{yy}^2 g(x,y^*(x))]^{-1} - \left [  \frac{\nabla_{yy}^2 \fL_{\lambda}(x,y_{\lambda}^*(x))}{\lambda} \right]^{-1}  \right \Vert \\
    &\quad + \left( \frac{7L_g \rho_g D_2}{\lambda \mu} + \frac{25L_g^2 }{\lambda \mu^2} \left( \frac{C_f \nu_g}{\mu} + \rho_f \right) \right) \left \Vert \left [  \frac{\nabla_{yy}^2 \fL_{\lambda}(x,y_{\lambda}^*(x))}{\lambda} \right]^{-1} \right \Vert \\
    &\le \frac{2\rho_g}{\lambda \mu^2}  \left( 1 + \frac{L_g}{\mu}\right)^2  \left(L_f + 
     \frac{C_f \rho_g}{\mu}
     \right) +   \frac{14 L_g \rho_g D_2}{\lambda \mu^2} + \frac{50 L_g^2 }{\lambda \mu^3} \left( \frac{C_f \nu_g}{\mu} + \rho_f \right) ,
\end{align*}}
where we use \Eqref{eq:diff-inv} in the second inequality.
\end{proof}

\begin{lem} \label{lem:nabla3-bound}
Under Assumption \ref{asm:sc} and \ref{asm:third}, for $ \lambda \ge 2 L_f /\mu$, 
$\nabla^2 \fL_{\lambda}^*(x) $ is $D_5$-Lipschitz,
 where
\begin{align*}
    D_5 &:=  \left( 1 + \frac{4L_g}{\mu} \right)^2 \left( 3\rho_f  + \frac{2 L_f \rho_g}{\mu} \right) + \left(1 + \frac{L_g}{\mu} \right)^2 \frac{C_f \nu_g}{\mu} \\
    &\quad + \left(2+ \frac{5L_g}{\mu} \right) D_2 \rho_g +  \left( 1+ \frac{L_g}{\mu} \right)^2 \frac{C_f \rho_g^2}{\mu^2} + L_g D_4   = \fO(\ell \kappa^5).
\end{align*}
\end{lem}

\begin{proof}
Similar to the proof of Lemma \ref{lem:nabla2-bound},  we split $\nabla^2 \fL_{\lambda}^*(x)$ into two terms:
\begin{align*}
    \nabla^2 \fL_{\lambda}^*(x) = \nabla A(x) + \nabla B(x),
\end{align*}
where both the mappings $A(x)$ and $B(x)$ both follow the definitions in \Eqref{eq:nabla-AB}. 
Taking total derivative on $A(x)$, we obtain
\begin{align*}
    \nabla A(x) = \nabla_{xx}^2 f(x,y_{\lambda}^*(x)) + \nabla y_{\lambda}^*(x) \nabla_{yx}^2 f(x,y_{\lambda}^*(x)).
\end{align*}
And recall \Eqref{eq:nabla-Bx} that 
\begin{align*}
    \nabla B(x) &= \lambda (\nabla_{xx}^2 g(x,y_{\lambda}^*(x))-\nabla_{xx}^2 g(x,y^*(x))) \\
    &\quad + \lambda \left(\nabla y_{\lambda}^*(x)  \nabla_{yx}^2 g(x,y_{\lambda}^*(x)) - \nabla y^*(x) \nabla_{yx}^2 g(x,y^*(x)) \right).
\end{align*}
Then we bound the Lipschitz coefficient of $\nabla A(x)$ and $\nabla B(x)$, respectively.

From \Eqref{eq:nabla2-yx} we can calculate that 
\begin{align} \label{bound:nabla2-yx}
    \Vert \nabla^2 y^*(x) \Vert \le \left( 1 + \frac{L_g}{\mu} \right)^2 \frac{\rho_g}{\mu}  .
\end{align}
Similarly, from \Eqref{eq:nabla2-yx-lambda} we can calculate that
\begin{align} \label{bound:nabla2-y-lambda}
    \Vert \nabla^2 y_{\lambda}^*(x)  \Vert &\le \left( 1 + \frac{4L_g}{\mu} \right)^2 \left( \frac{\rho_f}{\lambda}+ \rho_g \right) \frac{2}{\mu} \le \left( 1 + \frac{4L_g}{\mu} \right)^2 \left( \frac{2\rho_f}{L_f} + \frac{2 \rho_g}{\mu} \right). 
\end{align}
From Lemma \ref{lem:nabla-y-lambda-bound} we know that $y_{\lambda}^*(x)$ is $(4L_g/\mu)$-Lipschitz. This further implies that both $ \nabla_{xx}^2 f(x,y_{\lambda}^*(x))$ and $\nabla_{yx}^2 f(x,y_{\lambda}^*(x)) $ are $(1+ 4L_g/ \mu)\rho_f$-Lipschitz. Then, for any $x_1,x_2 \in \BR^{d_x}$, we have
\begin{align*}
    &\quad  \Vert \nabla A(x_1) - \nabla A(x_2) \Vert \\
   &\le \Vert  \nabla_{xx}^2 f(x_1,y_{\lambda}^*(x_1)) - \nabla_{xx}^2 f(x_2, y_{\lambda}^*(x_2)) \Vert  \\
   &\quad + \Vert \nabla y_{\lambda}^*(x_1) \nabla_{yx}^2 f(x_1,y_{\lambda}^*(x_1)) - \nabla y_{\lambda}^*(x_2) \nabla_{yx}^2 f(x_1,y_{\lambda}^*(x_1)) \Vert \\
   &\quad + \Vert \nabla y_{\lambda}^*(x_2) \nabla_{yx}^2 f(x_1,y_{\lambda}^*(x_1)) - \nabla y_{\lambda}^*(x_2) \nabla_{yx}^2 f(x_2,y_{\lambda}^*(x_2)) \Vert \\
   &\le \Vert \nabla_{xx}^2 f(x_1,y_{\lambda}^*(x_1)) - \nabla_{xx}^2 f(x_2, y_{\lambda}^*(x_2)) \Vert  \\
   &\quad + \Vert \nabla y_{\lambda}^*(x_1) - \nabla y_{\lambda}^*(x_2) \Vert \Vert \nabla_{yx}^2 f(x_1,y_{\lambda}^*(x_1)) \Vert \\
   &\quad + \Vert \nabla y_{\lambda}^*(x_2) \Vert \Vert \nabla_{yx}^2 f(x_1,y_{\lambda}^*(x_1)) - \nabla_{yx}^2 f(x_2,y_{\lambda}^*(x_2)) \Vert \\
   &\le \underbrace{\left(\left( 1+ \frac{4L_g}{\mu} \right)^2 \rho_f + \left( 1 + \frac{4L_g}{\mu} \right)^2 \left( \frac{2\rho_f}{L_f} + \frac{2 \rho_g}{\mu} \right) L_f \right)}_{C_1}\Vert x_1-x_2 \Vert,
\end{align*}
where $C_1$ gives the upper bound of the Lipschitz coefficient of the mapping $\nabla A(x)$.

To bound the Lipschitz coefficient of $\nabla B(x)$,
we first derive the explicit form of the mapping $\nabla^2 B(x): \BR^{d_x} \rightarrow \BR^{d_x} \times \BR^{d_x} \times \BR^{d_x} $ by:
\begin{align*}
    \nabla^2 B(x) 
    &= 
    \nabla^2 y_{\lambda}^*(x) \times_3 [\nabla_{yx}^2 f(x,y_{\lambda}^*(x))]^\top \\
    &\quad + \lambda ( \nabla_{xxx}^3 g(x,y_{\lambda}^*(x)) - \nabla_{xxx}^3 g(x,y^*(x)) ) \\
    &\quad + \lambda ( 
    \nabla_{yxx}^3 g(x,y_{\lambda}^*(x)) \times_1 \nabla y_{\lambda}^*(x) - \nabla_{yxx}^3 g(x,y^*(x)) \times_1 \nabla y^*(x) 
    ) \\
    &\quad + \lambda \left(  \nabla_{xyx}^3 g(x,y_{\lambda}^*(x)) \times_2 \nabla y_{\lambda}^*(x) - \nabla_{xyx}^3 g(x,y^*(x))   \times_2 \nabla y^*(x) \right)  \\
    &\quad + \lambda \left( 
     \nabla_{yyx}^3 g(x,y_{\lambda}^*(x)) \times_1 \nabla y_{\lambda}^*(x) \times_2  \nabla y_{\lambda}^*(x)   - 
    \nabla_{yyx}^3 g(x,y^*(x)) \times_1 \nabla y^*(x)  \times_2 \nabla y^*(x)
    \right) \\
    &\quad + \lambda \left(  \nabla^2 y_{\lambda}^*(x) \times_3 [\nabla_{yx}^2 g(x,y_{\lambda}^*(x))]^\top -  \nabla^2 y^*(x) \times_3 [\nabla_{yx}^2 g(x,y^*(x))]^\top\right)
    .
\end{align*}
Then we can bound the Lipschitz coefficient of $\nabla B(x)$
via the operator norm of $\nabla^2 B(x)$:
\begin{align*}
   C_2&:= \Vert \nabla^2 B(x) \Vert \\
    &\le 
    \Vert \nabla_{xxx}^3 g(x,y^*(x)) - 
    \nabla_{xxx}^3 g(x,y_{\lambda}^*(x))\Vert  \\
    &\quad + \lambda \Vert \nabla y^*(x) \Vert \Vert \nabla_{yxx}^3 g(x,y^*(x)) - \nabla_{yxx}^3 g(x,y_{\lambda}^*(x)) \Vert \\
    &\quad+ \lambda \Vert \nabla y_{\lambda}^*(x) - \nabla y^*(x) \Vert \Vert \nabla_{yxx}^3 g(x,y_{\lambda}^*(x)) \Vert \\
    &\quad + \lambda \Vert \nabla y^*(x) \Vert \Vert \nabla_{xyx}^3 g(x,y^*(x)) - \nabla_{xyx}^3 g(x,y_{\lambda}^*(x)) \Vert  \\
    &\quad + \lambda \Vert \nabla y^*(x) - \nabla y_{\lambda}^*(x) \Vert \Vert \nabla_{xyx}^3 g(x,y_{\lambda}^*(x)) \Vert \\
    &\quad + \lambda \Vert \nabla y^*(x) \Vert \Vert \nabla_{yyx}^3 g(x,y^*(x)) \Vert  \Vert \nabla y_{\lambda}^*(x) - \nabla y^*(x) \Vert  \\
    &\quad + \lambda \Vert \nabla y_{\lambda}^*(x) \Vert \Vert \nabla_{yyx}^3 g(x,y^*(x)) \Vert  \Vert \nabla y_{\lambda}^*(x) - \nabla y^*(x) \Vert \\
    &\quad + \lambda \Vert \nabla y^*(x) \Vert^2 \Vert \nabla_{yyx}^3 g(x,y^*(x)) - \nabla_{yyx}^3 g(x,y_{\lambda}^*(x)) \Vert \\
    &\quad + \lambda \Vert \nabla^2 y^*(x) \Vert \Vert \nabla_{yx}^2 g(x,y^*(x)) - \nabla_{yx}^2 g(x,y_{\lambda}^*(x)) \Vert \\
    &\quad + \lambda \Vert \nabla^2 y^*(x) - \nabla^2 y_{\lambda}^*(x) \Vert  \Vert \nabla_{yx}^2 g(x,y_{\lambda}^*(x)) \Vert.
\end{align*}
Then we can bound $C_2$ using the upper bound for
$\Vert \nabla y^*(x) \Vert$ given by \Eqref{eq:y-x}, that for 
$\Vert \nabla y_{\lambda}^*(x) \Vert$ by Lemma \ref{lem:nabla-y-lambda-bound}, that for $ \Vert \nabla y^*(x) - \nabla y_{\lambda}^*(x) \Vert$ by Lemma \ref{lem:yxyx-lambda}, that for $ \Vert \nabla^2 y^*(x) \Vert$ by \Eqref{bound:nabla2-yx} and that for $ \Vert \nabla^2 y^*(x) - \nabla^2 y_{\lambda}^*(x) \Vert$ by \Eqref{lem:nabla2-yx}. 
Finally, we can obtain 
\begin{align*}
    C_1+C_2 & \le \left( 1 + \frac{4L_g}{\mu} \right)^2 \left( 3\rho_f  + \frac{2 L_f \rho_g}{\mu} \right) + \left(1 + \frac{L_g}{\mu} \right)^2 \frac{C_f \nu_g}{\mu} \\
    &\quad + \left(2+ \frac{5L_g}{\mu} \right) D_2 \rho_g +  \left( 1+ \frac{L_g}{\mu} \right)^2 \frac{C_f \rho_g^2}{\mu^2} + L_g D_4.
\end{align*}
\end{proof}
 
\begin{lem} \label{lem:nabla2-phi}
Under Assumption \ref{asm:sc} and \ref{asm:third}, for $\lambda \ge 2 L_f /\mu$, we have $ \Vert \nabla^2 \fL_{\lambda}^*(x) - \nabla^2 \varphi(x) \Vert \le  D_6 /\lambda$,
where
\begin{align*}
    D_6 := 2L_g D_2^2 + \left( 1 + \frac{L_g}{\mu}\right)^2\left( \frac{C_f \rho_f}{\mu} + \frac{C_f L_f \rho_g}{\mu^2} + \frac{C_f^2 \nu_g}{2 \mu^2} + \frac{C_f^2 \rho_g^2}{2 \mu^3}  \right)  = \fO(\ell \kappa^6).
\end{align*}
\end{lem}

\begin{proof}
    Taking total derivative on $ \nabla \varphi(x) = \nabla_x f(x,y^*(x)) + \nabla y^*(x) \nabla_y f(x,y^*(x))$ yields
    \begin{align} \label{eq:phi-hess}
\begin{split}
    \nabla^2 \varphi(x) &= \nabla_{xx}^2 f(x,y^*(x)) + \nabla y^*(x) \nabla_{yx}^2 f(x,y^*(x))  + \nabla^2 y^*(x) \bar \times_3 \nabla_y f(x,y^*(x)) \\
    &\quad + \nabla_{xy}^2f(x,y^*(x)) [\nabla y^*(x)]^\top + \nabla y^*(x) \nabla_{yy}^2 f(x,y^*(x)) [\nabla y^*(x)]^\top.
\end{split}
\end{align}
Plug into the close form of $\nabla^2 y^*(x)$ given by \Eqref{eq:nabla2-yx}, we arrive at
    \begin{align*}
    &\quad \nabla^2 \varphi(x) \\
    &= \underbrace{\nabla_{xx }^2  f(x,y^*(x)) - \nabla_{xxy}^3 g(x,y^*(x)) \times_3 [  \nabla_{yy}^2 g(x,y^*(x))]^{-1} \bar \times_3 \nabla_y f(x,y^*(x)) }_{\rm (I)} \\
    &\quad +  
    \underbrace{  \left(\nabla_{yx}^2 f(x,y^*(x)) - \nabla_{yxy}^3 g(x,y^*(x)) \times_3 [\nabla_{yy}^2 g(x,y^*(x)]^{-1}  \bar \times_3 \nabla_y f(x,y^*(x)) \right) \times_1 \nabla y^*(x)  }_{\rm (II)}  \\
    &\quad +  
    \underbrace{  \left(\nabla_{xy}^2 f(x,y^*(x)) - \nabla_{xyy}^3 g(x,y^*(x)) \times_3 [\nabla_{yy}^2 g(x,y^*(x)]^{-1}  \bar \times_3 \nabla_y f(x,y^*(x)) \right) \times_2 \nabla y^*(x)  }_{\rm (III)}  \\
    &\quad + \underbrace{  \left( \nabla_{yy}^2 f(x,y^*(x)) -  \nabla_{yyy}^3 g(x,y^*(x)) \times_3 [\nabla_{yy}^2 g(x,y^*(x)]^{-1} \bar \times_3 \nabla_y f(x,y^*(x))   \right) \times_1 \nabla y^*(x) \times_2 \nabla y^*(x) }_{\rm (IV)}. 
\end{align*}
Recall that
\begin{align*}
    \nabla^2 \fL_{\lambda}^*(x) &= \nabla_{xx}^2  f(x,y_{\lambda}^*(x)) +  \lambda ( \nabla_{xx}^2 g(x,y_{\lambda}^*(x)) - \nabla_{xx}^2 g(x,y^*(x))) \\
    &\quad + \nabla y_{\lambda}^*(x) \nabla_{yx}^2 \fL_{\lambda}(x,y_{\lambda}^*(x))
    - \lambda \nabla y^*(x) \nabla_{yx}^2 g(x,y^*(x)) .
\end{align*}
Let
\begin{align*}
    \tilde \nabla^2 \fL_{\lambda}^*(x) &= \underbrace{\nabla_{xx}^2  f(x,y_{\lambda}^*(x)) +  \lambda ( \nabla_{xx}^2 g(x,y_{\lambda}^*(x)) - \nabla_{xx}^2 g(x,y^*(x)))}_{\rm (I')}\\
    &+  \underbrace{\nabla y^*(x) \left(\nabla_{yx}^2 \fL_{\lambda}(x,y_{\lambda}^*(x))- \lambda \nabla_{yx}^2 g(x,y^*(x))\right) }_{\rm (II')}\\ 
    & + \underbrace{ \left(\nabla_{xy}^2 \fL_{\lambda}(x,y_{\lambda}^*(x))- \lambda \nabla_{xy}^2 g(x,y^*(x))\right)[\nabla y^*(x)]^\top }_{\rm (III')} \\
    &+ \underbrace{\nabla y^*(x) \left(\nabla_{yy}^2 \fL_{\lambda} (x,y_{\lambda}^*(x))   - \lambda \nabla_{yy}^2 g(x,y^*(x)) \right) [\nabla y^*(x)]^\top }_{\rm (IV')}. 
\end{align*}
Recall \Eqref{eq:y-lambda-x} that
\begin{align*}
    \nabla y_{\lambda}^*(x) = - \nabla_{xy}^2 \fL_{\lambda}(x,y_{\lambda}^*(x)) [\nabla_{yy}^2 \fL_{\lambda}(x,y_{\lambda}^*(x)) ]^{-1}.
\end{align*}
Then using
equivalent transformation and the identity
\begin{align*}
    U A U^\top - V A V^\top = (U-V) A V^\top + VA (U-V)^\top + (U-V)A (U-V)^\top
\end{align*}
for a matrices $U,V$ and $A$,
we can show that
\begin{align*}
    &\quad \tilde \nabla^2 \fL_{\lambda}^*(x) - \nabla^2 \fL_{\lambda}^*(x) \\
    &= ( \nabla y^*(x) - \nabla y_{\lambda}^*(x) ) \nabla_{yx}^2 \fL_{\lambda} (x,y_{\lambda}^*(x)) + \\
    &\quad + \nabla_{xy}^2 \fL_{\lambda}(x,y_{\lambda}^*(x)) [ \nabla y^*(x)]^\top + \nabla y^*(x) \nabla_{yy}^2 \fL_{\lambda} (x,y_{\lambda}^*(x)) [\nabla y^*(x)]^\top \\
    &= ( \nabla y^*(x) - \nabla y_{\lambda}^*(x) ) \nabla_{yx}^2 \fL_{\lambda} (x,y_{\lambda}^*(x)) + \nabla_{xy}^2 \fL_{\lambda}(x,y_{\lambda}^*(x)) [ \nabla y^*(x) - \nabla y_{\lambda}^*(x)]^\top \\
    &\quad + \nabla y^*(x) \nabla_{yy}^2 \fL_{\lambda} (x,y_{\lambda}^*(x)) [\nabla y^*(x)]^\top - \nabla y_{\lambda}^*(x) \nabla_{yy}^2 \fL_{\lambda} (x,y_{\lambda}^*(x)) [\nabla y_{\lambda}^*(x)]^\top \\
    &= ( \nabla y^*(x) - \nabla y_{\lambda}^*(x) ) \nabla_{yx}^2 \fL_{\lambda} (x,y_{\lambda}^*(x)) + \nabla_{xy}^2 \fL_{\lambda}(x,y_{\lambda}^*(x)) [ \nabla y^*(x) - \nabla y_{\lambda}^*(x)]^\top \\ 
    &\quad + (\nabla y^*(x) - \nabla y_{\lambda}^*(x)) \nabla_{yy}^2 \fL_{\lambda} (x,y_{\lambda}^*(x)) [ \nabla y_{\lambda}^*(x)]^\top \\
    &\quad + \nabla y_{\lambda}^*(x) \nabla_{yy}^2 \fL_{\lambda} (x,y_{\lambda}^*(x)) (\nabla y^*(x) - \nabla y_{\lambda}^*(x))^\top \\
    &\quad + (\nabla y^*(x) - \nabla y_{\lambda}^*(x)) \nabla_{yy}^2 \fL_{\lambda} (x,y_{\lambda}^*(x)) [ \nabla y^*(x) - \nabla y_{\lambda}^*(x)]^\top \\
    &= (\nabla y^*(x) - \nabla y_{\lambda}^*(x)) \nabla_{yy}^2 \fL_{\lambda} (x,y_{\lambda}^*(x)) [ \nabla y^*(x) - \nabla y_{\lambda}^*(x)]^\top.
\end{align*}
Noting that $ \nabla_{yy}^2 \fL_{\lambda}(\,\cdot\, , \,\cdot\,) \preceq 2 \lambda L_g$, we have
\begin{align*}
    \Vert \tilde \nabla^2 \fL_{\lambda}^*(x) - \nabla^2 \fL_{\lambda}^*(x) \Vert \le 2 \lambda L_g \Vert \nabla y^*(x) - \nabla y_{\lambda}^*(x) \Vert^2 \le \frac{2L_g D_2^2}{\lambda }. 
\end{align*}
Next, we bound difference between 
$\tilde \nabla^2 \fL_{\lambda}^*(x) $ and $\nabla^2 \fL_{\lambda}^*(x)$. For the difference between ${\rm (I)}$ and ${\rm (I')}$, 
{\begin{align*}
    &\quad \rm (I) - {\rm (I')} \\
    &= \nabla_{xx}^2 f(x,y^*(x)) - \nabla_{xx}^2 f(x,y_{\lambda}^*(x))  \\
    & \quad + \nabla_{xxy}^3 g(x,y^*(x)) \times_3 [\nabla_{yy}^2 g(x,y^*(x))]^{-1} \bar \times_3 ( \nabla_y f(x,y_{\lambda}^*(x)) - \nabla_y f(x,y^*(x))) \\
    &\quad + \lambda \left( \nabla_{xx}^2 g(x,y^*(x)) -  \nabla_{xx}^2 g(x,y_{\lambda}^*(x))  + \nabla_{xxy}^3  g(x,y^*(x)) \bar \times_3
    (y_{\lambda}^*(x) - y^*(x)) \right) \\
    &\quad + \lambda \nabla_{xxy}^3 g(x,y^*(x)) \times_3 [\nabla_{yy}^2 g(x,y^*(x))]^{-1} 
 \\
 &~~~~~~~~~~~~~~~~~~~~~~~~~~~\bar \times_3 \left(\nabla_y g(x,y_{\lambda}^*(x)) 
    - \nabla_y g(x,y^*(x)) - \nabla_{yy}^2 g(x,y^*(x)) (y_{\lambda}^*(x) -y^*(x))  \right).
\end{align*}}
Therefore, 
\begin{align*}
    \Vert {\rm (I) } - {\rm (I')} \Vert &\le \rho_f \Vert y^*(x) - y_{\lambda}^*(x) \Vert  + \frac{L_f \rho_g}{\mu} \Vert y^*(x) - y_{\lambda}^*(x) \Vert \\
    &\quad + \frac{\lambda \nu_g}{2} \Vert y^*(x) - y_{\lambda}^*(x) \Vert^2 + \frac{\lambda \rho_g^2}{2\mu} \Vert y^*(x) - y_{\lambda}^*(x) \Vert^2  \\
    &\le \frac{C_f \rho_f}{\lambda \mu} + \frac{C_f L_f \rho_g}{\lambda \mu^2} +  \frac{C_f^2 \nu_g}{2\lambda \mu^2}+ \frac{C_f^2 \rho_g^2}{2 \lambda \mu^3} .
\end{align*}
Using $\Vert \nabla y^*(x) \Vert \le  L_g / \mu$, we can similarly bound the difference between ${\rm (II)}$ and ${\rm (II')}$ by
{ \begin{align*}
    &\quad \lVert{\rm (II)} - {\rm (II')}\rVert \\
    &\le 
      \Vert \nabla y^*(x) \Vert \Big \Vert  \nabla_{yx}^2 f(x,y^*(x)) - \nabla_{yxy}^3 g(x,y^*(x)) \times_3 [\nabla_{yy}^2 g(x,y^*(x)]^{-1} \bar \times_3 \nabla_y f(x,y^*(x))\\
     &\quad - \nabla_{yx}^2 f(x,y_{\lambda}^*(x))-\lambda  \left(\nabla_{yx}^2 g(x,y_{\lambda}^*(x))-\nabla_{yx}^2 g(x,y^*(x)) \right) \Big \Vert \\
     & \le \frac{L_g }{\mu} \left( \frac{C_f \rho_f}{\lambda \mu} + \frac{C_f L_f \rho_g}{\lambda \mu^2}  + \frac{C_f^2 \nu_g}{2\lambda \mu^2} +\frac{C_f^2 \rho_g^2}{2 \lambda \mu^3} \right),
\end{align*}}
And the bound for ${\rm(III)}$ and $ {\rm (III')}$ is the same. Finally, we know that
\begin{align*}
    &\quad \left \Vert {\rm (IV)} - {\rm (IV')} \right \Vert \\
    &\le \Vert \nabla y^*(x) \Vert^2 \Big \Vert \nabla_{yy}^2 f(x,y^*(x)) - \nabla_{yyy}^3 g(x,y^*(x)) \times_3 [ \nabla_{yy}^2 g(x,y^*(x))]^{-1} \bar \times_3 \nabla_y f(x,y^*(x))  \\
    &\quad - \nabla_{yy}^2 f(x,y_{\lambda}^*(x)) - \lambda \left( \nabla_{yy}^2 g(x,y_{\lambda}^*(x)) - \nabla_{yy}^2 g(x,y^*(x)) \right) \Big \Vert \\
    &\le \frac{L_g^2}{\mu^2} \left( \frac{C_f \rho_f}{\lambda \mu} + \frac{C_f L_f \rho_g}{\lambda \mu^2}  + \frac{C_f^2 \nu_g}{2\lambda \mu^2} +  \frac{C_f^2 \rho_g^2}{2 \lambda \mu^3}\right).
\end{align*}
Combing the above bounds completes our proof.
\end{proof}

\section{Proofs of Finding First-Order Stationarity} \label{apx:1st}

We recall the standard result for gradient descent for strongly convex functions, which is used for proving the linear convergence of $y_t^k \rightarrow y_{\lambda}^*(x_t) $ and $z_{t}^k \rightarrow y^*(x_t)$.
\begin{thm}[Theorem 3.10 by \citet{bubeck2015convex}] \label{thm:GD-SC}
Suppose $h(x): \BR^d \rightarrow \BR$ is  $\beta$-gradient Lipschitz and $\alpha$-strongly convex. Consider the following update of gradient descent:
\begin{align*}
    x_{t+1} = x_t - \frac{1}{\beta} \nabla h(x_t).
\end{align*}
Let $x^* = \arg \min_{x \in \BR^d} h(x)$. Then it holds that
\begin{align*}
    \Vert x_{t+1} - x^* \Vert^2 \le \left(1 - \frac{\alpha}{\beta} \right) \Vert x_t - x^* \Vert^2.
\end{align*}
\end{thm}

%\propGD*

% \begin{proof}
% It is a straightforward result of the following descent lemma for gradient descent.
% \begin{align*}
%     \fL_{\lambda}^*(x_{t+1}) &\le \fL_{\lambda}^*(x_t) + \langle \nabla \fL_{\lambda}^*(x_t) , x_{t+1} - x_t
%     \rangle + \frac{L}{2} \Vert x_{t+1 } - x_t \Vert^2 \\
%     &= \fL_{\lambda}^*(x_t) -
%     \frac{\eta}{2} \Vert \nabla \fL_{\lambda}^*(x_t) \Vert^2 -
%     \left( \frac{\eta}{2} - \frac{\eta^2 L}{2} \right) \Vert \hat \nabla \fL_{\lambda}^*(x_t) \Vert^2 + \frac{\eta}{2} \Vert \hat  \nabla \fL_{\lambda}^*(x_t) - \nabla \fL_{\lambda}^*(x_t) \Vert^2 \\
%     &\le \fL_{\lambda}^*(x_t) 
%     - \frac{\eta}{2} \Vert \nabla \fL_{\lambda}^*(x_t) \Vert^2 
%     + \frac{\eta}{2} \Vert  \hat \nabla \fL_{\lambda}^*(x_t) - \nabla \fL_{\lambda}^*(x_t) \Vert^2.
% \end{align*}
% A simple telescoping completes the proof.
% \end{proof}

% \propPGD*

% \begin{proof}
%     It is the same as the proof of Theorem 2.7 in \cite{huang2022efficiently}.
% \end{proof}



% \properror*

% \begin{proof}
% According to
% \begin{align*}
%     \Vert \hat \nabla \fL_{\lambda}^*(x_t) - \nabla \fL_{\lambda}^*(x_t) \Vert \le (L_f+ \lambda L_g) \Vert y_{t+1}- y_{\lambda}^*(x_t) \Vert + \lambda L_g \Vert z_{t+1} - y^*(x_t) \Vert,
% \end{align*}
% it suffices to let 
% \begin{align*}
% \Vert y_{t+1} - y_{\lambda}^*(x_t) \Vert & \le \frac{\zeta}{4\lambda L_g}   ~{\rm and}~
% \Vert z_{t+1} - y^*(x_t) \Vert \le \frac{\zeta}{2 \lambda L_g}
% \end{align*}
% holds for all $t$. Below, we only show how to control $\Vert y_{t+1} - y_{\lambda}^*(x_t) \Vert$ since the analysis for $ \Vert z_{t+1} - y^*(x_t) \Vert$ is very similar. 
% Since $\fL_{\lambda}(x,y)$ is   $(\lambda \mu /2)$-strongly convex and $(2 \lambda L_g)$-smooth in $y$.
% From Theorem 2.1.14 in \cite{nesterov2018lectures}, we know that
% \begin{align*}
%     \Vert y_{t+1} - y_{\lambda}^*(x_t) \Vert \le \left( 1 - \frac{\mu}{4L_g} \right)^{K} \Vert y_t - y^*(x_t) \Vert.
% \end{align*}
% This indicates that we can  set
% \begin{align*}
%     K_t \ge \frac{4 L_g}{\mu} \log \left( 
%     \frac{4 \lambda L_g \Vert y_t - y_{\lambda}^*(x_t) \Vert}{\zeta}
%     \right), \quad \forall 0 \le t \le T-1.
% \end{align*}
% Next, we need to bound $\Vert y_t - y_{\lambda}^*(x_t) \Vert$ for all $t$. For the case of $t=0$, we have
% \begin{align*}
%     \Vert y_0 - y_{\lambda}^*(x_0) \Vert \le \Vert y_0 - y^*(x_0) \Vert + \Vert y^*(x_0) - y_{\lambda}^*(x_0) \Vert \le \Vert y_0 - y^*(x_0) \Vert + \frac{C_f}{\lambda \mu}.
% \end{align*}
% For the case of $t \ge 1$, we can prove it by induction. Suppose that we have $\Vert y_{t} - y_{\lambda}^*(x_{t-1})\Vert \le \zeta / (4 \lambda L_g)$, then
% \begin{align*}
%     \Vert y_t - y_{\lambda}^*(x_t) \Vert  &\le  \Vert y_{t} - y_{\lambda}^*(x_{t-1}) \Vert + \Vert y_{\lambda}^*(x_{t}) - y_{\lambda}^*(x_{t-1}) \Vert \\
%     &\le \frac{\zeta}{4 \lambda L_g} + \frac{4 L_g}{\mu} \Vert x_t - x_{t-1} \Vert \\
%     &\le \frac{\zeta}{4 \lambda L_g} + \frac{4 L_g \eta}{\mu} \Vert \hat \nabla \fL_{\lambda}^*(x_{t-1})\Vert \\
%     &\le \frac{\zeta}{4 \lambda L_g} + \frac{4 L_g \eta \zeta}{\mu} + \frac{4 L_g \eta}{\mu} \Vert \nabla \fL_{\lambda}^*(x_{t-1})\Vert,
% \end{align*}
% where the second inequality uses the fact that $y_{\lambda}^*(\,\cdot\,)$ is $(4 L_g/\mu)$-Lipschitz.
%     Recall that
%     \begin{align*}
%         \nabla \fL_{\lambda}^*(x) = f(x,y_{\lambda}^*(x)) + \lambda ( 
%         g(x,y_{\lambda}^*(x) ) - g(x,y^*(x))
%         ).
%     \end{align*}
%     According to Lemma \ref{lem:yy-lambda}, we have
%     \begin{align*}
%         \Vert \nabla \fL_{\lambda}^*(x) \Vert \le C_f + \lambda L_g \Vert y_{\lambda}^*(x) - y^*(x) \Vert \le \left( 1 + \frac{L_g}{\mu}\right) C_f.
%     \end{align*}
% Thus,
% \begin{align*}
%     \Vert y_t - y_{\lambda}^*(x_t) \Vert \le \frac{\zeta}{4 \lambda L_g} + \frac{4 L_g \eta }{\mu} \left( \zeta +   \left( 1 + \frac{L_g}{\mu}\right) C_f\right).
% \end{align*}

% \end{proof}

Below, we prove that F${}^2$BA is indeed near-optimal.

\thmFBA* 
\begin{proof}
Let $L$ be the gradient Lipschitz coefficient of $\fL_{\lambda}^*(x)$.
According to Lemma \ref{lem:1st} and the setting of $\lambda$, we have
\begin{enumerate}[label=\alph*.]
    \item $\sup_{x \in \BR^{d_x}} \Vert \nabla \fL_{\lambda}^*(x) - \nabla \varphi(x) \Vert = \fO(\epsilon)$. 
    \item $\fL_{\lambda}^*(x_0) - \inf_{x \in \BR^{d_x}}  \fL_{\lambda}^*(x) = \fO(\Delta)$.
    \item $ L:=\sup_{x \in \BR^{d_x}}\Vert \nabla^2 \fL_{\lambda}^*(x) \Vert = \fO(\ell \kappa^3)$.
\end{enumerate}
Due to Lemma \ref{lem:yy-lambda}, we also have
\begin{align*}
    \Vert y_0 - y_{\lambda}^*(x_0) \Vert^2 + \Vert y_0 - y^*(x_0) \Vert^2 =\fO(R). 
\end{align*}
Now it suffices to show that the algorithm can find an $\epsilon$-first-order stationary of $\fL_{\lambda}^*(x)$. 

We begin from the following descent lemma for gradient descent. Let $\eta \le 1/(2L)$, then
\begin{align} \label{eq:des}
\begin{split}
    \fL_{\lambda}^*(x_{t+1}) &\le \fL_{\lambda}^*(x_t) + \langle \nabla \fL_{\lambda}^*(x_t) , x_{t+1} - x_t
    \rangle + \frac{L}{2} \Vert x_{t+1 } - x_t \Vert^2 \\
    &= \fL_{\lambda}^*(x_t) -
    \frac{\eta}{2} \Vert \nabla \fL_{\lambda}^*(x_t) \Vert^2 -
    \left( \frac{\eta}{2} - \frac{\eta^2 L}{2} \right) \Vert \hat \nabla \fL_{\lambda}^*(x_t) \Vert^2 + \frac{\eta}{2} \Vert \hat  \nabla \fL_{\lambda}^*(x_t) - \nabla \fL_{\lambda}^*(x_t) \Vert^2 \\
    &\le \fL_{\lambda}^*(x_t) 
    - \frac{\eta}{2} \Vert \nabla \fL_{\lambda}^*(x_t) \Vert^2 - 
    \frac{1}{4\eta} \Vert x_{t+1} - x_{t} \Vert^2
    + \frac{\eta}{2} \Vert  \hat \nabla \fL_{\lambda}^*(x_t) - \nabla \fL_{\lambda}^*(x_t) \Vert^2.
\end{split}
\end{align}
Note that
\begin{align*}
    \Vert \hat \nabla \fL_{\lambda}^*(x_t) - \nabla \fL_{\lambda}^*(x_t) \Vert \le 2 \lambda L_g \Vert y_{t}^K- y_{\lambda}^*(x_t) \Vert + \lambda L_g \Vert z_{t}^K - y^*(x_t) \Vert.
\end{align*}
By Theorem \ref{thm:GD-SC}, we have
\begin{align*}
    \Vert y_t^K - y_{\lambda}^*(x_t) \Vert^2 &\le \exp\left( - \frac{\mu K}{4 L_g} \right) \Vert y_t^0 - y_{\lambda}^*(x_t) \Vert^2 \\
     \Vert z_t^K - y^*(x_t) \Vert^2 &\le \exp\left( - \frac{\mu K}{L_g} \right) \Vert z_t^0 - y^*(x_t) \Vert^2
\end{align*}
Therefore, we have
\begin{align} \label{eq:err}
    \Vert \hat \nabla \fL_{\lambda}^*(x_t)- \nabla \fL_{\lambda}^*(x_t) \Vert^2 \le 4 \lambda^2 L_g^2 \exp\left( - \frac{\mu K}{4 L_g} \right) \left( \Vert y_t^0 - y_{\lambda}^*(x_t) \Vert^2 + \Vert z_t^0 - y^*(x_t) \Vert^2 \right)
\end{align}
By Young's inequality, 
\begin{align*}
    \Vert y_{t+1}^0 - y_{\lambda}^*(x_{t+1}) \Vert^2 & \le 2 \Vert y_t^K - y_{\lambda}^*(x_{t}) \Vert^2 + 2 \Vert y_{\lambda}^*(x_{t+1}) - y_{\lambda}^*(x_t) \Vert^2  \\
    &\le 2 \exp \left( - \frac{\mu K}{4 L_g} \right) \Vert y_t^0 - y_{\lambda}^*(x_t) \Vert^2 + \frac{32 L_g^2}{\mu^2} \Vert x_{t+1} - x_t \Vert^2,
\end{align*}
where the second inequality follows Lemma \ref{lem:nabla-y-lambda-bound} that $y_{\lambda}^*(x)$ is $(4 L_g/\mu)$-Lipschitz.  Similarly, we can derive the recursion about $\Vert z_t^0 - y^*(x_t) \Vert^2$. 

Put them together and let $K \ge 8 L_g/ \mu$, we have
\begin{align} \label{eq:delta_t}
    &\quad \Vert y_{t+1}^0 - y_{\lambda}^*(x_{t+1}) \Vert^2  + \Vert z_{t+1}^0 - y^*(x_{t+1}) \Vert^2 \\
    &\le 2 \exp \left( - \frac{\mu K}{4 L_g} \right) \left(\Vert y_t^0 - y_{\lambda}^*(x_t) \Vert^2 + \Vert z_t^0 - y^*(x_t) \Vert^2 \right) +\frac{34 L_g^2}{\mu^2} \Vert x_{t+1} - x_t \Vert^2 \\
    &\le \frac{1}{2} \left(\Vert y_t^0 - y_{\lambda}^*(x_t) \Vert^2 + \Vert z_t^0 - y^*(x_t) \Vert^2 \right) +\frac{34 L_g^2}{\mu^2} \Vert x_{t+1} - x_t \Vert^2.
\end{align}
Telescoping over $t$ yields
\begin{align*}
    &\quad \Vert y_t^0 - y_{\lambda}^*(x_t) \Vert^2 + \Vert z_t^0 - y^*(x_t) \Vert^2 \\
    &\le \underbrace{\left( \frac{1}{2} \right)^t \left(\Vert y_0 - y_{\lambda}^*(x_0) \Vert^2 + \Vert y_0 - y^*(x_0) \Vert^2 \right) + \frac{34 L_g^2}{\mu^2} \sum_{j=0}^{t-1}  \left( \frac{1}{2}\right)^{t-1-j} \Vert x_{j+1} - x_j \Vert^2}_{:=(*)}.
\end{align*}
Plug into \Eqref{eq:err}, which, in conjunction with \Eqref{eq:des}, yields that
\begin{align*}
     \fL_{\lambda}^*(x_{t+1}) &\le \fL_{\lambda}^*(x_t) 
    - \frac{\eta}{2} \Vert \nabla \fL_{\lambda}^*(x_t) \Vert^2 - 
    \frac{1}{4\eta} \Vert x_{t+1} - x_{t} \Vert^2 + 2 \eta \times \underbrace{ \lambda^2 L_g^2 \exp \left( - \frac{\mu K}{4L_g}\right)}_{:=\gamma} \times (*).
\end{align*}
Telescoping over $t$ further yields
\begin{align} \label{eq:final}
\begin{split}
    \frac{\eta}{2}\sum_{t=0}^{T-1} \Vert \nabla \fL_{\lambda}^*(x_t) \Vert^2 &\le \fL_{\lambda}^*(x_0) - \inf_{x \in \BR^{d_x}} \fL_{\lambda}^*(x) + 4 \eta \gamma \left(\Vert y_0 - y_{\lambda}^*(x_0) \Vert^2 + \Vert y_0 - y^*(x_0) \Vert^2 \right) \\
    &\quad - \left( \frac{1}{4 \eta} - \frac{136 \eta \gamma L_g^2}{\mu^2} \right) \sum_{t=0}^{T-1} \Vert x_{t+1} - x_t \Vert^2
\end{split}
\end{align}
Let $K = \fO(\kappa \log(\lambda \kappa)) = \fO(\kappa \log(\nicefrac{\ell \kappa}{\epsilon}))$ such that $\gamma$ is sufficiently small with
\begin{align*}
    \gamma \le \min \left\{ \frac{\mu^2}{1088 \eta^2 L_g^2}, \frac{1}{4 \eta}\right\}.
\end{align*}
Then we have,
\begin{align*}
    \frac{1}{T} \sum_{t=0}^{T-1} \Vert \nabla \fL_{\lambda}^*(x_t) \Vert^2 \le  \frac{2}{\eta T} \left( \fL_{\lambda}^*(x_0) - \inf_{x \in \BR^{d_x}} \fL_{\lambda}^*(x) + \Vert y_0 - y_{\lambda}^*(x_0) \Vert^2 + \Vert y_0 - y^*(x_0) \Vert^2 \right).
\end{align*}
This concludes the proof.
% The result follows from
% Lemma \ref{lem:1st}, Proposition \ref{prop:iGD} and Proposition \ref{prop:error}.
\end{proof}

% Next, we show that the logarithmic factor in $\epsilon$ can be removed by simply
% adding a normalization step in $x$. Below, we recall the descent lemma for normalized gradient descent.

% \begin{lem}[Lemma 2 by \citet{cutkosky2020momentum}] Suppose $h(x) : \BR^d \rightarrow \BR$ is $\beta$-gradient Lipschitz. For any given stepsize $\eta>0$, consider the following update of normalized gradient descent:
% \begin{align*}
%     x_{t+1} = x_t - \eta \frac{g_t}{\Vert g_t \Vert}.
% \end{align*}
% Then it holds that
% \begin{align*}
%     h(x_{t+1}) \le h(x_t) - \frac{\eta}{3}  \Vert \nabla h(x_t) \Vert + \frac{8 \eta}{3} \Vert g_t - \nabla h(x_t) \Vert^2 + \frac{L \eta^2}{2}.
% \end{align*}
% \end{lem}

\section{Proofs of Finding Second-Order Stationarity} \label{apx:2nd}

% In this section, we focus on how to escape saddle points in bilevel optimization. 
The following theorem generalizes the result of Perturbed Gradient Descent~\citep{jin2017escape}
to Inexact Perturbed Gradient Descent by
allowing an inexact gradient $\hat \nabla h(x)$ that is close to the exact gradient $\nabla h(x)$.

\begin{algorithm*}[htbp]  
\caption{Inexact Perturbed Gradient Descent} \label{alg:iPGD}
\begin{algorithmic}[1] 
\STATE \textbf{for} $t = 0,1,\cdots,T-1$ \\
\STATE \quad \textbf{if} $ \Vert \hat \nabla h(x_t) \Vert \le \frac{4}{5} \epsilon$ \textbf{and} no perturbation added in the last $\fT$ steps \\[1mm]
\STATE \quad \quad $x_{t} = x_t - \eta \xi_t $, where $\xi_t \sim \sB(r)$ \\[1mm]
\STATE \quad \textbf{end if} \\[1mm]
\STATE \quad $x_{t+1} = x_t -  \eta  \hat \nabla h(x_t)$ \\[1mm]
\STATE \textbf{end for} \\[1mm]
\end{algorithmic}
\end{algorithm*}

\begin{thm}[Lemma 3.9 by \citet{huang2022efficiently}]
    \label{thm:iPGD}
Suppose $h(z): \BR^d \rightarrow \BR$ is
$L$-gradient Lipschitz and $\rho$-Hessian Lipschitz. 
Set the parameters in Algorithm \ref{alg:iPGD} as 
\begin{align} \label{cond:iota}
    \eta = \frac{1}{L}, \quad r = \frac{\epsilon}{400 \iota^3}, \quad \fT = \frac{L}{\sqrt{\rho \epsilon}} \cdot \iota ~~{\rm with }~~ \iota \ge 1 ~~ {\rm and} ~~ \delta \ge  \frac{L \sqrt{d}}{\sqrt{\rho \epsilon}} \iota^2 2^{8-\iota}.
\end{align}
Once the following inequality holds in each
iteration:
\begin{align} \label{cond:PF2BA}
    \Vert \hat \nabla h(x_t) - \nabla h(x_t) \Vert \le \underbrace{\min \left\{ 
    \frac{1}{20 \iota^2},  \frac{1}{16 \iota^2 2^{\iota}} \right\} \epsilon}_{:=\zeta},
\end{align}
then we can find a point $x_t$ satisfying
\begin{align*}
    \Vert \nabla h(x_t) \Vert \le \epsilon, \quad  \nabla^2 h(x_t)  \succeq - \sqrt{\rho \epsilon}~ I_{d}
\end{align*}
within $T= \fO \left( {\iota^4  \Delta L}{\epsilon^{-2}}   \right)$ iterations with probability $1-\delta$, where
$\Delta = h(x_0) - \inf_{x \in \BR^{d}} h(x)$.    
\end{thm}

Then we can show the convergence of perturbed F${}^2$BA.
\thmPFBA*
\begin{proof}
% The result follows from
% Lemma \ref{lem:1st}, Lemma \ref{lem:2rd}, Proposition \ref{prop:iPGD} and Proposition \ref{prop:error}.
Let $L$ be the gradient Lipschitz coefficient and $\rho$ be the Hessian Lipschitz coefficient of $\fL_{\lambda}^*(x)$, respectively.
According to Lemma \ref{lem:1st}, Lemma \ref{lem:2rd}, Lemma \ref{lem:yy-lambda}  and the setting of $\lambda$, we have
\begin{enumerate}[label=\alph*.]
    \item $\sup_{x \in \BR^{d_x}} \Vert \nabla \fL_{\lambda}^*(x) - \nabla \varphi(x) \Vert = \fO(\epsilon)$. 
    \item $\fL_{\lambda}^*(x_0) - \inf_{x \in \BR^{d_x}}  \fL_{\lambda}^*(x) = \fO(\Delta)$.
    \item $\Vert y_0 - y_{\lambda}^*(x_0) \Vert^2 + \Vert y_0 - y^*(x_0) \Vert^2 =\fO(R)$.
    \item $ L:=\sup_{x \in \BR^{d_x}}\Vert \nabla^2 \fL_{\lambda}^*(x) \Vert = \fO(\ell \kappa^3)$.
    \item $\rho: = \sup_{x \in \BR^{d_x}} \Vert \nabla^3 \fL_{\lambda}^*(x) \Vert = \fO(\ell \kappa^5) $.
\end{enumerate}
Then it suffices to show that the algorithm can find an $\epsilon$-second-order stationary point of $\fL_{\lambda}^*(x)$ under our choice of parameters.
We apply Theorem \ref{thm:iPGD} to prove this result. 
Recall $\zeta$ defined in \Eqref{cond:PF2BA}.
It remains to show that 
\begin{align} \label{eq:cond}
    \Vert \nabla \fL_{\lambda}^*(x_t) - \nabla \fL_{\lambda}^*(x_t) \Vert \le \zeta
\end{align}
holds for all $t$. Recall that
\begin{align*}
    \Vert \hat \nabla \fL_{\lambda}^*(x_t) - \nabla \fL_{\lambda}^*(x_t) \Vert \le 2 \lambda L_g \Vert y_{t}^K- y_{\lambda}^*(x_t) \Vert + \lambda L_g \Vert z_{t}^K - y^*(x_t) \Vert.
\end{align*} 
 If we let
\begin{align} \label{eq:Kt}
    K_t = \frac{4 L_g}{\mu} \log \left( 
    \frac{4 \lambda L_g ( \Vert y_t^0 - y_{\lambda}^*(x_t) \Vert^2 + \Vert z_t^0 - y^*(x_t) \Vert^2 )}{\zeta} 
    \right) 
\end{align}
Then by Theorem \ref{thm:GD-SC}, we have
\begin{align} \label{eq:induction}
     \max \left\{ \Vert y_t^K - y_{\lambda}^*(x_t) \Vert , ~\Vert z_t^K-  y^*(x_t) \Vert   \right\} \le \frac{\zeta}{4 \lambda L_g},
\end{align}
which implies \Eqref{eq:cond}. The only remaining issue is that the algorithm does not know the value of $\delta_t:= \Vert y_t^0 - y_{\lambda}^*(x_t) \Vert^2 + \Vert z_t^0 - y^*(x_t) \Vert^2  $ in \Eqref{eq:Kt}. We address this issue by showing a upper bound of it.
Suppose that we have already upper bounded $\delta_t$, we can then use it to upper bound $\delta_{t+1}$.
By \Eqref{eq:delta_t}, if we let $K_t \ge 8L_g/\mu$ for all $t$, then we can show that
\begin{align*} 
     \delta_{t+1}
    &\le \frac{1}{2} \delta_t +\frac{34 L_g^2}{\mu^2} \Vert x_{t+1} - x_t \Vert^2.
\end{align*}
Note that the value of $\Vert x_{t+1} - x_t \Vert^2$ is known in advance. This recursion gives a valid upper bound of $\delta_t$ that can be used to specify $K_t$ in the algorithm. 
Algorithm using this choice of $K_t$ can provably find an $\epsilon$-second-order stationary point of $\fL_{\lambda}^*(x)$ within 
\begin{align} \label{plug:2}
\sum_{t=0}^{T-1} K_t \le \frac{4L_g}{\mu} \sum_{t=0}^{T-1} \log \left( 
\frac{4 \lambda L_g \delta_t}{\zeta}
\right)   \le \frac{4L_g T}{\mu} \log \left( \frac{4 \lambda L_g \sum_{t=0}^{T-1}\delta_t}{\zeta T} \right)
\end{align}
 first-order oracle calls.
Telescoping over $t$, we obtain
\begin{align} \label{plug:1}
    \sum_{t=0}^{T-1} \delta_{t} \le 2 \delta_0 + \frac{68L_g^2}{\mu^2} \sum_{t=0}^{T-1} \Vert x_{t+1} - x_t \Vert^2.
\end{align}
By \Eqref{eq:final}, if we let $K_t \ge \Omega( \kappa \log(\lambda \kappa))$ for all $t$, we can obtain
\begin{align*}
\frac{1}{8 \eta} \sum_{t=0}^{T-1} \Vert x_{t+1} - x_t \Vert^2 \le \fL_{\lambda}^*(x_0) - \inf_{x \in \BR^{d_x}} \fL_{\lambda}^*(x) + \Vert y_0 - y_{\lambda}^*(x_0) \Vert^2 + \Vert y_0 - y^*(x_0) \Vert^2 .
\end{align*}
Plugging into \Eqref{plug:1} and then \Eqref{plug:2} yields the upper complexity bound of $\tilde \fO(\ell \kappa^4 \epsilon^{-2})$ as claimed. 
\end{proof}

% \section{Proofs of Distributed Algorithms}

% First of all, we recall
% the convergence results of gradient tracking methods ~\cite{qu2017harnessing} for strongly convex objectives.

% \begin{prop}
% Suppose that Assumption \ref{asm:W} holds.
% Consider optimizing
% \begin{align*}
%     \min_{x \in \BR^d} \varphi(x):= \frac{1}{m} \sum_{i=1}^m f_i(x),
% \end{align*}
% where each $f_i(x)$ is  $L$-smooth and $\mu$-strongly convex.
% The update 
% \begin{align*} 
% X_{t+1} &= W^S \cdot (X_t - \eta G_t) \\
% G_{t+1} &= W^S \cdot (G_t + \nabla f(X_{t+1}) - \nabla f(X_t))
% \end{align*}
% with $\eta = 1/L, ~ S = \lceil \chi \log (10\kappa)\rceil$
% can find
% \begin{align*}
%     \sqrt{m }\Vert \bar x_{t} - x^* \Vert + \Vert X_{t} - \vone \bar x_{t} \Vert_F + \eta \Vert G_{t} - \vone \bar g_{t} \Vert_F \le \zeta
% \end{align*}
% within an iteration number bounded by
% \begin{align*}
%     \fO \left( 
%     \kappa \log \left( \dfrac{\sqrt{m}\Vert \bar x_0- x^* \Vert +\Vert X_0 - \vone \bar x_{0} \Vert_F +\eta \Vert G_{0} - \vone \bar g_{0} \Vert_F }{\zeta} \right)
%     \right).
% \end{align*}
% \end{prop}

% \begin{proof}
% We begin from
% \begin{align*}
% &\quad \Vert \bar x_{t+1} - x^* \Vert \\
% &=  \Vert \bar x_{t+1} - \eta \bar g_t - x^* \Vert \\
% &\le \Vert \bar x_{t+1} - \eta  \nabla f( \bar x_t) - x^* \Vert + \eta \Vert \nabla f(\bar x_t) -  \bar g_t\Vert\\
% &\le (1 - \mu/L)\Vert \bar x_t - x^* \Vert + \frac{\eta L}{\sqrt{m}} \Vert X_t - \vone \bar x_t \Vert_F, 
% \end{align*}
% where in the last step we use Lemma 10 and Lemma 8(c) in \cite{qu2017harnessing}.

% Also, it holds that
% \begin{align} \label{eq:cons}
% \Vert X_{t+1} - \vone \bar x_{t+1} \Vert \le \theta^S \Vert X_t - \vone \bar x_t \Vert + \theta^S \eta \Vert G_t - \vone \bar g_t \Vert. 
% \end{align}
% Furthermore,
% \begin{align*}
% &~\quad \Vert G_{t+1} - \vone \bar g_{t+1} \Vert_F \\
% &~\le \theta^S \Vert G_t - \vone \bar g_t \Vert_F + \theta^S \Vert \nabla f(X_{t+1}) - \nabla f(X_t) \Vert_F \\
% &~\le \theta^S \Vert G_t - \vone \bar g_t \Vert_F + \rho L \Vert X_{t+1} - X_t \Vert_F \\
% &~\le \theta^S \Vert G_t -  \vone \bar g_t \Vert_F + 3 \theta^S L \Vert X_{t} - \vone \bar x_t \Vert_F \\
% &~\quad + 3 \theta^S L \Vert \vone \bar x_{t+1} - \vone \bar x_t \Vert_F + 3\theta^S L \Vert X_{t+1} - \vone \bar x_{t+1} \Vert_F \\
% &\overset{(\ref{eq:cons})}{\le} \rho(1+ 3 \theta^S L \eta) \Vert G_t - \vone \bar g_t \Vert_F + 4 \theta^S L \Vert X_t - \vone \bar  x_t \Vert_F + 3 \theta^S L  \sqrt{m} \Vert
% \bar x_{t+1} - \bar x_t \Vert  \\
% &~\le~ \rho(1+ 3 \theta^S L \eta) \Vert G_t -  \vone \bar g_t \Vert_F + 4 \theta^S L \Vert X_t - \vone \bar x_t \Vert_F \\
% &~\quad +3 \theta^S L \eta \sqrt{m} \Vert  \nabla f(\bar x_t ) - \bar g_t \Vert + 3 \theta^S L \eta \sqrt{m} \Vert \nabla f(\bar x_t) \Vert \\
% &~\le  ~\theta^S(1+ 3 \theta^S L \eta) \Vert G_t -  \vone \bar g_t \Vert_F + \theta^S (4 L + 3  L^2 \eta) \Vert X_t - \vone \bar x_t \Vert_F   + 3 \theta^S L^2 \eta \sqrt{m} \Vert \bar x_t - x^* \Vert,
% \end{align*}
% where in the last step we use Lemma 8(c) in \cite{qu2017harnessing}. Combing the two inequalities above, we obtain
% \begin{align*}
% &\quad \Vert X_{t+1} - \vone \bar x_{t+1} \Vert_F + \eta \Vert G_{t+1} - \vone \bar g_{t+1} \Vert_F \\
% &\le \theta^S (1 + 4 L \eta + 3 L^2 \eta^2) \Vert X_t - \vone \bar x_t \Vert_F + \theta^S  (1 + 3 \rho L \eta) \eta \Vert G_t -  \vone \bar g_t \Vert_F + 3 \theta^S L^2\eta^2 \sqrt{m} \Vert \bar x_t - x^* \Vert. 
% \end{align*}
% Combing the above inequalities, we obtain
% \begin{align*}
% &\quad \sqrt{m} \Vert \bar x_{t+1} - x^* \Vert + \Vert X_{t+1} - \vone \bar x_{t+1} \Vert_F + \eta \Vert G_{t+1} - \vone \bar g_{t+1} \Vert_F \\
% &\le (1 - \mu/ L +  3 \theta^S L^2 \eta^2   ) \sqrt{m}\Vert \bar x_t - x^* \Vert + \theta^S (1+ 5L \eta + 3 L^2 \eta^2) \Vert X_t - \vone X_t \Vert_F +  \theta^S  (1 + 3 \theta^S L \eta) \eta \Vert G_t -  \vone \bar g_t \Vert_F.
% \end{align*}
% It can be verified the choice of $S$ ensures
% \begin{align*}
%     \theta^S \le \min \left \{ \frac{\mu}{6L}, ~ \frac{1}{10} \right\},
% \end{align*}
% then we can prove the following linear convergence:
% \begin{align*}
% &\quad   \sqrt{m} \Vert \bar x_{t+1} - x^* \Vert + \Vert X_{t+1} - \vone \bar x_{t+1} \Vert_F + \eta \Vert G_{t+1} - \vone \bar g_{t+1} \Vert_F \\
% &\le \left(1-\theta^S \right) \left( \sqrt{m}\Vert \bar x_t - x^* \Vert +  \Vert X_t - \vone X_t \Vert_F +  \eta \Vert G_t -  \vone \bar g_t \Vert_F \right).
% \end{align*}
% We arrive at the conclusion by telescoping the above inequality.
% \end{proof}

% \propGTGD*

% \begin{proof}
% In our proof, we define $\zeta: = \Vert \hat \nabla \fL_{\lambda}^*(x) - \nabla \fL_{\lambda}^*(x) \Vert_F$. The descent lemma tells us that
% \begin{align*}
% \fL_{\lambda}^*(\bar x_{t+1}) &\le \fL_{\lambda}^*(\bar x_t) + \langle \nabla \fL_{\lambda}^*(\bar x_t), \bar x_{t+1} - \bar x_t \rangle + \frac{L}{2}
% \Vert \bar x_{t+1} - \bar x_t \Vert^2 \\
% &= \fL_{\lambda}^*(\bar x_t) - \frac{\eta}{2} \Vert \nabla \fL_{\lambda}^*(\bar x_t) \Vert^2 -
% \left( \frac{\eta}{2} - \frac{L \eta^2}{2} \right) \Vert \bar g_t \Vert^2 + \frac{\eta}{2} \Vert \nabla \fL_{\lambda}^*(\bar x_t) - \bar g_t \Vert^2 \\
% &\le  \fL_{\lambda}^*(\bar x_t) - \frac{\eta}{2} \Vert \nabla \fL_{\lambda}^*(\bar x_t) \Vert^2 -
% \left( \frac{\eta}{2} - \frac{L \eta^2}{2} \right) \Vert \bar g_t \Vert^2 \\
% &\quad+ \eta \left\Vert \nabla \fL_{\lambda}^*(\bar x_t) - \frac{1}{m} \mathbf{1}^\top \nabla \fL_{\lambda}^*(X_t) \right\Vert^2
% + \frac{\eta}{m} \left \Vert  \hat \nabla \fL_{\lambda}^*(X_t) - \nabla \fL_{\lambda}^*(X_t)  \right \Vert_F^2. 
% \end{align*}
% Additionally,
% \begin{align*}
% &\quad \Vert G_{t+1} - \vone \bar g_{t+1} \Vert \\
% &\le \theta^{S_{\rm out}} \Vert G_t - \vone \bar g_t \Vert + \theta^{S_{\rm out}} \Vert \hat \nabla \fL_{\lambda}^*(X_{t+1}) - \hat \nabla \fL_{\lambda}^*(X_t) \Vert_F \\
% &\le \theta^{S_{\rm out}} \Vert G_t - \vone \bar g_t \Vert_F + \theta^{S_{\rm out}} L \Vert X_{t+1} - X_t \Vert_F + 2 \theta^{S_{\rm out}}  \zeta \\
% &\le \theta^{S_{\rm out}} \Vert G_t -  \vone \bar g_t \Vert_F + 3 \theta^{S_{\rm out}} L \Vert X_{t} - \vone \bar x_t \Vert_F  \\
% &\quad + 3 \theta^{S_{\rm out}} L \sqrt{m}\Vert  \bar x_{t+1} -  \bar x_t \Vert + 3\theta^{S_{\rm out}} L \Vert X_{t+1} - \vone \bar x_{t+1} \Vert_F + 2 \theta^{S_{\rm out}}  \zeta \\
% &\le \theta^{S_{\rm out}}(1+ 3 \theta^{S_{\rm out}} L \eta) \Vert G_t -  \vone \bar g_t \Vert_F + 4 \theta^{S_{\rm out}} L \Vert X_t - \vone \bar x_t \Vert_F + 3 \theta^{S_{\rm out}}   \eta \sqrt{m} \Vert \bar  g_{t} \Vert +2 \theta^{S_{\rm out}} \zeta.
% \end{align*}
% Combing this inequality with (\ref{eq:cons}) yields
% \begin{align*}
% &\quad \Vert X_{t+1} - \vone \bar x_{t+1} \Vert_F + \eta \Vert G_{t+1} - \vone \bar g_{t+1} \Vert_F \\
% &\le \rho (1 + 4 L \eta ) \Vert X_t - \vone \bar x_t \Vert_F + \theta^{S_{\rm out}}  (1 + 3 \theta^{S_{\rm out}} L \eta) \eta \Vert G_t -  \vone \bar g_t \Vert_F + 3 \theta^{S_{\rm out}} L \eta^2 \sqrt{m} \Vert \bar g_t \Vert + 2 \theta^{S_{\rm out}} \eta \zeta. 
% \end{align*}
% Taking square,
% \begin{align*}
%     &\quad \Vert X_{t+1} - \vone \bar x_{t+1} \Vert_F^2 + \eta^2 \Vert G_{t+1} - \vone \bar g_{t+1} \Vert_F^2 \\
% &\le 4 \theta^{2S_{\rm out}} (1 + 4 L \eta)^2 \Vert X_t - \vone \bar x_t \Vert_F^2 + 4 \theta^{2S_{\rm out}} (1+ 3 \theta^{S_{\rm out}} L \eta)^2 \eta^2 \Vert G_t - \vone  \bar g_t \Vert_F^2 \\
% &\quad + 36 \theta^{2S_{\rm out}} L^2 \eta^4 m \Vert \bar g_t \Vert^2 + 16 \theta^{2S_{\rm out}} \eta^2  \zeta^2.
% \end{align*}
% Therefore,
% \begin{align*}
%     &\quad \fL_{\lambda}^*(\bar x_{t+1})  + \frac{1}{m \eta } \Vert X_{t+1} - \vone \bar x_{t+1} \Vert_F^2 + \frac{\eta}{m } \Vert G_{t+1} - \vone \bar g_{t+1} \Vert_F^2 \\
%     &\le \fL_{\lambda}^*(\bar x_t) - \frac{\eta}{2} \Vert \nabla \fL_{\lambda}^*(\bar x_t) \Vert^2 - \left( \frac{\eta}{2} - \frac{L \eta^2}{2} - 36 \theta^{2S_{\rm out}} L^2 \eta^3  \right) \Vert \bar g_t \Vert^2 \\
%     &\quad +
% \end{align*}
% \end{proof}


% \bibliographystyle{plainnat}
% \bibliography{sample}

% \clearpage
% \appendix
% 
\newpage

\section{Notations for Tensors and Derivatives} \label{apx:tensor}
We follow the notations of tensors used in \cite{kolda2009tensor}.
For a three-way tensor $\fX \in \BR^{d_1 \times d_2 \times d_3}$,
we use $ [\fX]_{i_1,i_2,i_3}$ to represent its $(i_1,i_2,i_3)$-th element.
The inner product of two three-way tensors $\fX,\fY$ is defined by $\langle \fX, \fY \rangle := \sum_{i_1,i_2,i_3} [\fX]_{i_1,i_2,i_3} \cdot [\fY]_{i_1,i_2,i_3}$. The operator norm of three-way tensor $\fX \in \BR^{d_1 \times d_2 \times d_3}$ is defined by $\Vert \fX \Vert:= \sup_{\Vert x_i \Vert =1} \langle \fX, x_1 \circ x_2 \circ x_3 \rangle $, where the elements of  $x_1 \circ x_2 \circ x_3 \in \BR^{d_1 \times d_2 \times d_3}$ is $ [x_1 \circ x_2 \circ x_3]_{i_1,i_2,i_3} := [x_1]_{i_1} \cdot  [x_2]_{i_2} \cdot [x_3]_{i_3} $. 
It can be verified that such a definition generalizes the Euclidean norm of a vector and the spectral norm of a matrix to tensors.
For a three-way tensor  $ \fX \in \BR^{d_1 \times d_2 \times d_3}$ and a vector $v \in \BR^{d_1}$, their mode-1 product, denoted by $\fX \bar \times_1 v$, gives a matrix in $\BR^{d_2 \times d_3}$ with elements $ [\fX \bar \times_1 v]_{i_2,i_3} := \sum_{i_1} [\fX]_{i_1,i_2,i_3} \cdot [v]_{i_1} $.
We define $\bar \times_2$ and $\bar \times_3$ in a similar way, and it can  be verified that 
$ \Vert \fX \bar \times_i v \Vert \le \Vert \fX \Vert \Vert v \Vert$.
For a three-way tensor $\fX \in \BR^{d_1 \times d_2 \times d_3}$ and a matrix $A \in \BR^{d_1' \times d_1}$, their mode-1 product, denoted by $ \fX \times_1 A$, gives a tensor in $\BR^{d_1' \times d_2 \times d_3}$ with elements  $ [\fX \times_1 A]_{i_1',i_2,i_3} := \sum_{i_1} [\fX]_{i_1,i_2,i_3} \cdot [A]_{i_1',i_1} $. We define $ \times_2$ and $ \times_3$ in a similar way, and it can also be verified that 
$ \Vert \fX \times_i A \Vert \le \Vert \fX \Vert \Vert A \Vert$.

For a function $h(x,y) : \BR^{d_x} \times \BR^{d_y} \rightarrow \BR$,
we denote $\nabla_x h(x,y) \in \BR^{d_x}$ to be the partial gradient with respect to $x$, with elements given by $ [\nabla_x h(x,y) ]_{i} := \partial f(x,y) / \partial [x]_{i}$. And we define $\nabla_y h(x,y) \in \BR^{d_y}$ in a similar way.  We denote $ \nabla_{xx}^2 h(x,y) \in \BR^{d_x} \times \BR^{d_x}$ to be the partial Hessian with respect to $x$, with elements given by $ [\nabla_{xx}^2 h(x,y)]_{i_1,i_2} := \partial^2 h(x,y) / (\partial [x]_{i_1} \partial [x]_{i_2})$. And we define $ \nabla_{xy}^2 h(x,y) \in \BR^{d_x} \times \BR^{d_y} $, $\nabla_{yx}^2 h(x,y) \in \BR^{d_y} \times \BR^{d_x}$ and  $ \nabla_{yy}^2 h(x,y) \in \BR^{d_y} \times \BR^{d_y}$ in a similar way.  We denote $\nabla_{xxx}^3 h(x,y) \in \BR^{d_x} \times \BR^{d_x} \times \BR^{d_x}$ to be the partial third-order derivative with respect to $x$, with elements given by 
$ [\nabla_{xxx}^3 h(x,y)]_{i_1,i_2,i_3} := \partial^3 h(x,y) / (\partial [x]_{i_1} \partial [x]_{i_2} \partial [x]_{i_3})$. And we define $\nabla_{xxy}^3 h(x,y) \in \BR^{d_x} \times \BR^{d_x} \times \BR^{d_y}$, $\nabla_{xyx}^3 h(x,y) \in \BR^{d_x} \times \BR^{d_y} \times \BR^{d_x}$,
$\nabla_{xyy}^3 h(x,y) \in \BR^{d_x} \times \BR^{d_y} \times \BR^{d_y}$,
$\nabla_{yyy}^3 h(x,y) \in \BR^{d_y} \times \BR^{d_y} \times \BR^{d_y}$,
$\nabla_{yyx}^3 h(x,y) \in \BR^{d_y} \times \BR^{d_y} \times \BR^{d_x}$,
$\nabla_{yxy}^3 h(x,y) \in \BR^{d_y} \times \BR^{d_x} \times \BR^{d_y}$
and$\nabla_{yxx}^3 h(x,y) \in \BR^{d_y} \times \BR^{d_x} \times \BR^{d_x}$
in a similar way.
We denote $\nabla y^*(x) \in \BR^{d_x} \times \BR^{d_y} $ to be the matrix with elements given by $[\nabla y^*(x) ]_{i_1,i_2} :=  \partial [y^*(x)]_{i_2}/\partial [x]_{i_1}$. We denote $\nabla^2 y^*(x) \in \BR^{d_x} \times \BR^{d_x} \times \BR^{d_y} $ to be the three-way tensor with elements given by $[\nabla^2 y^*(x)]_{i_1,i_2,i_3}: = \partial^2 [y^*(x)]_{i_3} / (\partial [x]_{i_1} \partial [x]_{i_2})$. And we define $\nabla y_{\lambda}^*(x)$ and $\nabla^2 y_{\lambda}^*(x)$ in a similar way.

\section{Lemmas for Finding First-Order Stationarity} \label{apx:lem-1st}

\begin{lem}[From \citet{kwon2023fully}] \label{lem:Lag-sc}
Under Assumption \ref{asm:sc}, for $\lambda \ge 2L_f /\mu$,   $\fL_{\lambda}(x,\,\cdot\,)$ is  $(\lambda \mu/2)$-strongly convex.  
\end{lem}
% \begin{proof}
% It can be verified that
% \begin{align*}
%     \nabla_{yy}^2 \fL(x,y) = \nabla_{yy}^2 f(x,y) +  \lambda \nabla_{yy}^2 g(x,y) \succeq \frac{\lambda \mu}{2} I_{d_y}.
% \end{align*}
% % Since  $g^*(x)$ is $(1+ L_g / \mu)L_g$-smooth, then the smoothness constant of $\fL_{\lambda}(x,y)$ can be bounded by:
% % \begin{align*}
% %     L_f + \lambda \left(2+ \frac{L_g}{\mu} \right)L_g \le \frac{4 \lambda L_g^2}{\mu^2}.
% % \end{align*}
% \end{proof}


% Then it is shown in \cite{kwon2023fully} that the following results hold. We provide the proof here for completeness.

\begin{lem} \label{lem:yy-lambda}
    Under Assumption \ref{asm:sc}, for $\lambda \ge 2L_f/ \mu$, it holds that 
    \begin{align*}
        \Vert y_{\lambda}^*(x) - y^*(x) \Vert \le \frac{C_f}{\lambda \mu}
    \end{align*}
\end{lem}

\begin{proof}
    % It is a special case of Lemma 3.2 in \cite{kwon2023fully} with $x_1 = x_2, \lambda_1 = \lambda $ and $\lambda_2 = + \infty $.
By the first-order optimality condition, we know that
\begin{align*}
    \nabla_y f(x, y_{\lambda}^*(x)) + \lambda \nabla_y g(x,y_{\lambda}^*(x))  = 0.
\end{align*}
Then we have
\begin{align*}
    \Vert y_{\lambda}^*(x) - y^*(x) \Vert 
    \le \frac{1}{\mu} \Vert \nabla_y g(x,y_{\lambda}^*(x)) \Vert = \frac{1}{\lambda \mu} \Vert \nabla_y f(x,y_{\lambda}^*(x)) \Vert \le \frac{C_f}{\lambda \mu}
\end{align*}
\end{proof}

As a direct consequence, we can show that $\fL_{x}^*$ and $\varphi(x)$ are close.
\begin{lem} \label{lem:F-Lag-close}
 Under Assumption \ref{asm:sc}, for $\lambda \ge 2L_f / \mu$, it holds that $\vert \fL_{\lambda}^*(x) - \varphi(x) \vert \le D_0/ \lambda$, where
 \begin{align*}
     D_0 : = \left(C_f + \frac{C_f L_g}{2 \mu} \right) \frac{C_f }{\mu} = \fO(\ell \kappa^2).
 \end{align*}
\end{lem}

\begin{proof}
A simple calculus shows that
    \begin{align*}
    &\quad \vert \fL_{\lambda}^*(x) - \varphi(x) \vert \\
    &\le   \vert f(x,y_{\lambda}^*(x)) - f(x,y^*(x)) \vert + \lambda \vert g(x,y_{\lambda}^*(x)) - g(x,y^*(x)) \vert \\
    &\le C_f  \Vert y_{\lambda}^*(x) - y^*(x) \Vert + \frac{\lambda L_g}{2} \Vert y_{\lambda}^*(x) - y^*(x) \Vert^2 \\
    &\le \left(C_f + \frac{C_f L_g}{2 \mu} \right) \Vert y_{\lambda}^*(x) - y^*(x) \Vert \\
    &\le \left(C_f + \frac{C_f L_g}{2 \mu} \right) \frac{C_f }{\lambda \mu}.
    \end{align*}
\end{proof}
The following result is the key to designing fully first-order methods for bilevel optimization, first proved in Lemma 3.1 in \cite{kwon2023fully}. We provide the proof here for completeness.
\begin{lem} \label{lem:key-fully}
Under Assumption \ref{asm:sc}, for $\lambda \ge 2L_f/ \mu$,  it holds that $\Vert \nabla \fL_{\lambda}^*(x) - \nabla \varphi(x) \Vert \le D_1 / \lambda$, where
\begin{align} \label{eq:D-lambda}
    D_1 :=  \left( L_f 
   + \frac{\rho_g L_g}{\mu} + \frac{C_f L_g \rho_g}{2\mu^2} + \frac{C_f \rho_g}{2\mu} 
    \right) \frac{C_f}{\mu} = \fO(\ell \kappa^3).
\end{align}
\end{lem}

\begin{proof}
Taking total derivative on $\varphi(x) = f(x,y^*(x))$, we obtain the following result:
\begin{align} \label{eq:hyper-grad}
    \nabla \varphi(x) = \nabla_x f(x,y^*(x)) - \nabla_{xy}^2 g(x,y^*(x)) [ \nabla_{yy} g(x,y^*(x))]^{-1} \nabla_y f(x,y^*(x)). 
\end{align}
Also, we know that
\begin{align*}
    \nabla \fL_{\lambda}^*(x) = \nabla_x f(x,y_{\lambda}^*(x)) + \lambda (\nabla_x g(x,y_{\lambda}^*(x)) - \nabla_x g(x,y^*(x))).
\end{align*}
By simple calculus, we have
\begin{align*}
    &\quad \nabla \fL_{\lambda}^*(x) - \nabla \varphi(x) \\
    &= \nabla_x f(x,y_{\lambda}^*(x)) - \nabla_x f(x,y^*(x)) \\
    &\quad + \nabla_{xy}^2 g(x,y^*(x)) [ \nabla_{yy} g(x,y^*(x))]^{-1} (\nabla_y f(x,y^*(x)) - \nabla_y f(x,y_{\lambda}^*(x)) \\
    &\quad + \lambda \nabla_{xy}^2 g(x,y^*(x)) [ \nabla_{yy} g(x,y^*(x))]^{-1} \times \\
    &~~~~~~~~~~~~~~~( \nabla_{yy}^2 g(x,y^*(x)) (y_{\lambda}^*(x) - y^*(x)) + \nabla_y g(x,y^*(x)) - \nabla_y g(x,y_{\lambda}^*(x)) ) \\
    &\quad + \lambda ( \nabla_x g(x,y_{\lambda}^*(x)) - \nabla_x  g(x,y^*(x)) - \nabla_{xy}^2 g(x,y^*(x)) (y_{\lambda}^*(x) - y^*(x)) ).
\end{align*}
Taking norm on both sides,
\begin{align*}
    \Vert \nabla \fL_{\lambda}^*(x) - \nabla \varphi(x) \Vert &\le L_f \Vert y_{\lambda}^*(x) - y^*(x) \Vert + \frac{\rho_g L_g}{\mu} \Vert y_{\lambda}^*(x) - y^*(x) \Vert  \\
    &\quad + \frac{\lambda L_g \rho_g}{2\mu} \Vert y_{\lambda}^*(x) - y^*(x) \Vert^2  + \frac{\lambda \rho_g}{2} \Vert y_{\lambda}^*(x) - y^*(x) \Vert^2 \\
    &\le \left( L_f 
   + \frac{\rho_g L_g}{\mu} + \frac{C_f L_g \rho_g}{2\mu^2} + \frac{C_f \rho_g}{2\mu} 
    \right) \Vert y_{\lambda}^*(x) - y^*(x) \Vert \\
    &\le \left( L_f 
   + \frac{\rho_g L_g}{\mu} + \frac{C_f L_g \rho_g}{2\mu^2} + \frac{C_f \rho_g}{2\mu} 
    \right) \frac{C_f}{\lambda \mu}
\end{align*}


\end{proof}

\begin{lem} \label{lem:yxyx-lambda}
Under Assumption \ref{asm:sc}, for $\lambda \ge 2 L_f / \mu$, it holds that $ \Vert \nabla y^*(x) - \nabla y_{\lambda}^*(x) \Vert \le D_2/\lambda $, where
\begin{align*}
    D_2 := \left( \frac{1}{\mu} + \frac{2L_g}{\mu^2} \right) \left( L_f + \frac{C_f \rho_g}{\mu} \right) = \fO(\kappa^3).
\end{align*}

\end{lem}


\begin{proof}
 Taking derivative on both sides of $\nabla_y g (x,y^*(x)) = 0$ yields
 \begin{align} \label{eq:yxyy}
     \nabla_{xy}^2 g(x,y^*(x)) +  \nabla y^*(x)\nabla_{yy}^2 g(x,y^*(x)) = 0.
 \end{align}
 Rearranging, we can obtain
\begin{align} \label{eq:y-x}
  \nabla y^*(x) = - \nabla_{xy}^2 g(x,y^*(x)) [\nabla_{yy}^2 g(x,y^*(x)) ]^{-1} . 
  \end{align}
Similarly, we also have
\begin{align} \label{eq:y-lambda-x}
    \nabla y_{\lambda}^*(x) = -  \nabla_{xy}^2 \fL_{\lambda}(x,y_{\lambda}^*(x)) [\nabla_{yy}^2 \fL_{\lambda}(x,y_{\lambda}^*(x)) ]^{-1}.
\end{align}
Using the matrix identity $ A^{-1} - B^{-1} = A^{-1} (B -A) B^{-1}$, we have
\begin{align} \label{eq:diff-inv}
    \begin{split}
    &\quad  \left \Vert [\nabla_{yy}^2 g(x,y^*(x))]^{-1} - \left [  \frac{\nabla_{yy}^2 \fL_{\lambda}(x,y_{\lambda}^*(x))}{\lambda} \right]^{-1}  \right \Vert \\
    &\le  \left \Vert [\nabla_{yy}^2 g(x,y_{\lambda}^*(x))]^{-1} \right \Vert  
    \left \Vert  \frac{\nabla_{yy}^2 \fL_{\lambda}(x,y_{\lambda}^*(x))}{\lambda} -   \nabla_{yy}^2 g(x,y^*(x))  \right \Vert \left \Vert \left[\frac{\nabla_{yy}^2 \fL_{\lambda} (x,y_{\lambda}^*(x))}{\lambda}\right]^{-1}  \right \Vert \\
    &\le \frac{2}{\mu^2} \left \Vert 
    \frac{\nabla_{yy}^2 f(x,y_{\lambda}^*(x))}{\lambda}
     + \nabla_{yy}^2 g(x,y_{\lambda}^*(x)) - \nabla_{yy}^2 g(x,y^*(x)) \right \Vert \\
     &\le \frac{2}{\mu^2} \left( \frac{L_f}{\lambda} + \rho_g \Vert y_{\lambda}^*(x) - y^*(x) \Vert  \right)\\
     &\le \frac{2}{\lambda \mu^2} \left(L_f + 
     \frac{C_f \rho_g}{\mu}
     \right). 
    \end{split}
\end{align}
Note that the setting of $\lambda \ge 2L_f / \mu  $ implies $\Vert \nabla_{xy}^2 \fL(\,\cdot\,,\,\cdot\,) \Vert \le 2 \lambda L_g$, then we further have
\begin{align*}
    &\quad \Vert \nabla y_{\lambda}^*(x) - \nabla y^*(x) \Vert \\
    & \le\left  \Vert \nabla_{xy}^2 g(x,y^*(x)) - \frac{\nabla_{xy}^2 \fL_{\lambda}(x,y_{\lambda}^*(x))}{ \lambda} \right \Vert \left \Vert [\nabla_{yy}^2 g(x,y^*(x))]^{-1}  \right \Vert  \\
    &\quad + \left \Vert \frac{\nabla_{xy}^2 \fL_{\lambda}(x,y_{\lambda}^*(x))}{\lambda}  \right \Vert  \left \Vert [\nabla_{yy}^2 g(x,y^*(x))]^{-1} - \left [  \frac{\nabla_{yy}^2 \fL_{\lambda}(x,y_{\lambda}^*(x))}{\lambda} \right]^{-1}  \right \Vert  \\
    &\le \frac{1}{\mu} \left \Vert \nabla_{xy}^2 g(x,y^*(x)) - \nabla_{xy}^2 g(x,y_{\lambda}^*(x)) - \frac{\nabla_{xy}^2 f(x,y_{\lambda}^*(x))}{\lambda}  \right \Vert  +  \frac{2L_g}{\lambda \mu^2} \left(L_f + 
     \frac{C_f \rho_g}{\mu}
     \right)  \\
    &\le \left( \frac{1}{\lambda \mu} + \frac{2L_g}{\lambda \mu^2} \right) \left( L_f + \frac{C_f \rho_g}{\mu} \right).
    \end{align*}
\end{proof}

It is clear that $\Vert \nabla y^*(x) \Vert \le L_g/\mu$, therefore $y^*(x)$ is $(L_g/\mu)$-Lipschitz. Below, we show a similar result also holds for $y_{\lambda}^*(x)$.

\begin{lem} \label{lem:nabla-y-lambda-bound}
Under Assumption \ref{asm:sc}, for $\lambda \ge 2L_f /\mu$, it holds that $\Vert \nabla y_{\lambda}^*(x) \Vert \le 4L_g/\mu$.    
\end{lem}


\begin{proof}
    Recall \Eqref{eq:y-lambda-x} that 
\begin{align*}
    \nabla y_{\lambda}^*(x) = -  \nabla_{xy}^2 \fL_{\lambda}(x,y_{\lambda}^*(x)) [\nabla_{yy}^2 \fL_{\lambda}(x,y_{\lambda}^*(x)) ]^{-1}.
\end{align*}
We arrive at the conclusion by noting 
that $  \nabla_{xy}^2 \fL_{\lambda}(\,\cdot\,,\,\cdot\,) \preceq 2 \lambda L_g$ and  $ \nabla_{yy}^2 \fL_{\lambda}(\,\cdot\, , \,\cdot\,) \succeq \lambda \mu /2 $ by Lemma \ref{lem:Lag-sc}.

\end{proof}
This implies that $y_{\lambda}^*(x)$ is $(4L_g/\mu)$-Lipschitz.

\begin{lem} \label{lem:nabla2-bound}
Under Assumption \ref{asm:sc}, for $\lambda \ge 2L_f /\mu$, $ \nabla \fL_{\lambda}^*(x)$ is $D_3$-Lipschitz, where
\begin{align*}
    D_3 &:= L_f + \frac{4L_f L_g}{\mu} +\frac{C_f \rho_g}{\mu} +  \frac{C_f L_g \rho_g}{\mu^2} + L_g D_2 = \fO(\kappa^3).
\end{align*}
\end{lem}

\begin{proof}
Note that
\begin{align} \label{eq:nabla-AB}
    \nabla \fL_{\lambda}^*(x) = \underbrace{\nabla_x f(x,y_{\lambda}^*(x))}_{A(x)} + \underbrace{\lambda ( \nabla_x g(x,y_{\lambda}^*(x)) - \nabla_x g(x,y^*(x)))}_{B(x)}. 
\end{align}
where $A(x)$ and $B(x)$ are both mappings  $\BR^{d_x} \rightarrow \BR^{d_x}$. 
From Lemma \ref{lem:nabla-y-lambda-bound} we know that $y_{\lambda}^*(x)$ is $(4L_g/\mu)$-Lipschitz. This further implies that $A(x)$ is  $(1+4L_g/\mu)L_f$-Lipschitz. Next, we bound the Lipschitz  coefficient of $B(x)$ via its derivative $\nabla B(x): \BR^{d_x} \rightarrow \BR^{d_x} \times \BR^{d_x} $, which has the following form by 
taking total derivative on $B(x)$:
\begin{align} \label{eq:nabla-Bx}
\begin{split}
    \nabla B(x) &= 
    \lambda (\nabla_{xx}^2 g(x,y_{\lambda}^*(x))-\nabla_{xx}^2 g(x,y^*(x))) \\
    &\quad + \lambda \left(\nabla y_{\lambda}^*(x)  \nabla_{yx}^2 g(x,y_{\lambda}^*(x)) - \nabla y^*(x) \nabla_{yx}^2 g(x,y^*(x)) \right).
\end{split}
\end{align}
And we can bound the operator norm of $\nabla B(x)$ by:
\begin{align*}
    \Vert \nabla B(x) \Vert &\le  \lambda \Vert \nabla_{xx}^2 g(x,y_{\lambda}^*(x)) - \nabla_{xx}^2 g(x,y^*(x)) \Vert \\
    &\quad + \lambda \Vert \nabla y^*(x) \Vert \Vert \nabla_{yx}^2 g(x,y_{\lambda}^*(x)) - \nabla_{yx}^2 g(x,y^*(x)) \Vert \\
    &\quad  + \lambda \Vert \nabla y_{\lambda}^*(x) - \nabla y^*(x) \Vert \Vert \nabla_{yx}^2 g(x,y_{\lambda}^*(x)) \Vert\\
    &\le  \lambda \rho_g \left(1+ \frac{L_g}{\mu} \right) \Vert y_{\lambda}^*(x) - y^*(x) \Vert +\lambda L_g \Vert \nabla y_{\lambda}^*(x) - \nabla y^*(x) \Vert  \\
    &\le  \left(1+ \frac{L_g}{\mu} \right) \frac{C_f \rho_g}{\mu} +L_g D_2,
\end{align*}
where we use Lemma \ref{lem:nabla-y-lambda-bound}
in the second inequality; Lemma \ref{lem:yy-lambda} and \ref{lem:yxyx-lambda} in the third one.
\end{proof}


\section{Lemmas for Finding Second-Order Stationarity} \label{apx:lem-2nd}


\begin{lem} \label{lem:nabla2-yx}
Under Assumption \ref{asm:sc} and \ref{asm:third}, 
for $\lambda \ge  2L_f / \mu$,
we have $ \Vert \nabla^2 y^*(x) -\nabla^2 y_{\lambda}^*(x)  \Vert \le D_4 /\lambda$, where 
\begin{align*}
    D_4 := \frac{2\rho_g}{\lambda \mu^2}  \left( 1 + \frac{L_g}{\mu}\right)^2  \left(L_f + 
     \frac{C_f \rho_g}{\mu}
     \right) +   \frac{14 L_g \rho_g D_2}{\lambda \mu^2} + \frac{50 L_g^2 }{\lambda \mu^3} \left( \frac{C_f \nu_g}{\mu} + \rho_f \right) = \fO(\kappa^5).
\end{align*}
\end{lem}

\begin{proof}
    By taking the derivative with respect to $x$ on
    \begin{align*}
        \nabla_{xy}^2 g(x,y^*(x)) + \nabla y^*(x) \nabla_{yy}^2 g(x,y^*(x)) =0,
    \end{align*}
    we obtain
    \begin{align*}
        & \nabla_{xxy}^3 g(x,y^*(x)) + \nabla_{yxy}^3 g(x,y^*(x)) \times_{1}  \nabla y^*(x) + \nabla^2 y^*(x) \times_3 \nabla_{yy}^2 g(x,y^*(x)) \\
        &\quad + \nabla_{xyy}^3 g(x,y^*(x)) \times_2 \nabla y^*(x)  +\nabla_{yyy}^3 g(x,y^*(x)) \times_1 \nabla y^*(x) \times_2 \nabla y^*(x) = 0.
    \end{align*}
    Rearranging to get
    {\begin{align} \label{eq:nabla2-yx}
    \begin{split}
         \nabla^2 y^*(x) 
        &= -\left( \nabla_{xxy}^3 g(x,y^*(x)) + \nabla_{yxy}^3 g(x,y^*(x)) \times_1 \nabla y^*(x) 
        \right)  \times_3 [\nabla_{yy}^2 g(x,y^*(x))]^{-1} \\
        &\quad -
        \nabla_{xyy}^3 g(x,y^*(x)) \times_2 \nabla y^*(x) \times_3 [\nabla_{yy}^2 g(x,y^*(x))]^{-1} \\
        &\quad-  \nabla_{yyy}^3 g(x,y^*(x)) \times_1 \nabla y^*(x) \times_2 \nabla y^*(x)    \times_3 [\nabla_{yy}^2 g(x,y^*(x))]^{-1}.
    \end{split}
    \end{align}}
    Similarly,
    {
    \begin{align} \label{eq:nabla2-yx-lambda}
    \begin{split}
        \nabla^2 y_{\lambda}^*(x)  &= -\left( \nabla_{xxy}^3 \fL_{\lambda} (x, y_{\lambda}^*(x)) + \nabla_{yxy}^3 \fL_\lambda(x, y_{\lambda}^*(x)) \times_1 \nabla y_{\lambda}^*(x)  \right) \times_3
        [ \nabla_{yy}^2 \fL_{\lambda}(x,y_{\lambda}^*(x))  ]^{-1} \\
        &\quad - 
        \nabla_{xyy}^3 \fL_{\lambda}(x,y_{\lambda}^*(x))
        \times_2 \times_3 [ \nabla_{yy}^2 \fL_{\lambda}(x,y_{\lambda}^*(x))]^{-1} \\
        &\quad - 
        \nabla y_{\lambda}^*(x)   +  
         \nabla_{yyy}^3 \fL_{\lambda} (x,y_{\lambda}^*(x))
        \times_1 \nabla y_{\lambda}^*(x)
        \times_2 \nabla y_{\lambda}^*(x) 
         \times_3 [ \nabla_{yy}^2 \fL_{\lambda}(x,y_{\lambda}^*(x))]^{-1}.
    \end{split}
    \end{align}}
Note that 
\begin{align*}
    \left \Vert  \nabla_{xxy}^3  g(x,y^*(x)) - \frac{\nabla_{xxy}^3 \fL_{\lambda} (x,y_{\lambda}^*(x))  }{\lambda}   \right \Vert \le \nu_g \Vert y_{\lambda}^*(x) - y^*(x) \Vert + \frac{\rho_f}{\lambda} \le \frac{1}{\lambda} \left(\frac{C_f \nu_g}{\mu} +\rho_f \right),
\end{align*}
and
\begin{align*}
    &\quad \left \Vert 
    \nabla_{yxy}^3 g(x,y^*(x)) \times_1 \nabla y^*(x)  - \frac{\nabla_{yxy}^3 \fL_{\lambda} (x,y_{\lambda}^*(x)) \times_1 \nabla y_{\lambda}^*(x) }{\lambda}
    \right \Vert \\
    &\le \Vert \nabla y^*(x) - \nabla y_{\lambda}^*(x) \Vert \Vert \nabla_{yxy}^3 g(x,y^*(x)) \Vert + \Vert \nabla y_{\lambda}^*(x) \Vert \left\Vert \nabla_{yxy}^3 g(x,y^*(x)) - \frac{\nabla_{yxy}^3 \fL_{\lambda}(x, y_{\lambda}^*(x))}{\lambda} \right \Vert \\
    &\le \frac{\rho_g D_2}{\lambda} + \frac{4L_g}{\lambda \mu} \left( \frac{C_f \nu_g}{\mu} + \rho_f \right),
\end{align*}
and
{ 
\begin{align*}
    &\quad  \left \Vert
   \nabla_{yyy}^3 g(x,y^*(x)) \times_1  \nabla y^*(x)  \times_2 \nabla y^*(x)  - \frac{\nabla_{yyy}^3 \fL_{\lambda} (x,y_{\lambda}^*(x)) \times_1 \nabla y_{\lambda}^*(x) \times_2  \nabla y_{\lambda}^*(x)} { \lambda}
    \right \Vert \\
    &\le \Vert \nabla y^*(x) \Vert \Vert \nabla_{yyy}^3 g(x,y^*(x))\Vert \Vert \nabla y^*(x) - \nabla y_{\lambda}^*(x) \Vert  \\
    &\quad + \Vert \nabla y_{\lambda}^*(x) \Vert \Vert \nabla_{yyy}^3 g(x,y^*(x))\Vert \Vert \nabla y^*(x) - \nabla y_{\lambda}^*(x) \Vert  \\
    &\quad +  \Vert \nabla y_{\lambda}^*(x) \Vert^2  \left \Vert  \nabla_{xxy}^3  g(x,y^*(x)) - \frac{\nabla_{xxy}^3 \fL_{\lambda} (x,y_{\lambda}^*(x))  }{\lambda}   \right \Vert \\
    &\le  \frac{5L_g \rho_g D_2}{\lambda \mu} + \frac{16L_g^2 }{\lambda \mu^2} \left( \frac{C_f \nu_g}{\mu} + \rho_f \right),
\end{align*}}
we can obtain that
{
\begin{align*}
    &\quad \Vert \nabla^2 y^*(x) - \nabla^2 y_{\lambda}^*(x) \Vert  \\
    &\le \rho_g \left( 1 + \frac{L_g}{\mu}\right)^2 \left \Vert [\nabla_{yy}^2 g(x,y^*(x))]^{-1} - \left [  \frac{\nabla_{yy}^2 \fL_{\lambda}(x,y_{\lambda}^*(x))}{\lambda} \right]^{-1}  \right \Vert \\
    &\quad + \left( \frac{7L_g \rho_g D_2}{\lambda \mu} + \frac{25L_g^2 }{\lambda \mu^2} \left( \frac{C_f \nu_g}{\mu} + \rho_f \right) \right) \left \Vert \left [  \frac{\nabla_{yy}^2 \fL_{\lambda}(x,y_{\lambda}^*(x))}{\lambda} \right]^{-1} \right \Vert \\
    &\le \frac{2\rho_g}{\lambda \mu^2}  \left( 1 + \frac{L_g}{\mu}\right)^2  \left(L_f + 
     \frac{C_f \rho_g}{\mu}
     \right) +   \frac{14 L_g \rho_g D_2}{\lambda \mu^2} + \frac{50 L_g^2 }{\lambda \mu^3} \left( \frac{C_f \nu_g}{\mu} + \rho_f \right) ,
\end{align*}}
where we use \Eqref{eq:diff-inv} in the second inequality.
\end{proof}

\begin{lem} \label{lem:nabla3-bound}
Under Assumption \ref{asm:sc} and \ref{asm:third}, for $ \lambda \ge 2 L_f /\mu$, 
$\nabla^2 \fL_{\lambda}^*(x) $ is $D_5$-Lipschitz,
 where
\begin{align*}
    D_5 &:=  \left( 1 + \frac{4L_g}{\mu} \right)^2 \left( 3\rho_f  + \frac{2 L_f \rho_g}{\mu} \right) + \left(1 + \frac{L_g}{\mu} \right)^2 \frac{C_f \nu_g}{\mu} \\
    &\quad + \left(2+ \frac{5L_g}{\mu} \right) D_2 \rho_g +  \left( 1+ \frac{L_g}{\mu} \right)^2 \frac{C_f \rho_g^2}{\mu^2} + L_g D_4   = \fO(\ell \kappa^5).
\end{align*}
\end{lem}

\begin{proof}
Similar to the proof of Lemma \ref{lem:nabla2-bound},  we split $\nabla^2 \fL_{\lambda}^*(x)$ into two terms:
\begin{align*}
    \nabla^2 \fL_{\lambda}^*(x) = \nabla A(x) + \nabla B(x),
\end{align*}
where both the mappings $A(x)$ and $B(x)$ both follow the definitions in \Eqref{eq:nabla-AB}. 
Taking total derivative on $A(x)$, we obtain
\begin{align*}
    \nabla A(x) = \nabla_{xx}^2 f(x,y_{\lambda}^*(x)) + \nabla y_{\lambda}^*(x) \nabla_{yx}^2 f(x,y_{\lambda}^*(x)).
\end{align*}
And recall \Eqref{eq:nabla-Bx} that 
\begin{align*}
    \nabla B(x) &= \lambda (\nabla_{xx}^2 g(x,y_{\lambda}^*(x))-\nabla_{xx}^2 g(x,y^*(x))) \\
    &\quad + \lambda \left(\nabla y_{\lambda}^*(x)  \nabla_{yx}^2 g(x,y_{\lambda}^*(x)) - \nabla y^*(x) \nabla_{yx}^2 g(x,y^*(x)) \right).
\end{align*}
Then we bound the Lipschitz coefficient of $\nabla A(x)$ and $\nabla B(x)$, respectively.

From \Eqref{eq:nabla2-yx} we can calculate that 
\begin{align} \label{bound:nabla2-yx}
    \Vert \nabla^2 y^*(x) \Vert \le \left( 1 + \frac{L_g}{\mu} \right)^2 \frac{\rho_g}{\mu}  .
\end{align}
Similarly, from \Eqref{eq:nabla2-yx-lambda} we can calculate that
\begin{align} \label{bound:nabla2-y-lambda}
    \Vert \nabla^2 y_{\lambda}^*(x)  \Vert &\le \left( 1 + \frac{4L_g}{\mu} \right)^2 \left( \frac{\rho_f}{\lambda}+ \rho_g \right) \frac{2}{\mu} \le \left( 1 + \frac{4L_g}{\mu} \right)^2 \left( \frac{2\rho_f}{L_f} + \frac{2 \rho_g}{\mu} \right). 
\end{align}
From Lemma \ref{lem:nabla-y-lambda-bound} we know that $y_{\lambda}^*(x)$ is $(4L_g/\mu)$-Lipschitz. This further implies that both $ \nabla_{xx}^2 f(x,y_{\lambda}^*(x))$ and $\nabla_{yx}^2 f(x,y_{\lambda}^*(x)) $ are $(1+ 4L_g/ \mu)\rho_f$-Lipschitz. Then, for any $x_1,x_2 \in \BR^{d_x}$, we have
\begin{align*}
    &\quad  \Vert \nabla A(x_1) - \nabla A(x_2) \Vert \\
   &\le \Vert  \nabla_{xx}^2 f(x_1,y_{\lambda}^*(x_1)) - \nabla_{xx}^2 f(x_2, y_{\lambda}^*(x_2)) \Vert  \\
   &\quad + \Vert \nabla y_{\lambda}^*(x_1) \nabla_{yx}^2 f(x_1,y_{\lambda}^*(x_1)) - \nabla y_{\lambda}^*(x_2) \nabla_{yx}^2 f(x_1,y_{\lambda}^*(x_1)) \Vert \\
   &\quad + \Vert \nabla y_{\lambda}^*(x_2) \nabla_{yx}^2 f(x_1,y_{\lambda}^*(x_1)) - \nabla y_{\lambda}^*(x_2) \nabla_{yx}^2 f(x_2,y_{\lambda}^*(x_2)) \Vert \\
   &\le \Vert \nabla_{xx}^2 f(x_1,y_{\lambda}^*(x_1)) - \nabla_{xx}^2 f(x_2, y_{\lambda}^*(x_2)) \Vert  \\
   &\quad + \Vert \nabla y_{\lambda}^*(x_1) - \nabla y_{\lambda}^*(x_2) \Vert \Vert \nabla_{yx}^2 f(x_1,y_{\lambda}^*(x_1)) \Vert \\
   &\quad + \Vert \nabla y_{\lambda}^*(x_2) \Vert \Vert \nabla_{yx}^2 f(x_1,y_{\lambda}^*(x_1)) - \nabla_{yx}^2 f(x_2,y_{\lambda}^*(x_2)) \Vert \\
   &\le \underbrace{\left(\left( 1+ \frac{4L_g}{\mu} \right)^2 \rho_f + \left( 1 + \frac{4L_g}{\mu} \right)^2 \left( \frac{2\rho_f}{L_f} + \frac{2 \rho_g}{\mu} \right) L_f \right)}_{C_1}\Vert x_1-x_2 \Vert,
\end{align*}
where $C_1$ gives the upper bound of the Lipschitz coefficient of the mapping $\nabla A(x)$.

To bound the Lipschitz coefficient of $\nabla B(x)$,
we first derive the explicit form of the mapping $\nabla^2 B(x): \BR^{d_x} \rightarrow \BR^{d_x} \times \BR^{d_x} \times \BR^{d_x} $ by:
\begin{align*}
    \nabla^2 B(x) 
    &= 
    \nabla^2 y_{\lambda}^*(x) \times_3 [\nabla_{yx}^2 f(x,y_{\lambda}^*(x))]^\top \\
    &\quad + \lambda ( \nabla_{xxx}^3 g(x,y_{\lambda}^*(x)) - \nabla_{xxx}^3 g(x,y^*(x)) ) \\
    &\quad + \lambda ( 
    \nabla_{yxx}^3 g(x,y_{\lambda}^*(x)) \times_1 \nabla y_{\lambda}^*(x) - \nabla_{yxx}^3 g(x,y^*(x)) \times_1 \nabla y^*(x) 
    ) \\
    &\quad + \lambda \left(  \nabla_{xyx}^3 g(x,y_{\lambda}^*(x)) \times_2 \nabla y_{\lambda}^*(x) - \nabla_{xyx}^3 g(x,y^*(x))   \times_2 \nabla y^*(x) \right)  \\
    &\quad + \lambda \left( 
     \nabla_{yyx}^3 g(x,y_{\lambda}^*(x)) \times_1 \nabla y_{\lambda}^*(x) \times_2  \nabla y_{\lambda}^*(x)   - 
    \nabla_{yyx}^3 g(x,y^*(x)) \times_1 \nabla y^*(x)  \times_2 \nabla y^*(x)
    \right) \\
    &\quad + \lambda \left(  \nabla^2 y_{\lambda}^*(x) \times_3 [\nabla_{yx}^2 g(x,y_{\lambda}^*(x))]^\top -  \nabla^2 y^*(x) \times_3 [\nabla_{yx}^2 g(x,y^*(x))]^\top\right)
    .
\end{align*}
Then we can bound the Lipschitz coefficient of $\nabla B(x)$
via the operator norm of $\nabla^2 B(x)$:
\begin{align*}
   C_2&:= \Vert \nabla^2 B(x) \Vert \\
    &\le 
    \Vert \nabla_{xxx}^3 g(x,y^*(x)) - 
    \nabla_{xxx}^3 g(x,y_{\lambda}^*(x))\Vert  \\
    &\quad + \lambda \Vert \nabla y^*(x) \Vert \Vert \nabla_{yxx}^3 g(x,y^*(x)) - \nabla_{yxx}^3 g(x,y_{\lambda}^*(x)) \Vert \\
    &\quad+ \lambda \Vert \nabla y_{\lambda}^*(x) - \nabla y^*(x) \Vert \Vert \nabla_{yxx}^3 g(x,y_{\lambda}^*(x)) \Vert \\
    &\quad + \lambda \Vert \nabla y^*(x) \Vert \Vert \nabla_{xyx}^3 g(x,y^*(x)) - \nabla_{xyx}^3 g(x,y_{\lambda}^*(x)) \Vert  \\
    &\quad + \lambda \Vert \nabla y^*(x) - \nabla y_{\lambda}^*(x) \Vert \Vert \nabla_{xyx}^3 g(x,y_{\lambda}^*(x)) \Vert \\
    &\quad + \lambda \Vert \nabla y^*(x) \Vert \Vert \nabla_{yyx}^3 g(x,y^*(x)) \Vert  \Vert \nabla y_{\lambda}^*(x) - \nabla y^*(x) \Vert  \\
    &\quad + \lambda \Vert \nabla y_{\lambda}^*(x) \Vert \Vert \nabla_{yyx}^3 g(x,y^*(x)) \Vert  \Vert \nabla y_{\lambda}^*(x) - \nabla y^*(x) \Vert \\
    &\quad + \lambda \Vert \nabla y^*(x) \Vert^2 \Vert \nabla_{yyx}^3 g(x,y^*(x)) - \nabla_{yyx}^3 g(x,y_{\lambda}^*(x)) \Vert \\
    &\quad + \lambda \Vert \nabla^2 y^*(x) \Vert \Vert \nabla_{yx}^2 g(x,y^*(x)) - \nabla_{yx}^2 g(x,y_{\lambda}^*(x)) \Vert \\
    &\quad + \lambda \Vert \nabla^2 y^*(x) - \nabla^2 y_{\lambda}^*(x) \Vert  \Vert \nabla_{yx}^2 g(x,y_{\lambda}^*(x)) \Vert.
\end{align*}
Then we can bound $C_2$ using the upper bound for
$\Vert \nabla y^*(x) \Vert$ given by \Eqref{eq:y-x}, that for 
$\Vert \nabla y_{\lambda}^*(x) \Vert$ by Lemma \ref{lem:nabla-y-lambda-bound}, that for $ \Vert \nabla y^*(x) - \nabla y_{\lambda}^*(x) \Vert$ by Lemma \ref{lem:yxyx-lambda}, that for $ \Vert \nabla^2 y^*(x) \Vert$ by \Eqref{bound:nabla2-yx} and that for $ \Vert \nabla^2 y^*(x) - \nabla^2 y_{\lambda}^*(x) \Vert$ by \Eqref{lem:nabla2-yx}. 
Finally, we can obtain 
\begin{align*}
    C_1+C_2 & \le \left( 1 + \frac{4L_g}{\mu} \right)^2 \left( 3\rho_f  + \frac{2 L_f \rho_g}{\mu} \right) + \left(1 + \frac{L_g}{\mu} \right)^2 \frac{C_f \nu_g}{\mu} \\
    &\quad + \left(2+ \frac{5L_g}{\mu} \right) D_2 \rho_g +  \left( 1+ \frac{L_g}{\mu} \right)^2 \frac{C_f \rho_g^2}{\mu^2} + L_g D_4.
\end{align*}
\end{proof}
 
\begin{lem} \label{lem:nabla2-phi}
Under Assumption \ref{asm:sc} and \ref{asm:third}, for $\lambda \ge 2 L_f /\mu$, we have $ \Vert \nabla^2 \fL_{\lambda}^*(x) - \nabla^2 \varphi(x) \Vert \le  D_6 /\lambda$,
where
\begin{align*}
    D_6 := 2L_g D_2^2 + \left( 1 + \frac{L_g}{\mu}\right)^2\left( \frac{C_f \rho_f}{\mu} + \frac{C_f L_f \rho_g}{\mu^2} + \frac{C_f^2 \nu_g}{2 \mu^2} + \frac{C_f^2 \rho_g^2}{2 \mu^3}  \right)  = \fO(\ell \kappa^6).
\end{align*}
\end{lem}

\begin{proof}
    Taking total derivative on $ \nabla \varphi(x) = \nabla_x f(x,y^*(x)) + \nabla y^*(x) \nabla_y f(x,y^*(x))$ yields
    \begin{align} \label{eq:phi-hess}
\begin{split}
    \nabla^2 \varphi(x) &= \nabla_{xx}^2 f(x,y^*(x)) + \nabla y^*(x) \nabla_{yx}^2 f(x,y^*(x))  + \nabla^2 y^*(x) \bar \times_3 \nabla_y f(x,y^*(x)) \\
    &\quad + \nabla_{xy}^2f(x,y^*(x)) [\nabla y^*(x)]^\top + \nabla y^*(x) \nabla_{yy}^2 f(x,y^*(x)) [\nabla y^*(x)]^\top.
\end{split}
\end{align}
Plug into the close form of $\nabla^2 y^*(x)$ given by \Eqref{eq:nabla2-yx}, we arrive at
    \begin{align*}
    &\quad \nabla^2 \varphi(x) \\
    &= \underbrace{\nabla_{xx }^2  f(x,y^*(x)) - \nabla_{xxy}^3 g(x,y^*(x)) \times_3 [  \nabla_{yy}^2 g(x,y^*(x))]^{-1} \bar \times_3 \nabla_y f(x,y^*(x)) }_{\rm (I)} \\
    &\quad +  
    \underbrace{  \left(\nabla_{yx}^2 f(x,y^*(x)) - \nabla_{yxy}^3 g(x,y^*(x)) \times_3 [\nabla_{yy}^2 g(x,y^*(x)]^{-1}  \bar \times_3 \nabla_y f(x,y^*(x)) \right) \times_1 \nabla y^*(x)  }_{\rm (II)}  \\
    &\quad +  
    \underbrace{  \left(\nabla_{xy}^2 f(x,y^*(x)) - \nabla_{xyy}^3 g(x,y^*(x)) \times_3 [\nabla_{yy}^2 g(x,y^*(x)]^{-1}  \bar \times_3 \nabla_y f(x,y^*(x)) \right) \times_2 \nabla y^*(x)  }_{\rm (III)}  \\
    &\quad + \underbrace{  \left( \nabla_{yy}^2 f(x,y^*(x)) -  \nabla_{yyy}^3 g(x,y^*(x)) \times_3 [\nabla_{yy}^2 g(x,y^*(x)]^{-1} \bar \times_3 \nabla_y f(x,y^*(x))   \right) \times_1 \nabla y^*(x) \times_2 \nabla y^*(x) }_{\rm (IV)}. 
\end{align*}
Recall that
\begin{align*}
    \nabla^2 \fL_{\lambda}^*(x) &= \nabla_{xx}^2  f(x,y_{\lambda}^*(x)) +  \lambda ( \nabla_{xx}^2 g(x,y_{\lambda}^*(x)) - \nabla_{xx}^2 g(x,y^*(x))) \\
    &\quad + \nabla y_{\lambda}^*(x) \nabla_{yx}^2 \fL_{\lambda}(x,y_{\lambda}^*(x))
    - \lambda \nabla y^*(x) \nabla_{yx}^2 g(x,y^*(x)) .
\end{align*}
Let
\begin{align*}
    \tilde \nabla^2 \fL_{\lambda}^*(x) &= \underbrace{\nabla_{xx}^2  f(x,y_{\lambda}^*(x)) +  \lambda ( \nabla_{xx}^2 g(x,y_{\lambda}^*(x)) - \nabla_{xx}^2 g(x,y^*(x)))}_{\rm (I')}\\
    &+  \underbrace{\nabla y^*(x) \left(\nabla_{yx}^2 \fL_{\lambda}(x,y_{\lambda}^*(x))- \lambda \nabla_{yx}^2 g(x,y^*(x))\right) }_{\rm (II')}\\ 
    & + \underbrace{ \left(\nabla_{xy}^2 \fL_{\lambda}(x,y_{\lambda}^*(x))- \lambda \nabla_{xy}^2 g(x,y^*(x))\right)[\nabla y^*(x)]^\top }_{\rm (III')} \\
    &+ \underbrace{\nabla y^*(x) \left(\nabla_{yy}^2 \fL_{\lambda} (x,y_{\lambda}^*(x))   - \lambda \nabla_{yy}^2 g(x,y^*(x)) \right) [\nabla y^*(x)]^\top }_{\rm (IV')}. 
\end{align*}
Recall \Eqref{eq:y-lambda-x} that
\begin{align*}
    \nabla y_{\lambda}^*(x) = - \nabla_{xy}^2 \fL_{\lambda}(x,y_{\lambda}^*(x)) [\nabla_{yy}^2 \fL_{\lambda}(x,y_{\lambda}^*(x)) ]^{-1}.
\end{align*}
Then using
equivalent transformation and the identity
\begin{align*}
    U A U^\top - V A V^\top = (U-V) A V^\top + VA (U-V)^\top + (U-V)A (U-V)^\top
\end{align*}
for a matrices $U,V$ and $A$,
we can show that
\begin{align*}
    &\quad \tilde \nabla^2 \fL_{\lambda}^*(x) - \nabla^2 \fL_{\lambda}^*(x) \\
    &= ( \nabla y^*(x) - \nabla y_{\lambda}^*(x) ) \nabla_{yx}^2 \fL_{\lambda} (x,y_{\lambda}^*(x)) + \\
    &\quad + \nabla_{xy}^2 \fL_{\lambda}(x,y_{\lambda}^*(x)) [ \nabla y^*(x)]^\top + \nabla y^*(x) \nabla_{yy}^2 \fL_{\lambda} (x,y_{\lambda}^*(x)) [\nabla y^*(x)]^\top \\
    &= ( \nabla y^*(x) - \nabla y_{\lambda}^*(x) ) \nabla_{yx}^2 \fL_{\lambda} (x,y_{\lambda}^*(x)) + \nabla_{xy}^2 \fL_{\lambda}(x,y_{\lambda}^*(x)) [ \nabla y^*(x) - \nabla y_{\lambda}^*(x)]^\top \\
    &\quad + \nabla y^*(x) \nabla_{yy}^2 \fL_{\lambda} (x,y_{\lambda}^*(x)) [\nabla y^*(x)]^\top - \nabla y_{\lambda}^*(x) \nabla_{yy}^2 \fL_{\lambda} (x,y_{\lambda}^*(x)) [\nabla y_{\lambda}^*(x)]^\top \\
    &= ( \nabla y^*(x) - \nabla y_{\lambda}^*(x) ) \nabla_{yx}^2 \fL_{\lambda} (x,y_{\lambda}^*(x)) + \nabla_{xy}^2 \fL_{\lambda}(x,y_{\lambda}^*(x)) [ \nabla y^*(x) - \nabla y_{\lambda}^*(x)]^\top \\ 
    &\quad + (\nabla y^*(x) - \nabla y_{\lambda}^*(x)) \nabla_{yy}^2 \fL_{\lambda} (x,y_{\lambda}^*(x)) [ \nabla y_{\lambda}^*(x)]^\top \\
    &\quad + \nabla y_{\lambda}^*(x) \nabla_{yy}^2 \fL_{\lambda} (x,y_{\lambda}^*(x)) (\nabla y^*(x) - \nabla y_{\lambda}^*(x))^\top \\
    &\quad + (\nabla y^*(x) - \nabla y_{\lambda}^*(x)) \nabla_{yy}^2 \fL_{\lambda} (x,y_{\lambda}^*(x)) [ \nabla y^*(x) - \nabla y_{\lambda}^*(x)]^\top \\
    &= (\nabla y^*(x) - \nabla y_{\lambda}^*(x)) \nabla_{yy}^2 \fL_{\lambda} (x,y_{\lambda}^*(x)) [ \nabla y^*(x) - \nabla y_{\lambda}^*(x)]^\top.
\end{align*}
Noting that $ \nabla_{yy}^2 \fL_{\lambda}(\,\cdot\, , \,\cdot\,) \preceq 2 \lambda L_g$, we have
\begin{align*}
    \Vert \tilde \nabla^2 \fL_{\lambda}^*(x) - \nabla^2 \fL_{\lambda}^*(x) \Vert \le 2 \lambda L_g \Vert \nabla y^*(x) - \nabla y_{\lambda}^*(x) \Vert^2 \le \frac{2L_g D_2^2}{\lambda }. 
\end{align*}
Next, we bound difference between 
$\tilde \nabla^2 \fL_{\lambda}^*(x) $ and $\nabla^2 \fL_{\lambda}^*(x)$. For the difference between ${\rm (I)}$ and ${\rm (I')}$, 
{\begin{align*}
    &\quad \rm (I) - {\rm (I')} \\
    &= \nabla_{xx}^2 f(x,y^*(x)) - \nabla_{xx}^2 f(x,y_{\lambda}^*(x))  \\
    & \quad + \nabla_{xxy}^3 g(x,y^*(x)) \times_3 [\nabla_{yy}^2 g(x,y^*(x))]^{-1} \bar \times_3 ( \nabla_y f(x,y_{\lambda}^*(x)) - \nabla_y f(x,y^*(x))) \\
    &\quad + \lambda \left( \nabla_{xx}^2 g(x,y^*(x)) -  \nabla_{xx}^2 g(x,y_{\lambda}^*(x))  + \nabla_{xxy}^3  g(x,y^*(x)) \bar \times_3
    (y_{\lambda}^*(x) - y^*(x)) \right) \\
    &\quad + \lambda \nabla_{xxy}^3 g(x,y^*(x)) \times_3 [\nabla_{yy}^2 g(x,y^*(x))]^{-1} 
 \\
 &~~~~~~~~~~~~~~~~~~~~~~~~~~~\bar \times_3 \left(\nabla_y g(x,y_{\lambda}^*(x)) 
    - \nabla_y g(x,y^*(x)) - \nabla_{yy}^2 g(x,y^*(x)) (y_{\lambda}^*(x) -y^*(x))  \right).
\end{align*}}
Therefore, 
\begin{align*}
    \Vert {\rm (I) } - {\rm (I')} \Vert &\le \rho_f \Vert y^*(x) - y_{\lambda}^*(x) \Vert  + \frac{L_f \rho_g}{\mu} \Vert y^*(x) - y_{\lambda}^*(x) \Vert \\
    &\quad + \frac{\lambda \nu_g}{2} \Vert y^*(x) - y_{\lambda}^*(x) \Vert^2 + \frac{\lambda \rho_g^2}{2\mu} \Vert y^*(x) - y_{\lambda}^*(x) \Vert^2  \\
    &\le \frac{C_f \rho_f}{\lambda \mu} + \frac{C_f L_f \rho_g}{\lambda \mu^2} +  \frac{C_f^2 \nu_g}{2\lambda \mu^2}+ \frac{C_f^2 \rho_g^2}{2 \lambda \mu^3} .
\end{align*}
Using $\Vert \nabla y^*(x) \Vert \le  L_g / \mu$, we can similarly bound the difference between ${\rm (II)}$ and ${\rm (II')}$ by
{ \begin{align*}
    &\quad \lVert{\rm (II)} - {\rm (II')}\rVert \\
    &\le 
      \Vert \nabla y^*(x) \Vert \Big \Vert  \nabla_{yx}^2 f(x,y^*(x)) - \nabla_{yxy}^3 g(x,y^*(x)) \times_3 [\nabla_{yy}^2 g(x,y^*(x)]^{-1} \bar \times_3 \nabla_y f(x,y^*(x))\\
     &\quad - \nabla_{yx}^2 f(x,y_{\lambda}^*(x))-\lambda  \left(\nabla_{yx}^2 g(x,y_{\lambda}^*(x))-\nabla_{yx}^2 g(x,y^*(x)) \right) \Big \Vert \\
     & \le \frac{L_g }{\mu} \left( \frac{C_f \rho_f}{\lambda \mu} + \frac{C_f L_f \rho_g}{\lambda \mu^2}  + \frac{C_f^2 \nu_g}{2\lambda \mu^2} +\frac{C_f^2 \rho_g^2}{2 \lambda \mu^3} \right),
\end{align*}}
And the bound for ${\rm(III)}$ and $ {\rm (III')}$ is the same. Finally, we know that
\begin{align*}
    &\quad \left \Vert {\rm (IV)} - {\rm (IV')} \right \Vert \\
    &\le \Vert \nabla y^*(x) \Vert^2 \Big \Vert \nabla_{yy}^2 f(x,y^*(x)) - \nabla_{yyy}^3 g(x,y^*(x)) \times_3 [ \nabla_{yy}^2 g(x,y^*(x))]^{-1} \bar \times_3 \nabla_y f(x,y^*(x))  \\
    &\quad - \nabla_{yy}^2 f(x,y_{\lambda}^*(x)) - \lambda \left( \nabla_{yy}^2 g(x,y_{\lambda}^*(x)) - \nabla_{yy}^2 g(x,y^*(x)) \right) \Big \Vert \\
    &\le \frac{L_g^2}{\mu^2} \left( \frac{C_f \rho_f}{\lambda \mu} + \frac{C_f L_f \rho_g}{\lambda \mu^2}  + \frac{C_f^2 \nu_g}{2\lambda \mu^2} +  \frac{C_f^2 \rho_g^2}{2 \lambda \mu^3}\right).
\end{align*}
Combing the above bounds completes our proof.
\end{proof}

\section{Proofs of Finding First-Order Stationarity} \label{apx:1st}

We recall the standard result for gradient descent for strongly convex functions, which is used for proving the linear convergence of $y_t^k \rightarrow y_{\lambda}^*(x_t) $ and $z_{t}^k \rightarrow y^*(x_t)$.
\begin{thm}[Theorem 3.10 by \citet{bubeck2015convex}] \label{thm:GD-SC}
Suppose $h(x): \BR^d \rightarrow \BR$ is  $\beta$-gradient Lipschitz and $\alpha$-strongly convex. Consider the following update of gradient descent:
\begin{align*}
    x_{t+1} = x_t - \frac{1}{\beta} \nabla h(x_t).
\end{align*}
Let $x^* = \arg \min_{x \in \BR^d} h(x)$. Then it holds that
\begin{align*}
    \Vert x_{t+1} - x^* \Vert^2 \le \left(1 - \frac{\alpha}{\beta} \right) \Vert x_t - x^* \Vert^2.
\end{align*}
\end{thm}

%\propGD*

% \begin{proof}
% It is a straightforward result of the following descent lemma for gradient descent.
% \begin{align*}
%     \fL_{\lambda}^*(x_{t+1}) &\le \fL_{\lambda}^*(x_t) + \langle \nabla \fL_{\lambda}^*(x_t) , x_{t+1} - x_t
%     \rangle + \frac{L}{2} \Vert x_{t+1 } - x_t \Vert^2 \\
%     &= \fL_{\lambda}^*(x_t) -
%     \frac{\eta}{2} \Vert \nabla \fL_{\lambda}^*(x_t) \Vert^2 -
%     \left( \frac{\eta}{2} - \frac{\eta^2 L}{2} \right) \Vert \hat \nabla \fL_{\lambda}^*(x_t) \Vert^2 + \frac{\eta}{2} \Vert \hat  \nabla \fL_{\lambda}^*(x_t) - \nabla \fL_{\lambda}^*(x_t) \Vert^2 \\
%     &\le \fL_{\lambda}^*(x_t) 
%     - \frac{\eta}{2} \Vert \nabla \fL_{\lambda}^*(x_t) \Vert^2 
%     + \frac{\eta}{2} \Vert  \hat \nabla \fL_{\lambda}^*(x_t) - \nabla \fL_{\lambda}^*(x_t) \Vert^2.
% \end{align*}
% A simple telescoping completes the proof.
% \end{proof}

% \propPGD*

% \begin{proof}
%     It is the same as the proof of Theorem 2.7 in \cite{huang2022efficiently}.
% \end{proof}



% \properror*

% \begin{proof}
% According to
% \begin{align*}
%     \Vert \hat \nabla \fL_{\lambda}^*(x_t) - \nabla \fL_{\lambda}^*(x_t) \Vert \le (L_f+ \lambda L_g) \Vert y_{t+1}- y_{\lambda}^*(x_t) \Vert + \lambda L_g \Vert z_{t+1} - y^*(x_t) \Vert,
% \end{align*}
% it suffices to let 
% \begin{align*}
% \Vert y_{t+1} - y_{\lambda}^*(x_t) \Vert & \le \frac{\zeta}{4\lambda L_g}   ~{\rm and}~
% \Vert z_{t+1} - y^*(x_t) \Vert \le \frac{\zeta}{2 \lambda L_g}
% \end{align*}
% holds for all $t$. Below, we only show how to control $\Vert y_{t+1} - y_{\lambda}^*(x_t) \Vert$ since the analysis for $ \Vert z_{t+1} - y^*(x_t) \Vert$ is very similar. 
% Since $\fL_{\lambda}(x,y)$ is   $(\lambda \mu /2)$-strongly convex and $(2 \lambda L_g)$-smooth in $y$.
% From Theorem 2.1.14 in \cite{nesterov2018lectures}, we know that
% \begin{align*}
%     \Vert y_{t+1} - y_{\lambda}^*(x_t) \Vert \le \left( 1 - \frac{\mu}{4L_g} \right)^{K} \Vert y_t - y^*(x_t) \Vert.
% \end{align*}
% This indicates that we can  set
% \begin{align*}
%     K_t \ge \frac{4 L_g}{\mu} \log \left( 
%     \frac{4 \lambda L_g \Vert y_t - y_{\lambda}^*(x_t) \Vert}{\zeta}
%     \right), \quad \forall 0 \le t \le T-1.
% \end{align*}
% Next, we need to bound $\Vert y_t - y_{\lambda}^*(x_t) \Vert$ for all $t$. For the case of $t=0$, we have
% \begin{align*}
%     \Vert y_0 - y_{\lambda}^*(x_0) \Vert \le \Vert y_0 - y^*(x_0) \Vert + \Vert y^*(x_0) - y_{\lambda}^*(x_0) \Vert \le \Vert y_0 - y^*(x_0) \Vert + \frac{C_f}{\lambda \mu}.
% \end{align*}
% For the case of $t \ge 1$, we can prove it by induction. Suppose that we have $\Vert y_{t} - y_{\lambda}^*(x_{t-1})\Vert \le \zeta / (4 \lambda L_g)$, then
% \begin{align*}
%     \Vert y_t - y_{\lambda}^*(x_t) \Vert  &\le  \Vert y_{t} - y_{\lambda}^*(x_{t-1}) \Vert + \Vert y_{\lambda}^*(x_{t}) - y_{\lambda}^*(x_{t-1}) \Vert \\
%     &\le \frac{\zeta}{4 \lambda L_g} + \frac{4 L_g}{\mu} \Vert x_t - x_{t-1} \Vert \\
%     &\le \frac{\zeta}{4 \lambda L_g} + \frac{4 L_g \eta}{\mu} \Vert \hat \nabla \fL_{\lambda}^*(x_{t-1})\Vert \\
%     &\le \frac{\zeta}{4 \lambda L_g} + \frac{4 L_g \eta \zeta}{\mu} + \frac{4 L_g \eta}{\mu} \Vert \nabla \fL_{\lambda}^*(x_{t-1})\Vert,
% \end{align*}
% where the second inequality uses the fact that $y_{\lambda}^*(\,\cdot\,)$ is $(4 L_g/\mu)$-Lipschitz.
%     Recall that
%     \begin{align*}
%         \nabla \fL_{\lambda}^*(x) = f(x,y_{\lambda}^*(x)) + \lambda ( 
%         g(x,y_{\lambda}^*(x) ) - g(x,y^*(x))
%         ).
%     \end{align*}
%     According to Lemma \ref{lem:yy-lambda}, we have
%     \begin{align*}
%         \Vert \nabla \fL_{\lambda}^*(x) \Vert \le C_f + \lambda L_g \Vert y_{\lambda}^*(x) - y^*(x) \Vert \le \left( 1 + \frac{L_g}{\mu}\right) C_f.
%     \end{align*}
% Thus,
% \begin{align*}
%     \Vert y_t - y_{\lambda}^*(x_t) \Vert \le \frac{\zeta}{4 \lambda L_g} + \frac{4 L_g \eta }{\mu} \left( \zeta +   \left( 1 + \frac{L_g}{\mu}\right) C_f\right).
% \end{align*}

% \end{proof}

Below, we prove that F${}^2$BA is indeed near-optimal.

\thmFBA* 
\begin{proof}
Let $L$ be the gradient Lipschitz coefficient of $\fL_{\lambda}^*(x)$.
According to Lemma \ref{lem:1st} and the setting of $\lambda$, we have
\begin{enumerate}[label=\alph*.]
    \item $\sup_{x \in \BR^{d_x}} \Vert \nabla \fL_{\lambda}^*(x) - \nabla \varphi(x) \Vert = \fO(\epsilon)$. 
    \item $\fL_{\lambda}^*(x_0) - \inf_{x \in \BR^{d_x}}  \fL_{\lambda}^*(x) = \fO(\Delta)$.
    \item $ L:=\sup_{x \in \BR^{d_x}}\Vert \nabla^2 \fL_{\lambda}^*(x) \Vert = \fO(\ell \kappa^3)$.
\end{enumerate}
Due to Lemma \ref{lem:yy-lambda}, we also have
\begin{align*}
    \Vert y_0 - y_{\lambda}^*(x_0) \Vert^2 + \Vert y_0 - y^*(x_0) \Vert^2 =\fO(R). 
\end{align*}
Now it suffices to show that the algorithm can find an $\epsilon$-first-order stationary of $\fL_{\lambda}^*(x)$. 

We begin from the following descent lemma for gradient descent. Let $\eta \le 1/(2L)$, then
\begin{align} \label{eq:des}
\begin{split}
    \fL_{\lambda}^*(x_{t+1}) &\le \fL_{\lambda}^*(x_t) + \langle \nabla \fL_{\lambda}^*(x_t) , x_{t+1} - x_t
    \rangle + \frac{L}{2} \Vert x_{t+1 } - x_t \Vert^2 \\
    &= \fL_{\lambda}^*(x_t) -
    \frac{\eta}{2} \Vert \nabla \fL_{\lambda}^*(x_t) \Vert^2 -
    \left( \frac{\eta}{2} - \frac{\eta^2 L}{2} \right) \Vert \hat \nabla \fL_{\lambda}^*(x_t) \Vert^2 + \frac{\eta}{2} \Vert \hat  \nabla \fL_{\lambda}^*(x_t) - \nabla \fL_{\lambda}^*(x_t) \Vert^2 \\
    &\le \fL_{\lambda}^*(x_t) 
    - \frac{\eta}{2} \Vert \nabla \fL_{\lambda}^*(x_t) \Vert^2 - 
    \frac{1}{4\eta} \Vert x_{t+1} - x_{t} \Vert^2
    + \frac{\eta}{2} \Vert  \hat \nabla \fL_{\lambda}^*(x_t) - \nabla \fL_{\lambda}^*(x_t) \Vert^2.
\end{split}
\end{align}
Note that
\begin{align*}
    \Vert \hat \nabla \fL_{\lambda}^*(x_t) - \nabla \fL_{\lambda}^*(x_t) \Vert \le 2 \lambda L_g \Vert y_{t}^K- y_{\lambda}^*(x_t) \Vert + \lambda L_g \Vert z_{t}^K - y^*(x_t) \Vert.
\end{align*}
By Theorem \ref{thm:GD-SC}, we have
\begin{align*}
    \Vert y_t^K - y_{\lambda}^*(x_t) \Vert^2 &\le \exp\left( - \frac{\mu K}{4 L_g} \right) \Vert y_t^0 - y_{\lambda}^*(x_t) \Vert^2 \\
     \Vert z_t^K - y^*(x_t) \Vert^2 &\le \exp\left( - \frac{\mu K}{L_g} \right) \Vert z_t^0 - y^*(x_t) \Vert^2
\end{align*}
Therefore, we have
\begin{align} \label{eq:err}
    \Vert \hat \nabla \fL_{\lambda}^*(x_t)- \nabla \fL_{\lambda}^*(x_t) \Vert^2 \le 4 \lambda^2 L_g^2 \exp\left( - \frac{\mu K}{4 L_g} \right) \left( \Vert y_t^0 - y_{\lambda}^*(x_t) \Vert^2 + \Vert z_t^0 - y^*(x_t) \Vert^2 \right)
\end{align}
By Young's inequality, 
\begin{align*}
    \Vert y_{t+1}^0 - y_{\lambda}^*(x_{t+1}) \Vert^2 & \le 2 \Vert y_t^K - y_{\lambda}^*(x_{t}) \Vert^2 + 2 \Vert y_{\lambda}^*(x_{t+1}) - y_{\lambda}^*(x_t) \Vert^2  \\
    &\le 2 \exp \left( - \frac{\mu K}{4 L_g} \right) \Vert y_t^0 - y_{\lambda}^*(x_t) \Vert^2 + \frac{32 L_g^2}{\mu^2} \Vert x_{t+1} - x_t \Vert^2,
\end{align*}
where the second inequality follows Lemma \ref{lem:nabla-y-lambda-bound} that $y_{\lambda}^*(x)$ is $(4 L_g/\mu)$-Lipschitz.  Similarly, we can derive the recursion about $\Vert z_t^0 - y^*(x_t) \Vert^2$. 

Put them together and let $K \ge 8 L_g/ \mu$, we have
\begin{align} \label{eq:delta_t}
    &\quad \Vert y_{t+1}^0 - y_{\lambda}^*(x_{t+1}) \Vert^2  + \Vert z_{t+1}^0 - y^*(x_{t+1}) \Vert^2 \\
    &\le 2 \exp \left( - \frac{\mu K}{4 L_g} \right) \left(\Vert y_t^0 - y_{\lambda}^*(x_t) \Vert^2 + \Vert z_t^0 - y^*(x_t) \Vert^2 \right) +\frac{34 L_g^2}{\mu^2} \Vert x_{t+1} - x_t \Vert^2 \\
    &\le \frac{1}{2} \left(\Vert y_t^0 - y_{\lambda}^*(x_t) \Vert^2 + \Vert z_t^0 - y^*(x_t) \Vert^2 \right) +\frac{34 L_g^2}{\mu^2} \Vert x_{t+1} - x_t \Vert^2.
\end{align}
Telescoping over $t$ yields
\begin{align*}
    &\quad \Vert y_t^0 - y_{\lambda}^*(x_t) \Vert^2 + \Vert z_t^0 - y^*(x_t) \Vert^2 \\
    &\le \underbrace{\left( \frac{1}{2} \right)^t \left(\Vert y_0 - y_{\lambda}^*(x_0) \Vert^2 + \Vert y_0 - y^*(x_0) \Vert^2 \right) + \frac{34 L_g^2}{\mu^2} \sum_{j=0}^{t-1}  \left( \frac{1}{2}\right)^{t-1-j} \Vert x_{j+1} - x_j \Vert^2}_{:=(*)}.
\end{align*}
Plug into \Eqref{eq:err}, which, in conjunction with \Eqref{eq:des}, yields that
\begin{align*}
     \fL_{\lambda}^*(x_{t+1}) &\le \fL_{\lambda}^*(x_t) 
    - \frac{\eta}{2} \Vert \nabla \fL_{\lambda}^*(x_t) \Vert^2 - 
    \frac{1}{4\eta} \Vert x_{t+1} - x_{t} \Vert^2 + 2 \eta \times \underbrace{ \lambda^2 L_g^2 \exp \left( - \frac{\mu K}{4L_g}\right)}_{:=\gamma} \times (*).
\end{align*}
Telescoping over $t$ further yields
\begin{align} \label{eq:final}
\begin{split}
    \frac{\eta}{2}\sum_{t=0}^{T-1} \Vert \nabla \fL_{\lambda}^*(x_t) \Vert^2 &\le \fL_{\lambda}^*(x_0) - \inf_{x \in \BR^{d_x}} \fL_{\lambda}^*(x) + 4 \eta \gamma \left(\Vert y_0 - y_{\lambda}^*(x_0) \Vert^2 + \Vert y_0 - y^*(x_0) \Vert^2 \right) \\
    &\quad - \left( \frac{1}{4 \eta} - \frac{136 \eta \gamma L_g^2}{\mu^2} \right) \sum_{t=0}^{T-1} \Vert x_{t+1} - x_t \Vert^2
\end{split}
\end{align}
Let $K = \fO(\kappa \log(\lambda \kappa)) = \fO(\kappa \log(\nicefrac{\ell \kappa}{\epsilon}))$ such that $\gamma$ is sufficiently small with
\begin{align*}
    \gamma \le \min \left\{ \frac{\mu^2}{1088 \eta^2 L_g^2}, \frac{1}{4 \eta}\right\}.
\end{align*}
Then we have,
\begin{align*}
    \frac{1}{T} \sum_{t=0}^{T-1} \Vert \nabla \fL_{\lambda}^*(x_t) \Vert^2 \le  \frac{2}{\eta T} \left( \fL_{\lambda}^*(x_0) - \inf_{x \in \BR^{d_x}} \fL_{\lambda}^*(x) + \Vert y_0 - y_{\lambda}^*(x_0) \Vert^2 + \Vert y_0 - y^*(x_0) \Vert^2 \right).
\end{align*}
This concludes the proof.
% The result follows from
% Lemma \ref{lem:1st}, Proposition \ref{prop:iGD} and Proposition \ref{prop:error}.
\end{proof}

% Next, we show that the logarithmic factor in $\epsilon$ can be removed by simply
% adding a normalization step in $x$. Below, we recall the descent lemma for normalized gradient descent.

% \begin{lem}[Lemma 2 by \citet{cutkosky2020momentum}] Suppose $h(x) : \BR^d \rightarrow \BR$ is $\beta$-gradient Lipschitz. For any given stepsize $\eta>0$, consider the following update of normalized gradient descent:
% \begin{align*}
%     x_{t+1} = x_t - \eta \frac{g_t}{\Vert g_t \Vert}.
% \end{align*}
% Then it holds that
% \begin{align*}
%     h(x_{t+1}) \le h(x_t) - \frac{\eta}{3}  \Vert \nabla h(x_t) \Vert + \frac{8 \eta}{3} \Vert g_t - \nabla h(x_t) \Vert^2 + \frac{L \eta^2}{2}.
% \end{align*}
% \end{lem}

\section{Proofs of Finding Second-Order Stationarity} \label{apx:2nd}

% In this section, we focus on how to escape saddle points in bilevel optimization. 
The following theorem generalizes the result of Perturbed Gradient Descent~\citep{jin2017escape}
to Inexact Perturbed Gradient Descent by
allowing an inexact gradient $\hat \nabla h(x)$ that is close to the exact gradient $\nabla h(x)$.

\begin{algorithm*}[htbp]  
\caption{Inexact Perturbed Gradient Descent} \label{alg:iPGD}
\begin{algorithmic}[1] 
\STATE \textbf{for} $t = 0,1,\cdots,T-1$ \\
\STATE \quad \textbf{if} $ \Vert \hat \nabla h(x_t) \Vert \le \frac{4}{5} \epsilon$ \textbf{and} no perturbation added in the last $\fT$ steps \\[1mm]
\STATE \quad \quad $x_{t} = x_t - \eta \xi_t $, where $\xi_t \sim \sB(r)$ \\[1mm]
\STATE \quad \textbf{end if} \\[1mm]
\STATE \quad $x_{t+1} = x_t -  \eta  \hat \nabla h(x_t)$ \\[1mm]
\STATE \textbf{end for} \\[1mm]
\end{algorithmic}
\end{algorithm*}

\begin{thm}[Lemma 3.9 by \citet{huang2022efficiently}]
    \label{thm:iPGD}
Suppose $h(z): \BR^d \rightarrow \BR$ is
$L$-gradient Lipschitz and $\rho$-Hessian Lipschitz. 
Set the parameters in Algorithm \ref{alg:iPGD} as 
\begin{align} \label{cond:iota}
    \eta = \frac{1}{L}, \quad r = \frac{\epsilon}{400 \iota^3}, \quad \fT = \frac{L}{\sqrt{\rho \epsilon}} \cdot \iota ~~{\rm with }~~ \iota \ge 1 ~~ {\rm and} ~~ \delta \ge  \frac{L \sqrt{d}}{\sqrt{\rho \epsilon}} \iota^2 2^{8-\iota}.
\end{align}
Once the following inequality holds in each
iteration:
\begin{align} \label{cond:PF2BA}
    \Vert \hat \nabla h(x_t) - \nabla h(x_t) \Vert \le \underbrace{\min \left\{ 
    \frac{1}{20 \iota^2},  \frac{1}{16 \iota^2 2^{\iota}} \right\} \epsilon}_{:=\zeta},
\end{align}
then we can find a point $x_t$ satisfying
\begin{align*}
    \Vert \nabla h(x_t) \Vert \le \epsilon, \quad  \nabla^2 h(x_t)  \succeq - \sqrt{\rho \epsilon}~ I_{d}
\end{align*}
within $T= \fO \left( {\iota^4  \Delta L}{\epsilon^{-2}}   \right)$ iterations with probability $1-\delta$, where
$\Delta = h(x_0) - \inf_{x \in \BR^{d}} h(x)$.    
\end{thm}

Then we can show the convergence of perturbed F${}^2$BA.
\thmPFBA*
\begin{proof}
% The result follows from
% Lemma \ref{lem:1st}, Lemma \ref{lem:2rd}, Proposition \ref{prop:iPGD} and Proposition \ref{prop:error}.
Let $L$ be the gradient Lipschitz coefficient and $\rho$ be the Hessian Lipschitz coefficient of $\fL_{\lambda}^*(x)$, respectively.
According to Lemma \ref{lem:1st}, Lemma \ref{lem:2rd}, Lemma \ref{lem:yy-lambda}  and the setting of $\lambda$, we have
\begin{enumerate}[label=\alph*.]
    \item $\sup_{x \in \BR^{d_x}} \Vert \nabla \fL_{\lambda}^*(x) - \nabla \varphi(x) \Vert = \fO(\epsilon)$. 
    \item $\fL_{\lambda}^*(x_0) - \inf_{x \in \BR^{d_x}}  \fL_{\lambda}^*(x) = \fO(\Delta)$.
    \item $\Vert y_0 - y_{\lambda}^*(x_0) \Vert^2 + \Vert y_0 - y^*(x_0) \Vert^2 =\fO(R)$.
    \item $ L:=\sup_{x \in \BR^{d_x}}\Vert \nabla^2 \fL_{\lambda}^*(x) \Vert = \fO(\ell \kappa^3)$.
    \item $\rho: = \sup_{x \in \BR^{d_x}} \Vert \nabla^3 \fL_{\lambda}^*(x) \Vert = \fO(\ell \kappa^5) $.
\end{enumerate}
Then it suffices to show that the algorithm can find an $\epsilon$-second-order stationary point of $\fL_{\lambda}^*(x)$ under our choice of parameters.
We apply Theorem \ref{thm:iPGD} to prove this result. 
Recall $\zeta$ defined in \Eqref{cond:PF2BA}.
It remains to show that 
\begin{align} \label{eq:cond}
    \Vert \nabla \fL_{\lambda}^*(x_t) - \nabla \fL_{\lambda}^*(x_t) \Vert \le \zeta
\end{align}
holds for all $t$. Recall that
\begin{align*}
    \Vert \hat \nabla \fL_{\lambda}^*(x_t) - \nabla \fL_{\lambda}^*(x_t) \Vert \le 2 \lambda L_g \Vert y_{t}^K- y_{\lambda}^*(x_t) \Vert + \lambda L_g \Vert z_{t}^K - y^*(x_t) \Vert.
\end{align*} 
 If we let
\begin{align} \label{eq:Kt}
    K_t = \frac{4 L_g}{\mu} \log \left( 
    \frac{4 \lambda L_g ( \Vert y_t^0 - y_{\lambda}^*(x_t) \Vert^2 + \Vert z_t^0 - y^*(x_t) \Vert^2 )}{\zeta} 
    \right) 
\end{align}
Then by Theorem \ref{thm:GD-SC}, we have
\begin{align} \label{eq:induction}
     \max \left\{ \Vert y_t^K - y_{\lambda}^*(x_t) \Vert , ~\Vert z_t^K-  y^*(x_t) \Vert   \right\} \le \frac{\zeta}{4 \lambda L_g},
\end{align}
which implies \Eqref{eq:cond}. The only remaining issue is that the algorithm does not know the value of $\delta_t:= \Vert y_t^0 - y_{\lambda}^*(x_t) \Vert^2 + \Vert z_t^0 - y^*(x_t) \Vert^2  $ in \Eqref{eq:Kt}. We address this issue by showing a upper bound of it.
Suppose that we have already upper bounded $\delta_t$, we can then use it to upper bound $\delta_{t+1}$.
By \Eqref{eq:delta_t}, if we let $K_t \ge 8L_g/\mu$ for all $t$, then we can show that
\begin{align*} 
     \delta_{t+1}
    &\le \frac{1}{2} \delta_t +\frac{34 L_g^2}{\mu^2} \Vert x_{t+1} - x_t \Vert^2.
\end{align*}
Note that the value of $\Vert x_{t+1} - x_t \Vert^2$ is known in advance. This recursion gives a valid upper bound of $\delta_t$ that can be used to specify $K_t$ in the algorithm. 
Algorithm using this choice of $K_t$ can provably find an $\epsilon$-second-order stationary point of $\fL_{\lambda}^*(x)$ within 
\begin{align} \label{plug:2}
\sum_{t=0}^{T-1} K_t \le \frac{4L_g}{\mu} \sum_{t=0}^{T-1} \log \left( 
\frac{4 \lambda L_g \delta_t}{\zeta}
\right)   \le \frac{4L_g T}{\mu} \log \left( \frac{4 \lambda L_g \sum_{t=0}^{T-1}\delta_t}{\zeta T} \right)
\end{align}
 first-order oracle calls.
Telescoping over $t$, we obtain
\begin{align} \label{plug:1}
    \sum_{t=0}^{T-1} \delta_{t} \le 2 \delta_0 + \frac{68L_g^2}{\mu^2} \sum_{t=0}^{T-1} \Vert x_{t+1} - x_t \Vert^2.
\end{align}
By \Eqref{eq:final}, if we let $K_t \ge \Omega( \kappa \log(\lambda \kappa))$ for all $t$, we can obtain
\begin{align*}
\frac{1}{8 \eta} \sum_{t=0}^{T-1} \Vert x_{t+1} - x_t \Vert^2 \le \fL_{\lambda}^*(x_0) - \inf_{x \in \BR^{d_x}} \fL_{\lambda}^*(x) + \Vert y_0 - y_{\lambda}^*(x_0) \Vert^2 + \Vert y_0 - y^*(x_0) \Vert^2 .
\end{align*}
Plugging into \Eqref{plug:1} and then \Eqref{plug:2} yields the upper complexity bound of $\tilde \fO(\ell \kappa^4 \epsilon^{-2})$ as claimed. 
\end{proof}

% \section{Proofs of Distributed Algorithms}

% First of all, we recall
% the convergence results of gradient tracking methods ~\cite{qu2017harnessing} for strongly convex objectives.

% \begin{prop}
% Suppose that Assumption \ref{asm:W} holds.
% Consider optimizing
% \begin{align*}
%     \min_{x \in \BR^d} \varphi(x):= \frac{1}{m} \sum_{i=1}^m f_i(x),
% \end{align*}
% where each $f_i(x)$ is  $L$-smooth and $\mu$-strongly convex.
% The update 
% \begin{align*} 
% X_{t+1} &= W^S \cdot (X_t - \eta G_t) \\
% G_{t+1} &= W^S \cdot (G_t + \nabla f(X_{t+1}) - \nabla f(X_t))
% \end{align*}
% with $\eta = 1/L, ~ S = \lceil \chi \log (10\kappa)\rceil$
% can find
% \begin{align*}
%     \sqrt{m }\Vert \bar x_{t} - x^* \Vert + \Vert X_{t} - \vone \bar x_{t} \Vert_F + \eta \Vert G_{t} - \vone \bar g_{t} \Vert_F \le \zeta
% \end{align*}
% within an iteration number bounded by
% \begin{align*}
%     \fO \left( 
%     \kappa \log \left( \dfrac{\sqrt{m}\Vert \bar x_0- x^* \Vert +\Vert X_0 - \vone \bar x_{0} \Vert_F +\eta \Vert G_{0} - \vone \bar g_{0} \Vert_F }{\zeta} \right)
%     \right).
% \end{align*}
% \end{prop}

% \begin{proof}
% We begin from
% \begin{align*}
% &\quad \Vert \bar x_{t+1} - x^* \Vert \\
% &=  \Vert \bar x_{t+1} - \eta \bar g_t - x^* \Vert \\
% &\le \Vert \bar x_{t+1} - \eta  \nabla f( \bar x_t) - x^* \Vert + \eta \Vert \nabla f(\bar x_t) -  \bar g_t\Vert\\
% &\le (1 - \mu/L)\Vert \bar x_t - x^* \Vert + \frac{\eta L}{\sqrt{m}} \Vert X_t - \vone \bar x_t \Vert_F, 
% \end{align*}
% where in the last step we use Lemma 10 and Lemma 8(c) in \cite{qu2017harnessing}.

% Also, it holds that
% \begin{align} \label{eq:cons}
% \Vert X_{t+1} - \vone \bar x_{t+1} \Vert \le \theta^S \Vert X_t - \vone \bar x_t \Vert + \theta^S \eta \Vert G_t - \vone \bar g_t \Vert. 
% \end{align}
% Furthermore,
% \begin{align*}
% &~\quad \Vert G_{t+1} - \vone \bar g_{t+1} \Vert_F \\
% &~\le \theta^S \Vert G_t - \vone \bar g_t \Vert_F + \theta^S \Vert \nabla f(X_{t+1}) - \nabla f(X_t) \Vert_F \\
% &~\le \theta^S \Vert G_t - \vone \bar g_t \Vert_F + \rho L \Vert X_{t+1} - X_t \Vert_F \\
% &~\le \theta^S \Vert G_t -  \vone \bar g_t \Vert_F + 3 \theta^S L \Vert X_{t} - \vone \bar x_t \Vert_F \\
% &~\quad + 3 \theta^S L \Vert \vone \bar x_{t+1} - \vone \bar x_t \Vert_F + 3\theta^S L \Vert X_{t+1} - \vone \bar x_{t+1} \Vert_F \\
% &\overset{(\ref{eq:cons})}{\le} \rho(1+ 3 \theta^S L \eta) \Vert G_t - \vone \bar g_t \Vert_F + 4 \theta^S L \Vert X_t - \vone \bar  x_t \Vert_F + 3 \theta^S L  \sqrt{m} \Vert
% \bar x_{t+1} - \bar x_t \Vert  \\
% &~\le~ \rho(1+ 3 \theta^S L \eta) \Vert G_t -  \vone \bar g_t \Vert_F + 4 \theta^S L \Vert X_t - \vone \bar x_t \Vert_F \\
% &~\quad +3 \theta^S L \eta \sqrt{m} \Vert  \nabla f(\bar x_t ) - \bar g_t \Vert + 3 \theta^S L \eta \sqrt{m} \Vert \nabla f(\bar x_t) \Vert \\
% &~\le  ~\theta^S(1+ 3 \theta^S L \eta) \Vert G_t -  \vone \bar g_t \Vert_F + \theta^S (4 L + 3  L^2 \eta) \Vert X_t - \vone \bar x_t \Vert_F   + 3 \theta^S L^2 \eta \sqrt{m} \Vert \bar x_t - x^* \Vert,
% \end{align*}
% where in the last step we use Lemma 8(c) in \cite{qu2017harnessing}. Combing the two inequalities above, we obtain
% \begin{align*}
% &\quad \Vert X_{t+1} - \vone \bar x_{t+1} \Vert_F + \eta \Vert G_{t+1} - \vone \bar g_{t+1} \Vert_F \\
% &\le \theta^S (1 + 4 L \eta + 3 L^2 \eta^2) \Vert X_t - \vone \bar x_t \Vert_F + \theta^S  (1 + 3 \rho L \eta) \eta \Vert G_t -  \vone \bar g_t \Vert_F + 3 \theta^S L^2\eta^2 \sqrt{m} \Vert \bar x_t - x^* \Vert. 
% \end{align*}
% Combing the above inequalities, we obtain
% \begin{align*}
% &\quad \sqrt{m} \Vert \bar x_{t+1} - x^* \Vert + \Vert X_{t+1} - \vone \bar x_{t+1} \Vert_F + \eta \Vert G_{t+1} - \vone \bar g_{t+1} \Vert_F \\
% &\le (1 - \mu/ L +  3 \theta^S L^2 \eta^2   ) \sqrt{m}\Vert \bar x_t - x^* \Vert + \theta^S (1+ 5L \eta + 3 L^2 \eta^2) \Vert X_t - \vone X_t \Vert_F +  \theta^S  (1 + 3 \theta^S L \eta) \eta \Vert G_t -  \vone \bar g_t \Vert_F.
% \end{align*}
% It can be verified the choice of $S$ ensures
% \begin{align*}
%     \theta^S \le \min \left \{ \frac{\mu}{6L}, ~ \frac{1}{10} \right\},
% \end{align*}
% then we can prove the following linear convergence:
% \begin{align*}
% &\quad   \sqrt{m} \Vert \bar x_{t+1} - x^* \Vert + \Vert X_{t+1} - \vone \bar x_{t+1} \Vert_F + \eta \Vert G_{t+1} - \vone \bar g_{t+1} \Vert_F \\
% &\le \left(1-\theta^S \right) \left( \sqrt{m}\Vert \bar x_t - x^* \Vert +  \Vert X_t - \vone X_t \Vert_F +  \eta \Vert G_t -  \vone \bar g_t \Vert_F \right).
% \end{align*}
% We arrive at the conclusion by telescoping the above inequality.
% \end{proof}

% \propGTGD*

% \begin{proof}
% In our proof, we define $\zeta: = \Vert \hat \nabla \fL_{\lambda}^*(x) - \nabla \fL_{\lambda}^*(x) \Vert_F$. The descent lemma tells us that
% \begin{align*}
% \fL_{\lambda}^*(\bar x_{t+1}) &\le \fL_{\lambda}^*(\bar x_t) + \langle \nabla \fL_{\lambda}^*(\bar x_t), \bar x_{t+1} - \bar x_t \rangle + \frac{L}{2}
% \Vert \bar x_{t+1} - \bar x_t \Vert^2 \\
% &= \fL_{\lambda}^*(\bar x_t) - \frac{\eta}{2} \Vert \nabla \fL_{\lambda}^*(\bar x_t) \Vert^2 -
% \left( \frac{\eta}{2} - \frac{L \eta^2}{2} \right) \Vert \bar g_t \Vert^2 + \frac{\eta}{2} \Vert \nabla \fL_{\lambda}^*(\bar x_t) - \bar g_t \Vert^2 \\
% &\le  \fL_{\lambda}^*(\bar x_t) - \frac{\eta}{2} \Vert \nabla \fL_{\lambda}^*(\bar x_t) \Vert^2 -
% \left( \frac{\eta}{2} - \frac{L \eta^2}{2} \right) \Vert \bar g_t \Vert^2 \\
% &\quad+ \eta \left\Vert \nabla \fL_{\lambda}^*(\bar x_t) - \frac{1}{m} \mathbf{1}^\top \nabla \fL_{\lambda}^*(X_t) \right\Vert^2
% + \frac{\eta}{m} \left \Vert  \hat \nabla \fL_{\lambda}^*(X_t) - \nabla \fL_{\lambda}^*(X_t)  \right \Vert_F^2. 
% \end{align*}
% Additionally,
% \begin{align*}
% &\quad \Vert G_{t+1} - \vone \bar g_{t+1} \Vert \\
% &\le \theta^{S_{\rm out}} \Vert G_t - \vone \bar g_t \Vert + \theta^{S_{\rm out}} \Vert \hat \nabla \fL_{\lambda}^*(X_{t+1}) - \hat \nabla \fL_{\lambda}^*(X_t) \Vert_F \\
% &\le \theta^{S_{\rm out}} \Vert G_t - \vone \bar g_t \Vert_F + \theta^{S_{\rm out}} L \Vert X_{t+1} - X_t \Vert_F + 2 \theta^{S_{\rm out}}  \zeta \\
% &\le \theta^{S_{\rm out}} \Vert G_t -  \vone \bar g_t \Vert_F + 3 \theta^{S_{\rm out}} L \Vert X_{t} - \vone \bar x_t \Vert_F  \\
% &\quad + 3 \theta^{S_{\rm out}} L \sqrt{m}\Vert  \bar x_{t+1} -  \bar x_t \Vert + 3\theta^{S_{\rm out}} L \Vert X_{t+1} - \vone \bar x_{t+1} \Vert_F + 2 \theta^{S_{\rm out}}  \zeta \\
% &\le \theta^{S_{\rm out}}(1+ 3 \theta^{S_{\rm out}} L \eta) \Vert G_t -  \vone \bar g_t \Vert_F + 4 \theta^{S_{\rm out}} L \Vert X_t - \vone \bar x_t \Vert_F + 3 \theta^{S_{\rm out}}   \eta \sqrt{m} \Vert \bar  g_{t} \Vert +2 \theta^{S_{\rm out}} \zeta.
% \end{align*}
% Combing this inequality with (\ref{eq:cons}) yields
% \begin{align*}
% &\quad \Vert X_{t+1} - \vone \bar x_{t+1} \Vert_F + \eta \Vert G_{t+1} - \vone \bar g_{t+1} \Vert_F \\
% &\le \rho (1 + 4 L \eta ) \Vert X_t - \vone \bar x_t \Vert_F + \theta^{S_{\rm out}}  (1 + 3 \theta^{S_{\rm out}} L \eta) \eta \Vert G_t -  \vone \bar g_t \Vert_F + 3 \theta^{S_{\rm out}} L \eta^2 \sqrt{m} \Vert \bar g_t \Vert + 2 \theta^{S_{\rm out}} \eta \zeta. 
% \end{align*}
% Taking square,
% \begin{align*}
%     &\quad \Vert X_{t+1} - \vone \bar x_{t+1} \Vert_F^2 + \eta^2 \Vert G_{t+1} - \vone \bar g_{t+1} \Vert_F^2 \\
% &\le 4 \theta^{2S_{\rm out}} (1 + 4 L \eta)^2 \Vert X_t - \vone \bar x_t \Vert_F^2 + 4 \theta^{2S_{\rm out}} (1+ 3 \theta^{S_{\rm out}} L \eta)^2 \eta^2 \Vert G_t - \vone  \bar g_t \Vert_F^2 \\
% &\quad + 36 \theta^{2S_{\rm out}} L^2 \eta^4 m \Vert \bar g_t \Vert^2 + 16 \theta^{2S_{\rm out}} \eta^2  \zeta^2.
% \end{align*}
% Therefore,
% \begin{align*}
%     &\quad \fL_{\lambda}^*(\bar x_{t+1})  + \frac{1}{m \eta } \Vert X_{t+1} - \vone \bar x_{t+1} \Vert_F^2 + \frac{\eta}{m } \Vert G_{t+1} - \vone \bar g_{t+1} \Vert_F^2 \\
%     &\le \fL_{\lambda}^*(\bar x_t) - \frac{\eta}{2} \Vert \nabla \fL_{\lambda}^*(\bar x_t) \Vert^2 - \left( \frac{\eta}{2} - \frac{L \eta^2}{2} - 36 \theta^{2S_{\rm out}} L^2 \eta^3  \right) \Vert \bar g_t \Vert^2 \\
%     &\quad +
% \end{align*}
% \end{proof}


\end{document}